\documentclass[11pt]{amsart}
\usepackage[margin=1.05in]{geometry}
\usepackage{amsmath,amssymb,amsthm,mathtools}
\usepackage{booktabs,longtable,array}
\usepackage[hidelinks]{hyperref}
\hypersetup{pdftitle={Open problems settled by transferred methods},pdfauthor={Keio}}

\newtheorem{theorem}{Theorem}[section]
\newtheorem{lemma}[theorem]{Lemma}
\newtheorem{proposition}[theorem]{Proposition}
\newtheorem{corollary}[theorem]{Corollary}
\newtheorem{conjecture}[theorem]{Conjecture}
\theoremstyle{definition}
\newtheorem{definition}[theorem]{Definition}
\newtheorem{problem}[theorem]{Problem}
\theoremstyle{remark}
\newtheorem{remark}[theorem]{Remark}

\providecommand{\C}{\mathbb C}
\providecommand{\R}{\mathbb R}
\providecommand{\Q}{\mathbb Q}
\providecommand{\Z}{\mathbb Z}
\providecommand{\N}{\mathbb N}
\providecommand{\Qb}{\overline{\mathbb Q}}
\usepackage{url}
\setlength{\emergencystretch}{2.5em}
\begin{document}
\title[Problems from the literature]{Problems from the literature settled by transferred methods:\\
rook walks, HCM laws, hypocycloid division points,\\
mixture likelihoods, polygamma Wronskians and composite sequences}
\author{Keio}
\date{23 October 2026}

\begin{abstract}
We take up twenty-five conjectures, questions and guesses stated in the
literature by other authors and settle sixteen of them as posed (one
conditionally on an earlier, not yet refereed manuscript of this series,
and one by combining known results); for the remaining nine we prove
partial results whose scope is stated precisely.
The proofs transfer a small set of methods developed in earlier
manuscripts of this series, sometimes in combination. Highlights:
(i) for every $d$ the numbers of diagonal rook paths in $d$ dimensions
satisfy a linear recurrence of order at most $d$ with polynomial
coefficients, produced by an explicit algorithm, which answers a challenge
of Kauers and Zeilberger; we also prove their conjectured large-$d$
asymptotics and the guessed recurrences for $d\le7$ (Problem~17 of Chen and
Kauers for $4\le d\le7$);
(ii) the Gompertz critical-point conjecture of Polyanskiy and Wu is false,
and mixtures of two analytic strongly log-concave densities can have
arbitrarily many modes;
(iii) Bosch's conjecture that $\sqrt{\gamma_t\gamma_s}$ is hyperbolically
completely monotone fails exactly when $|t-s|>1/2$, the condition
$c\le a+b$ is necessary for $\Gamma_c\times\mathbf B_{a,b}$ to be HCM, as
Bosch and Simon believed, and $\mathrm{HM}\cap\mathrm{GGC}\ne\mathrm{HCM}$;
(iv) contrary to a speculation of Mani and Rubinstein-Salzedo, individual
division points of the tricuspoid can be constructible when $n\nmid6$ (the
first case is $n=32$), and we give the complete criterion;
(v) several guesses and degree questions of Qi and coauthors on
Wronskians of polygamma products are settled;
(vi) assuming estimates from an earlier manuscript of this series, the
binomial nonparametric MLE attains the minimax Wasserstein-$p$ rate for
every $p\ge1$;
(vii) every real biquadratic field contains infinitely many independent
Pisot numbers all of whose integer parts of powers are composite and
$\equiv3,4\pmod 6$, and Šiurys's conjecture on composite $k$-step
Fibonacci-like sequences holds for all $3\le k\le130$ and for a set of
$k$ of lower density at least $3/4$.
The partial results also cover a conjecture and a problem of Faust et al.\
on Floquet isospectrality (the conjecture is proved for every period with a
divisor in $[5,69]$, and a structure theorem holds for all period lattices)
and two questions of Erchenko, Jakobson and Tsypin on spectral volumes of
weighted Laplacians. Each problem is quoted from
its source and each outcome is labelled.
\end{abstract}

\maketitle
\tableofcontents

\section{Introduction}\label{sec:intro}

\subsection{What this paper does}
An earlier review of the manuscripts of this series singled out ten
results, each of which answered a problem posed by other authors. Here we
ask how far the \emph{methods} behind those results reach. For each
method we searched the literature for explicitly stated conjectures,
questions and guesses, checked that they were still open as of October
2026, discarded those already treated in the series, and attacked the
rest. This paper records what came out: complete solutions where we found
them, counterexamples where the expectation was wrong, and partial results
where neither was within reach.

The sections are independent. Each begins with the problem quoted from its
source, followed by the main theorem and its proof. Labels and notation are
local to each section.

\subsection{The transferred methods}
We refer to the ten methods by number; the earlier manuscripts contain the
original applications.
\begin{enumerate}
\item[M1.] \emph{Singularity budgets for recurrences}: a recurrence of
order $r$ forces an ODE with at most $r$ nonzero singular points, and
critical values of a real Laurent polynomial on a torus produce genuine
singularities (Franel's conjecture \cite{K-006,K-034}).
\item[M2.] \emph{Likelihood transport for mixtures}: optimality of the
NPMLE, transport--entropy inequalities and Chebyshev--Bernstein bounds
(binomial mixtures \cite{K-029}).
\item[M3.] \emph{De Rham and $\mathcal D$-module counts} for minimal
orders (banana maximal cuts \cite{K-032}).
\item[M4.] \emph{Perturbative Floquet rigidity}: coordinate averaging,
induction on dimension and eigenvalue perturbation \cite{K-034}.
\item[M5.] \emph{Spectral cones of weighted Laplacians}: norm
obstructions, first-contact and double-contact arguments \cite{K-015}.
\item[M6.] \emph{Tower composition} for Ewald points \cite{K-041}.
\item[M7.] \emph{Local monodromy and irreducibility} (isotropic Fuchs
squeeze) \cite{K-037}.
\item[M8.] \emph{Galois and transcendence classifications of division
points and quantiles} \cite{K-024}.
\item[M9.] \emph{Laplace densities and complete monotonicity}: Bernstein's
theorem turns complete monotonicity into the sign of a density, and a
vanishing integral forces a sign change \cite[Part~VIII]{K-033}.
\item[M10.] \emph{Finite-field return and trace congruences} for composite
values of recurrences \cite{K-007}.
\end{enumerate}
Section~\ref{R1:sec} combines M1, M3 and M7; Section~\ref{R7-sec} combines
M8 and M9; Section~\ref{R5:sec} combines M2 with Fourier windows and
degenerate critical points; Section~\ref{R2:sec} combines M9 with a
likelihood-ratio ordering; Section~\ref{R6:sec} uses M10 with
Bang--Zsigmondy primitive divisors; Section~\ref{R3:sec} uses M4 with
interval certification; Section~\ref{R4:sec} uses M5. Method M6 did not
lead to any problem that we could treat.

\subsection{Outcomes}
Table~\ref{tab:outcomes} lists every problem we treat. \emph{Closed} means
that the statement asked in the source is proved or disproved as posed;
\emph{partial} means that only the stated part is proved. Some closures are
easy once the problem is seen from the right angle, and we say so in the
table and in the text.

{\small
\begin{longtable}{@{}p{0.03\textwidth}>{\raggedright\arraybackslash}p{0.37\textwidth}>{\raggedright\arraybackslash}p{0.2\textwidth}>{\raggedright\arraybackslash}p{0.33\textwidth}@{}}
\caption{Problems treated and outcomes.}\label{tab:outcomes}\\
\toprule
 & Problem & Source & Outcome\\
\midrule\endfirsthead
\toprule
 & Problem & Source & Outcome\\
\midrule\endhead
\bottomrule\endfoot
1 & Order-$d$ recurrence for $d$-dimensional diagonal rook paths, every $d$ & Kauers--Zeilberger \cite[\S2]{R1-KZ} & \textbf{Closed}: order $\le d$ for all $d$, explicit algorithm (Thm.~\ref{R1:thm:upper})\\
2 & Large-$d$ asymptotics of rook walks & Kauers--Zeilberger \cite[\S8]{R1-KZ} & \textbf{Closed} (Thm.~\ref{R1:thm:largeD})\\
3 & Prove the guessed recurrences for $d=4,\dots,12$ & Chen--Kauers \cite[Problem~17]{R1-CK} & \emph{Partial}: proved for $d\le7$, computer-assisted exact arithmetic (Prop.~\ref{R1:prop:soft})\\
4 & Minimal order is exactly $d$ & Kauers--Zeilberger (empirical) & \emph{Partial}: $d\le3$ (Prop.~\ref{R1:prop:small})\\
5 & Constructible individual division points of the tricuspoid & Mani--Rubinstein-Salzedo \cite[Rem.~5.7]{R7-MRS16} & \textbf{Closed}: all three speculations fail; complete criterion (Thm.~\ref{R7-thm-tricuspoid})\\
6 & Division of the drawn lemniscate & \cite[Rem.~1.4]{R7-MRS16} & \textbf{Closed} by combining known results (Videla; Baek et al.; Cox--Shurman) (Thm.~\ref{R7-thm-drawn})\\
7 & Is $\mathrm{HM}\cap\mathrm{GGC}=\mathrm{HCM}$? & Jedidi--Simon \cite[Rem.~4.1(b)]{R7-JS13} & \textbf{Closed}: no; an elementary example (Thm.~\ref{R7-thm-JS})\\
8 & $\sqrt{\gamma_t\gamma_s}$ is HCM for all $s,t$; $n$-fold version & Bosch \cite[Rem.~1]{R7-Bosch15} & \textbf{Closed}: false; HCM iff $|t-s|\le1/2$; $n$-fold version fails for every $n\ge2$ (Thm.~\ref{R7-thm-bosch}, Prop.~\ref{R7-prop-nfold})\\
9 & $c\le a+b$ is necessary for $\Gamma_c\times\mathbf B_{a,b}$ to be HCM & Bosch--Simon \cite{R7-BoschSimon16} & \textbf{Closed}: true (Thm.~\ref{R7-thm-gammabeta})\\
10 & At most $O(a)$ critical points of Gompertz convolutions & Polyanskiy--Wu \cite[\S5.2]{R5-PW20} & \textbf{Closed}: false, $\Omega(a^2)$ (Thm.~\ref{R5:thm:gompertz})\\
11 & Is the $\Omega(k^2)$ mode count tight? & Polyanskiy--Wu \cite[\S5.5]{R5-PW20} & \textbf{Closed} as posed: unbounded already for $k=2$ (Thm.~\ref{R5:thm:modes})\\
12 & Wasserstein-$p$ guarantees for the binomial NPMLE & Vinayak--Kong--Valiant--Kakade \cite[\S5]{R5-VKVK19} & \textbf{Closed} conditionally on four estimates of \cite{K-029} (Thm.~\ref{R5:thm:wp})\\
13 & Remark~17 guess and $\Omega$-problem & Qi \cite{R2-QiCR2022} & \textbf{Closed}; easy via a splitting identity (Thm.~\ref{R2:thm:A})\\
14 & Remark~18 guess and $\Omega$-problem & Qi \cite{R2-QiCR2022} & \textbf{Closed} (guess); threshold: one half closed, the other characterized (Thm.~\ref{R2:thm:B})\\
15 & Remark~4 guess and $\Lambda$-problem & Yin--Han--Qi \cite{R2-YHQ2024} & \textbf{Closed} (guess); threshold: one half closed, the other characterized (Thm.~\ref{R2:thm:C})\\
16 & Remark~2 sign guesses & Yin--Han--Qi \cite{R2-YHQ2024} & \textbf{Closed}, elementary (Prop.~\ref{R2:prop:D})\\
17 & CM degrees of $H_2,H_3$ & Qi--Agarwal \cite[\S5.4]{R2-QiAgarwal2019} & \textbf{Closed} (Thm.~\ref{R2:thm:E})\\
18 & CM degrees of the function (5.2) & Qi--Agarwal \cite[\S5.5]{R2-QiAgarwal2019} & \emph{Partial}: closed for $\mu\ge4$ and $\lambda=0$; for $\lambda<0$ non-integer in $(0,1)$, values not in closed form (Thm.~\ref{R2:thm:F})\\
19 & Sharp threshold and lower bound in (3.1) & Qi \cite[Rem.~4.1]{R2-QiMIA2021} & \emph{Partial} (Prop.~\ref{R2:prop:G})\\
20 & Infinitely many biquadratic Pisot $\alpha$ with $\lfloor\alpha^n\rfloor\equiv3,4\pmod6$ & Hahn--Ismailescu--Kim--Kim \cite[\S5]{R6-HIKK} & \textbf{Closed}, in a stronger field-by-field form; the literal statement is easy (Thm.~\ref{R6:thm:biquad})\\
21 & Composite $k$-step Fibonacci-like sequences for every $k>2$ & Šiurys \cite{R6-Siurys} & \emph{Partial}: all $3\le k\le130$, explicit residue classes, lower density $\ge3/4$ (Thm.~\ref{R6:thm:siurys})\\
22 & Nonzero conjugate-reflection-symmetric potentials Floquet isospectral to $0$ for every period $q>3$ & Faust et al.\ \cite[Conj.~2]{R3-FFORS} & \emph{Partial}: every $q$ that is even (known) or has a divisor in $[5,69]$; exact for $q=5,9$, interval certificates for odd $q\le69$; open e.g.\ for primes $q>69$ (Cor.~\ref{R3:cor:conj2})\\
23 & Classify the period lattices $\Gamma\subset\Z^d$ admitting nonzero potentials Floquet isospectral to $0$ & Faust et al.\ \cite[Problem~1]{R3-FFORS} & \emph{Partial}: structure theorem for all lattices; rigidity for all $\Gamma\supseteq2\Z^d$ or $3\Z^d$ and an infinite family of non-rectangular lattices; reduction to a core statement; planar lattices of index $\le12$ (Thm.~\ref{R3:thm:structure}, Cor.~\ref{R3:cor:criterion}, Prop.~\ref{R3:prop:computations})\\
24 & Does adding an edge increase the spectral volume $P(G)$? & Erchenko--Jakobson--Tsypin \cite{R4-EJT} & \emph{Partial}: connected graphs on at most $6$ vertices; two infinite families (Thms.~\ref{R4:thm:small}, \ref{R4:thm:matching}, \ref{R4:thm:hadwiger})\\
25 & Does suspension never decrease $P(G)$? & Erchenko--Jakobson--Tsypin \cite{R4-EJT} & \emph{Partial}: graphs on at most $4$ vertices; $P(G\vee K_1)\ge\gamma_nP(G)$ in general (Thms.~\ref{R4:thm:Q4small}, \ref{R4:thm:Q4partial})\\
\end{longtable}
}

\subsection{Verification}
Every computation used in a proof is exact (rational or algebraic
arithmetic, certified interval arithmetic, or finite certificates checked
by an independent verifier); numerical experiments that only illustrate a
statement are labelled as such. The scripts named in the text are collected
in the supplementary archive that accompanies this paper; each prints the
checks it performs and was run under \texttt{python3 -I}.

\subsection{Scope}
Several closures are modest, and we flag them: rows 13, 16 and 20 are easy
once the right identity or example is found, row~6 combines known results,
and row~12 depends on the binomial-mixture theorem of \cite{K-029}, which
has not been refereed. The most substantial results are those in rows 1, 2,
5 and 8--12 of Table~\ref{tab:outcomes} and, among the partial results, rows
3 and 21--23. Statements about the literature are as of October 2026, and a
missed recent preprint cannot be ruled out.

%% ==================================================================
%% Section: Rook walks
%% ==================================================================
% R1 section: minimal recurrence orders (area R1).  LaTeX fragment, no preamble.
% Requires amsmath, amssymb, amsthm and the environments theorem, lemma, proposition,
% corollary, definition, remark, conjecture, problem.
\providecommand{\Rone}{\mathrm{R1}}
\providecommand{\RoneJ}{\mathcal{J}}
\providecommand{\RoneN}{\mathbb{N}}
\providecommand{\RoneQ}{\mathbb{Q}}
\providecommand{\RoneC}{\mathbb{C}}
\providecommand{\texorpdfstring}[2]{#1}

\section{Recurrences of order at most \texorpdfstring{$d$}{d} for diagonal rook paths in every dimension}\label{R1:sec}

\subsection{The problems and the status of this section}\label{R1:sec:problems}

For $d\ge1$ let
\[
R_d(x_1,\dots,x_d)=\frac{1}{1-\sum_{i=1}^{d}\frac{x_i}{1-x_i}}
=\sum_{m\in\RoneN^d}\rho_d(m)\,x^m ,
\qquad a^{(d)}_n=\rho_d(n,n,\dots,n).
\]
Here $\rho_d(m)$ is the number of ways a rook can move from the origin to $m$ in
$\RoneN^d$ when every move increases exactly one coordinate by a positive integer; thus
$a^{(2)}_n$ is OEIS A051708, $a^{(3)}_n$ is A144045, and the table $a^{(d)}_n$ is A181731
\cite{R1-OEIS}.
The sequence $(a^{(d)}_n)_n$ is P-recursive by Lipshitz's theorem on diagonals
\cite{R1-Lipshitz}. Kauers and Zeilberger computed (by guessing) recurrences for
$2\le d\le 12$ and formulated the following challenges.

\begin{problem}[{Kauers--Zeilberger \cite[\S2]{R1-KZ}}]\label{R1:prob:KZ}
``The recurrences for $d$-dimensional rook walks turn out (empirically) to have order $d$
(for $2\le d\le 12$). It is very possible that it won't be too hard to prove this sharp
upper bound in general, or even a weaker---but still realistic---one. This may enable one to
give a `soft' proof that the empirically `guessed' recurrences are indeed rigorously proved.''
\end{problem}

\begin{problem}[{Chen--Kauers \cite[Problem 17]{R1-CK}}]\label{R1:prob:CK}
``For $d=4,\dots,12$, prove recurrence equations for the diagonals of the rational series
$1/(1-\sum_{i=1}^{d}\frac{x_i}{1-x_i})$ conjectured in [Kauers--Zeilberger].''
(The source adds: ``For $d=1,2$, the problem is easy. For $d=3$, it was solved in''
\cite{R1-BCvHP}.)
\end{problem}

\begin{problem}[{Kauers--Zeilberger \cite[\S8]{R1-KZ}}]\label{R1:prob:KZ8}
``Let $w_n(d)$ be the number of ways a rook can positively walk from $[0^d]$ to $[n^d]$.''
``The asymptotic formulas for the individual $n$ (for $1\le n\le 4$) lead one to conjecture
that the leading asymptotics for $w_n(d)$ as $d\to\infty$ is
$e^{n-1}\frac{(nd)!}{n!^d}\bigl(1+O(\frac1d)\bigr)$.''
``We leave the rigorous proof of this as yet another challenge to the reader.''
\end{problem}

Note that $w_n(d)=a^{(d)}_n$.

\paragraph{Status.}
Theorem~\ref{R1:thm:upper} proves, for \emph{every} $d\ge1$, that $a^{(d)}_n$ satisfies a
nonzero linear recurrence of order at most $d$ with polynomial coefficients, valid for all
$n\ge1$, and gives a finite algorithm producing it. This proves the upper bound asked for in
Problem~\ref{R1:prob:KZ} and supplies rigorously proved recurrences of order at most $d$ for every $d$.
Together with an explicit degree bound (Lemma~\ref{R1:lem:degbound}) it yields exactly the
``soft proof'' anticipated in Problem~\ref{R1:prob:KZ}: Proposition~\ref{R1:prop:soft} proves,
for $2\le d\le7$, the unique recurrence of order $d$ and of the coefficient degree reported in
\cite[\S5]{R1-KZ} that fits the initial terms. This settles Problem~\ref{R1:prob:CK} for
$d=4,5,6,7$, with one caveat: we have not compared coefficient lists with the published
files \cite{R1-KZfiles}; the identification with the recurrences of \cite{R1-KZ}, which were
fitted to initial terms, rests on this uniqueness. For $8\le d\le12$ Problem~\ref{R1:prob:CK}
remains open: rigorous recurrences of order at most $d$ are proved to exist and are
computable, and for every $d\le12$ their coefficient degree, computed modulo a large prime,
equals the degree reported in \cite[\S5]{R1-KZ} (Remark~\ref{R1:rem:computations}), but the
specific recurrences of \cite{R1-KZ} are not proved here for $d\ge8$.
Theorem~\ref{R1:thm:largeD} proves the conjectured asymptotics of Problem~\ref{R1:prob:KZ8}
(an explicit second term is recorded, on numerical evidence only, in
Remark~\ref{R1:rem:secondterm}). The \emph{sharpness} of the bound (that no recurrence of
order $<d$ exists) is proved here only for $d\le3$ (Proposition~\ref{R1:prop:small}); for
$d\ge4$ it remains open. We reduce it to a statement about local monodromy
(Lemma~\ref{R1:lem:index}, Conjecture~\ref{R1:conj:sharp}), which holds for the differential
operators found by guessing for $d\le5$; this is evidence only, since these operators are not
certified to be minimal.

\subsection{A Laguerre integral for the diagonal}\label{R1:sec:laguerre}

For $n\ge1$ put
\begin{equation}\label{R1:eq:ell}
\ell_n(s)=\sum_{k=1}^{n}\binom{n-1}{k-1}\frac{s^k}{k!}=L^{(-1)}_n(-s),
\end{equation}
where $L^{(\alpha)}_n$ is the Laguerre polynomial. (Integrals of products of Laguerre polynomials against
$e^{-s}$ have a classical combinatorial meaning, cf.\ Even--Gillis \cite{R1-EG}; we only use the elementary
identities below.)

\begin{lemma}\label{R1:lem:rep}
For all $d\ge1$ and $n\ge1$,
\[
a^{(d)}_n=\sum_{k\in\{1,\dots,n\}^d}\frac{(k_1+\dots+k_d)!}{k_1!\cdots k_d!}
\prod_{i=1}^{d}\binom{n-1}{k_i-1}
=\int_0^\infty e^{-s}\,\ell_n(s)^d\,ds .
\]
\end{lemma}

\begin{proof}
Put $u_i=x_i/(1-x_i)$. For $|x_i|$ small, $R_d=\sum_{K\ge0}(u_1+\dots+u_d)^K$, and the
multinomial theorem gives $R_d=\sum_{k\in\RoneN^d}\frac{|k|!}{k_1!\cdots k_d!}\prod_iu_i^{k_i}$,
an absolutely convergent expansion near $0$. For $n\ge1$ we have $[x^n]u^k=[x^{n-k}](1-x)^{-k}
=\binom{n-1}{k-1}$ if $1\le k\le n$, and $[x^n]u^k=0$ if $k=0$ or $k>n$. Extracting the
coefficient of $x_1^n\cdots x_d^n$ gives the finite sum. Finally insert
$|k|!=\int_0^\infty e^{-s}s^{|k|}ds$ and factor the finite sum over $k$:
$\sum_k\frac{|k|!}{\prod k_i!}\prod_i\binom{n-1}{k_i-1}=\int_0^\infty e^{-s}
\prod_i\bigl(\sum_{k_i=1}^n\binom{n-1}{k_i-1}\frac{s^{k_i}}{k_i!}\bigr)ds$.
\end{proof}

\begin{lemma}\label{R1:lem:ell}
For $n\ge1$:
\begin{enumerate}
\item[(a)] $\ell_n(0)=0$, $\ell_n'(0)=1$;
\item[(b)] $s\,\ell_n''+s\,\ell_n'-n\,\ell_n=0$;
\item[(c)] $(n+1)\ell_{n+1}=(n+s)\ell_n+s\ell_n'$ and $\ell_{n+1}'=\ell_n+\ell_n'$.
\end{enumerate}
\end{lemma}

\begin{proof}
Write $\ell_n=\sum_kc_ks^k$, $c_k=\binom{n-1}{k-1}/k!$. (a) is clear. (b) The coefficient of
$s^k$ in $s\ell''+s\ell'-n\ell$ is $(k+1)kc_{k+1}+(k-n)c_k=\bigl(k\binom{n-1}{k}-(n-k)
\binom{n-1}{k-1}\bigr)/k!=0$. (c) For $1\le k\le n$ the coefficient of $s^k$ in $(n+s)\ell_n+s\ell_n'$ is
$(n+k)c_k+c_{k-1}=\binom{n-1}{k-1}\frac{(n+k)(n-k+1)+k(k-1)}{(n-k+1)k!}
=\binom{n-1}{k-1}\frac{n(n+1)}{(n-k+1)k!}=(n+1)\binom{n}{k-1}\frac1{k!}$, which is the
coefficient of $s^k$ in $(n+1)\ell_{n+1}$; for $k=n+1$ both sides have coefficient $1/n!$. The second identity is Pascal's rule for the
coefficient of $s^{k-1}$: $\binom{n}{k-1}=\binom{n-1}{k-2}+\binom{n-1}{k-1}$.
\end{proof}

\subsection{Integration by parts and a reduction to \texorpdfstring{$d$}{d} integrals}\label{R1:sec:reduction}

Fix $d\ge1$. For $n\ge1$, $0\le i\le d$ and integers $k\ge i-d$ put
\[
\RoneJ_n(i,k)=\int_0^\infty e^{-s}s^k\,\ell_n(s)^{d-i}\ell_n'(s)^i\,ds .
\]
The integral converges: by Lemma~\ref{R1:lem:ell}(a) the integrand is
$O(s^{k+d-i})$ at $0$, and $k+d-i\ge0$. Note $\RoneJ_n(0,0)=a^{(d)}_n$.

\begin{lemma}[integration by parts]\label{R1:lem:ibp}
For $n\ge1$, $0\le i\le d$ and $k\ge i-d+1$,
\begin{equation}\label{R1:eq:ibp}
(1+i)\RoneJ_n(i,k)=k\,\RoneJ_n(i,k-1)+(d-i)\,\RoneJ_n(i+1,k)+i\,n\,\RoneJ_n(i-1,k-1),
\end{equation}
where a term whose coefficient vanishes ($i=d$, resp.\ $i=0$) is omitted.
\end{lemma}

\begin{proof}
Let $m_i=\ell^{d-i}\ell'^i$ ($\ell=\ell_n$) and $g=e^{-s}s^km_i$. Since $k+d-i\ge1$,
$g(0)=0$; also $g(s)\to0$ as $s\to\infty$. Hence $\int_0^\infty g'\,ds=0$. By
Lemma~\ref{R1:lem:ell}(b), $\ell''=n\ell/s-\ell'$, so
$m_i'=(d-i)m_{i+1}+i\ell^{d-i}\ell'^{i-1}\ell''=(d-i)m_{i+1}+\frac{in}{s}m_{i-1}-im_i$, and
$g'=e^{-s}\bigl(-s^km_i+ks^{k-1}m_i+(d-i)s^km_{i+1}+ins^{k-1}m_{i-1}-is^km_i\bigr)$. Each of
the four integrals obtained by splitting converges (the exponents of $s$ at $0$ are
$k+d-i$, $k-1+d-i$, $k+d-i-1$, $k+d-i$, all $\ge0$), which gives \eqref{R1:eq:ibp}.
\end{proof}

\begin{lemma}[reduction]\label{R1:lem:reduction}
There are polynomials $\gamma^{(i,k)}_b\in\RoneQ[n]$, for $0\le i\le d$, $k\ge0$ and
$1-d\le b\le0$, depending only on $d,i,k,b$, such that for every integer $n\ge1$
\[
\RoneJ_n(i,k)=\sum_{b=1-d}^{0}\gamma^{(i,k)}_b(n)\,\RoneJ_n(0,b).
\]
\end{lemma}

\begin{proof}
All identities below hold for every integer $n\ge1$, with coefficients that are polynomials
in $n$ independent of $n$; they are obtained from \eqref{R1:eq:ibp} by linear elimination in
which we only divide by the nonzero integers $d-i+1$ and $d+1$. ``Span'' means
$\RoneQ[n]$-linear span.

\emph{Step 1 (lowering $i$).} We claim: for $0\le i\le d$ and $k\ge i-d$,
\begin{equation}\label{R1:eq:claimA}
\begin{gathered}
\RoneJ_n(i,k)=\sum_{b=k-i}^{k}\lambda^{(i,k)}_b(n)\,\RoneJ_n(0,b),\\
\lambda^{(i,k)}_k=\binom{d}{i}^{-1},\qquad
\lambda^{(i,k)}_{k-i}=(-1)^i\binom{k}{i}\binom{d}{i}^{-1}
\end{gathered}
\end{equation}
(for $i=0$ both formulas read $\lambda^{(0,k)}_k=1$). This is trivial for $i=0$. For $i\ge1$,
relation \eqref{R1:eq:ibp} with $(i-1,k)$ in place of $(i,k)$ is valid because
$k\ge i-d=(i-1)-d+1$, and reads
\[
(d-i+1)\RoneJ_n(i,k)=i\,\RoneJ_n(i-1,k)-k\,\RoneJ_n(i-1,k-1)-(i-1)n\,\RoneJ_n(i-2,k-1).
\]
By induction the three terms on the right lie in the span of $\RoneJ_n(0,b)$ for
$b\in[k-i+1,k]$, $[k-i,k-1]$ and $[k-i+1,k-1]$ respectively. The index $b=k$ occurs only in
the first term and $b=k-i$ only in the second, which gives
$\lambda^{(i,k)}_k=\frac{i}{d-i+1}\lambda^{(i-1,k)}_k$ and
$\lambda^{(i,k)}_{k-i}=\frac{-k}{d-i+1}\lambda^{(i-1,k-1)}_{k-i}$, and hence the formulas in
\eqref{R1:eq:claimA}. In particular $\lambda^{(d,k)}_k=1$.

\emph{Step 2 (a recurrence in $k$).} Let $k\ge1$. Relation \eqref{R1:eq:ibp} with $i=d$
(valid since $k\ge1$) gives $(1+d)\RoneJ_n(d,k)=k\RoneJ_n(d,k-1)+dn\RoneJ_n(d-1,k-1)$.
Insert \eqref{R1:eq:claimA}: the left side is $(1+d)\RoneJ_n(0,k)$ plus a combination of
$\RoneJ_n(0,b)$, $k-d\le b\le k-1$; on the right, $\RoneJ_n(d,k-1)$ involves $b\in[k-1-d,k-1]$,
the coefficient of $\RoneJ_n(0,k-1-d)$ being $(-1)^d\binom{k-1}{d}$, and $\RoneJ_n(d-1,k-1)$
involves $b\in[k-d,k-1]$. Therefore
\begin{equation}\label{R1:eq:krec}
\RoneJ_n(0,k)=\sum_{b=k-d-1}^{k-1}\mu^{(k)}_b(n)\RoneJ_n(0,b),\qquad
\mu^{(k)}_{k-d-1}=\frac{(-1)^dk}{d+1}\binom{k-1}{d},
\end{equation}
and $\mu^{(k)}_{k-d-1}=0$ when $1\le k\le d$.

\emph{Step 3.} By induction on $k$: $\RoneJ_n(0,k)$ lies in the span $B$ of
$\RoneJ_n(0,1-d),\dots,\RoneJ_n(0,0)$ for every $k\ge1-d$. This is trivial for $k\le0$. For
$1\le k\le d$, \eqref{R1:eq:krec} only involves indices $b$ in $[k-d,k-1]\subset[1-d,k-1]$;
for $k\ge d+1$ it involves indices $b$ in $[k-d-1,k-1]\subset[0,k-1]$.

\emph{Step 4.} Let $k\ge0$. By \eqref{R1:eq:claimA}, $\RoneJ_n(i,k)$ is a combination of
$\RoneJ_n(0,b)$ with $b\ge k-i\ge-d$, and $b=-d$ can only occur for $(i,k)=(d,0)$, where its
coefficient is $(-1)^d\binom{0}{d}=0$. Hence only $b\ge1-d$ occur, and Step 3 finishes the
proof.
\end{proof}

\subsection{The upper bound}\label{R1:sec:upper}

\begin{theorem}\label{R1:thm:upper}
For every $d\ge1$ there are polynomials $c_0,\dots,c_d\in\mathbb{Z}[n]$, not all zero, such
that
\[
\sum_{j=0}^{d}c_j(n)\,a^{(d)}_{n+j}=0\qquad\text{for all integers }n\ge1 .
\]
In particular the minimal order of a recurrence with polynomial coefficients satisfied by
$(a^{(d)}_n)_{n\ge0}$ is at most $d$. The polynomials $c_j$ are produced by the finite
procedure in the proof.
\end{theorem}

\begin{proof}
Iterating Lemma~\ref{R1:lem:ell}(c), for each $j\ge0$ there are $\alpha_j,\beta_j\in
\RoneQ(n)[s]$, of degree at most $j$ in $s$ and with denominators dividing
$(n+1)(n+2)\cdots(n+j)$, such that $\ell_{n+j}=\alpha_j\ell_n+\beta_j\ell_n'$ for all $n\ge1$
(start with $\alpha_0=1,\beta_0=0$ and use
$(\ell_{m+1},\ell_{m+1}')=(\frac{m+s}{m+1}\ell_m+\frac{s}{m+1}\ell_m',\ \ell_m+\ell_m')$ with
$m=n,\dots,n+j-1$). By Lemma~\ref{R1:lem:rep}, for $0\le j\le d$ and $n\ge1$,
\[
a^{(d)}_{n+j}=\int_0^\infty e^{-s}(\alpha_j\ell_n+\beta_j\ell_n')^d\,ds
=\sum_{i=0}^{d}\binom{d}{i}\sum_{k\ge0}[s^k]\bigl(\alpha_j^{d-i}\beta_j^{i}\bigr)\,\RoneJ_n(i,k),
\]
a finite sum. By Lemma~\ref{R1:lem:reduction},
$a^{(d)}_{n+j}=\sum_{b=1-d}^{0}V_{j,b}(n)\RoneJ_n(0,b)$ with $V_{j,b}\in\RoneQ(n)$ whose
denominators divide $\bigl((n+1)\cdots(n+d)\bigr)^d$; these identities hold for all $n\ge1$.
The $(d+1)\times d$ matrix $V=(V_{j,b})$ over the field $\RoneQ(n)$ has a nonzero left kernel
vector; multiplying it by a suitable nonzero element of $\RoneQ(n)$ we get
$c\in\mathbb{Z}[n]^{d+1}\setminus\{0\}$ with $\sum_jc_jV_{j,b}=0$ in $\RoneQ(n)$ for every $b$.
Evaluating at an integer $n\ge1$ (where all denominators of $V$ are nonzero),
$\sum_jc_j(n)a^{(d)}_{n+j}=\sum_b\bigl(\sum_jc_j(n)V_{j,b}(n)\bigr)\RoneJ_n(0,b)=0$.
\end{proof}

\begin{remark}[what the method is]\label{R1:rem:method}
Heuristically (this is not used anywhere), the count $d$ is a de Rham cohomology count for
the connection $e^{-s}\otimes\mathrm{Sym}^d(\text{Laguerre})$ on $\mathbb{G}_m$: the Laguerre
connection of Lemma~\ref{R1:lem:ell}(b) is regular at $s=0$ with a single unipotent Jordan
block and has exponential factors $1,e^{-s}$ at $\infty$, so $e^{-s}\mathrm{Sym}^d$ has
irregularity $d+1$ at $\infty$ and $H^1$ has dimension $d+1$; the integrals start at the
boundary point $s=0$, and Step~4 of the proof shows that the class of $\RoneJ_n(0,-d)$ never
contributes. The shift $n\mapsto n+1$ acts on the remaining $d$-dimensional space, which
bounds the order. This is the
mechanism of the banana upper bound of \cite{K-032} (method~M3), transported from a Laplace variable
to the discrete variable $n$. The proof above is the elementary version of this count and
does not use $\mathcal D$-module theory.
\end{remark}

\begin{remark}[computations]\label{R1:rem:computations}
The script \texttt{rook\_\allowbreak{}verify.py} checks Lemmas~\ref{R1:lem:rep}--\ref{R1:lem:ibp} exactly
on many instances, and that for $2\le d\le7$ and $n=1,\dots,12$ the left kernel of $V(n)$ is
one-dimensional, the coefficients $c_0(n),c_d(n)$ are nonzero (so the relation has order
exactly $d$), and the relation annihilates the exact rook numbers. For $2\le d\le5$ the
constructive relation is proportional, at $n=1,\dots,40$, to the minimal-degree order-$d$
recurrence recomputed by guessing (degrees $1,4,9,18$, as in \cite[\S5]{R1-KZ}).
The script \texttt{rook\_\allowbreak{}degrees\_\allowbreak{}modp.py} computes the constructive recurrence modulo
$p=2^{61}-1$ by evaluation and interpolation, without using the numbers $a^{(d)}_n$; for
every $2\le d\le12$ its minimal coefficient degree is $1,4,9,18,31,50,75,108,149,200,261$,
exactly the degrees listed in \cite[\S5]{R1-KZ}. For $8\le d\le12$ this is evidence, not proof,
that the constructive
recurrence \emph{is} the one of \cite{R1-KZ}: if the recurrence of \cite{R1-KZ} is the only one
of its order and degree fitting the initial terms used to guess it, and if the degree of the
constructive recurrence over $\RoneQ$ equals its degree modulo $p$, then the two agree up to a
constant factor. For $d\le7$ the identification is proved in
Proposition~\ref{R1:prop:soft} below.
\end{remark}

\subsection{A soft proof of the guessed recurrences for \texorpdfstring{$d\le7$}{d<=7}}\label{R1:sec:soft}

Kauers and Zeilberger hoped that an upper bound would yield a ``soft'' proof of the guessed
recurrences. The proof of Theorem~\ref{R1:thm:upper} does this once one bounds degrees. Let
$P_j(n)=(n+1)(n+2)\cdots(n+j)$ ($P_0=1$), and let $V_{j,b}\in\RoneQ(n)$ ($0\le j\le d$,
$1-d\le b\le0$) be the rational functions produced in that proof from the polynomials of
Lemma~\ref{R1:lem:reduction}, so that $a^{(d)}_{n+j}=\sum_bV_{j,b}(n)\RoneJ_n(0,b)$ for $n\ge1$.

\begin{lemma}[degree bound]\label{R1:lem:degbound}
Put $e=\lceil d/2\rceil$. Then $P_j^dV_{j,b}$ is a polynomial in $n$ of degree at most
$\lfloor d/2\rfloor+e\,dj$. Consequently, for every $K=(K_0,\dots,K_d)\in\RoneQ[n]^{d+1}$ with
$\deg K_j\le D$, each
\[
N_b(n)=P_d(n)^d\sum_{j=0}^{d}K_j(n)V_{j,b}(n)\qquad(1-d\le b\le0)
\]
is a polynomial of degree at most $B_d(D)=D+\lfloor d/2\rfloor+e\,d^2$.
\end{lemma}

\begin{proof}
We follow degrees in $n$ through the proof of Lemma~\ref{R1:lem:reduction}. In Step~1,
$\lambda^{(i,k)}_b$ has degree $\le\lfloor i/2\rfloor$, by induction on $i$: the recursion
multiplies by $n$ only in the term $(i-1)n\RoneJ_n(i-2,k-1)$, and
$1+\lfloor(i-2)/2\rfloor=\lfloor i/2\rfloor$. In Step~2 the coefficients come from
$\RoneJ_n(d,k)$, $k\RoneJ_n(d,k-1)$ and $dn\RoneJ_n(d-1,k-1)$, so $\mu^{(k)}_b$ has degree
$\le\max(\lfloor d/2\rfloor,1+\lfloor(d-1)/2\rfloor)=e$. By Step~3 and induction on $k$,
$\RoneJ_n(0,k)=\sum_b\gamma^{(0,k)}_b\RoneJ_n(0,b)$ with $\deg\gamma^{(0,k)}_b\le ek$ for $k\ge1$
(and $\gamma^{(0,k)}_b$ constant for $k\le0$); with Step~4,
$\deg\gamma^{(i,k)}_b\le\lfloor i/2\rfloor+ek$ for $k\ge0$.
Next, by induction on $j$ with Lemma~\ref{R1:lem:ell}(c),
$P_j\ell_{n+j}=A_j\ell_n+B_j\ell_n'$ and $P_j\ell_{n+j}'=C_j\ell_n+D_j\ell_n'$ with
$A_j,B_j,C_j,D_j\in\RoneQ[n,s]$ of total degree $\le j$ (start with $A_0=D_0=1$,
$B_0=C_0=0$; then $A_{j+1}=(n+j+s)A_j+sC_j$, $B_{j+1}=(n+j+s)B_j+sD_j$,
$C_{j+1}=(n+j+1)(A_j+C_j)$, $D_{j+1}=(n+j+1)(B_j+D_j)$). Hence
$P_j^dV_{j,b}=\sum_{i=0}^{d}\binom di\sum_{k=0}^{dj}[s^k]\bigl(A_j^{d-i}B_j^{i}\bigr)\,
\gamma^{(i,k)}_b$, where $[s^k](A_j^{d-i}B_j^i)$ has degree $\le dj-k$ in $n$; the total is
$\le\max_{i,k}(dj-k+\lfloor i/2\rfloor+ek)\le\lfloor d/2\rfloor+e\,dj$ because $e\ge1$ and
$k\le dj$. Finally $N_b=\sum_jK_j\,(P_d/P_j)^d\,P_j^dV_{j,b}$ has degree at most
$D+\max_j\bigl(d(d-j)+\lfloor d/2\rfloor+e\,dj\bigr)=D+\lfloor d/2\rfloor+e\,d^2$.
\end{proof}

\begin{proposition}[the guessed recurrences for $d\le7$]\label{R1:prop:soft}
Let $2\le d\le7$, $D_d=1,4,9,18,31,50$ and $N_d=48,63,94,159,270,455$ respectively. Up to a
constant factor
there is exactly one nonzero $K\in\RoneQ[n]^{d+1}$ with $\deg K_j\le D_d$ and
$\sum_{j=0}^dK_j(n)a^{(d)}_{n+j}=0$ for $1\le n\le N_d-d$, and this $K$ satisfies
$\sum_jK_j(n)a^{(d)}_{n+j}=0$ for \emph{all} $n\ge0$. In particular, for $d=4,5,6,7$ the
recurrences of order $d$ and coefficient degree $9$, $18$, $31$, $50$ conjectured in
\cite{R1-KZ} are proved, provided they are satisfied by the terms
$a^{(d)}_1,\dots,a^{(d)}_{N_d}$ (any recurrence of that order and degree satisfied by these
terms equals $K$ up to a constant factor; we have not compared coefficients with the published
files \cite{R1-KZfiles}).
\end{proposition}

\begin{proof}
This is a computer-assisted proof (script \texttt{rook\_\allowbreak{}softproof.py}); all arithmetic below is
exact except the rank computation modulo a prime, which is only used as a lower bound for
the rank over $\RoneQ$. The homogeneous linear system (equations $1\le n\le N_d-d$) for the
coefficients of $K$ has a one-dimensional solution space modulo a prime $p$; since a nonzero
minor modulo $p$ comes from a nonzero integer minor, the solution space over $\RoneQ$ has
dimension at most one. A nonzero solution $K$ is obtained by Chinese remaindering and rational
reconstruction (five primes for $d\le5$, seven for $d=6$, twelve for $d=7$; the integer
coefficients of the primitive $K$ have at most $4,18,42,104,175,317$ bits) and is checked to
satisfy the equations for $0\le n\le N_d-d$ exactly, so the dimension is one. By
Lemma~\ref{R1:lem:degbound}, each $N_b$ is a polynomial of degree at most
$B_d(D_d)=6,23,43,95,142,249$ for $d=2,\dots,7$. The script
evaluates the reduction of Lemma~\ref{R1:lem:reduction} and the shift of
Theorem~\ref{R1:thm:upper} at $n=n_0$ with exact rational arithmetic (evaluation commutes with
these operations, which only divide by nonzero integers and by $P_j(n_0)\ne0$) and checks
$\sum_jK_j(n_0)V_{j,b}(n_0)=0$ for $n_0=1,\dots,B_d(D_d)+1$ and all $b$. Hence $N_b\equiv0$,
and for every $n\ge1$,
$\sum_jK_j(n)a^{(d)}_{n+j}=\sum_b\bigl(\sum_jK_j(n)V_{j,b}(n)\bigr)\RoneJ_n(0,b)=0$; the case
$n=0$ was checked directly above.
\end{proof}

\begin{remark}\label{R1:rem:softlarger}
The same procedure applies verbatim to any $d$ and any candidate recurrence. For
$8\le d\le12$ it needs the candidate (recomputed by guessing, or taken from the files
\cite{R1-KZfiles}) and $B_d(D_d)+1=336,518,655,932,1132$ exact evaluations of the reduction;
we did not carry this out for $d\ge8$.
\end{remark}

\subsection{The large-dimension asymptotics}\label{R1:sec:largeD}

\begin{theorem}\label{R1:thm:largeD}
For every fixed $n\ge1$,
$w_n(d)=a^{(d)}_n=e^{n-1}\dfrac{(nd)!}{n!^d}\bigl(1+O(\tfrac1d)\bigr)$ as $d\to\infty$.
\end{theorem}

\begin{proof}
For $n=1$, $a^{(d)}_1=d!$ exactly. Let $n\ge2$. By \eqref{R1:eq:ell},
$\ell_n(s)=\frac{s^n}{n!}P(1/s)$ with
$P(u)=\sum_{j=0}^{n-1}\binom{n-1}{j}\frac{n!}{(n-j)!}u^j$, so $P(0)=1$, $P'(0)=n(n-1)$,
$P$ has nonnegative coefficients, and $1\le P(u)\le(1+nu)^{n-1}$ for $u\ge0$ (because
$n!/(n-j)!\le n^j$). Put $N=nd$ and let $S$ have the Gamma density $e^{-s}s^N/N!$ on
$(0,\infty)$. Lemma~\ref{R1:lem:rep} gives
\[
\frac{n!^d}{(nd)!}a^{(d)}_n=\mathbb{E}\bigl[P(1/S)^d\bigr].
\]
\emph{Bulk.} On $[0,1]$ the function $\log P$ is smooth with $\log P(u)=n(n-1)u+r(u)$,
$|r(u)|\le C_nu^2$. On the event $S\ge N/2$ (and $d\ge4/n$, so $1/S\le1$) we get
$d\log P(1/S)=(n-1)N/S+\epsilon$ with $|\epsilon|\le4C_n/(n^2d)$. Write $S=N(1+X)$; then
$\mathbb{E}X=1/N$, $\mathbb{E}X^2=(N+2)/N^2$, and on $X\ge-1/2$ the function
$\phi(X)=\exp\bigl((n-1)(\frac1{1+X}-1)\bigr)$ satisfies $|\phi(X)-1+(n-1)X|\le K_nX^2$.
Since $\mathbb{P}(S<N/2)\le e^{-cN}$ and $\mathbb{E}[|X|;S<N/2]\le e^{-cN}$ for some $c>0$
(Chernoff bound for the Gamma law),
$\mathbb{E}[e^{(n-1)N/S};S\ge N/2]=e^{n-1}\bigl(1+O(1/N)\bigr)$, and therefore
$\mathbb{E}[P(1/S)^d;S\ge N/2]=e^{n-1}\bigl(1+O(1/d)\bigr)$.

\emph{Small $S$.} Using $P(1/s)^d\le(1+n/s)^{(n-1)d}$ and $N-(n-1)d=d$,
\begin{align*}
\mathbb{E}[P(1/S)^d;S<N/2]&\le\frac1{N!}\int_0^{N/2}e^{-s}s^{d}(s+n)^{(n-1)d}ds\\
&\le\frac{e^{n}}{N!}\int_{n}^{N/2+n}e^{-v}v^{N}dv\le e^{n}\,\mathbb{P}(S\le N/2+n),
\end{align*}
which is $O(e^{-c'N})$. Adding the two parts proves the theorem.
\end{proof}

\begin{remark}[second term; numerical only]\label{R1:rem:secondterm}
Expanding one order further ($\log P(u)=n(n-1)u-n(n-1)^2u^2+O(u^3)$ and
$\mathrm{Var}(S/N)\sim1/N$) suggests
$a^{(d)}_n=e^{n-1}\frac{(nd)!}{n!^d}\bigl(1-\frac{(n-1)^2}{2nd}+O(d^{-2})\bigr)$. Exact values
for $n=2,3,4$ and $d\le80$ agree (\texttt{rook\_\allowbreak{}verify.py}, part 6: $d(\text{ratio}-1)\to
-\tfrac14,-\tfrac23,-\tfrac98$). We did not write out the error analysis of this second term.
Combinatorially, $a^{(d)}_n=\sum_w2^{A(w)}$ over the arrangements $w$ of the multiset
$\{1^n,\dots,d^n\}$, where $A(w)$ counts adjacent equal letters (split each run into a
composition); this explains the factor $e^{n-1}=\mathbb{E}[2^{\mathrm{Poisson}(n-1)}]$.
\end{remark}

\subsection{Lower bounds: what is proved and what remains}\label{R1:sec:lower}

The following refinement of the singularity budget used in \cite{K-006,K-034} (Franel order,
method~M1) counts local monodromy with multiplicity.

\begin{lemma}[index budget]\label{R1:lem:index}
Let $(a_n)$ satisfy a nonzero recurrence $\sum_{j=0}^{r}p_j(n)a_{n+j}=0$ ($n\ge n_0$,
$p_j\in\RoneC[n]$), let $F(t)=\sum a_nt^n$ have positive radius of convergence, let $V$ be
the solution space (at a base point) of the minimal differential operator of $F$, and for
$p\in\RoneC^\times$ let $h_p$ be the dimension of the space of local solutions at $p$ (of
that operator, continued to a punctured disc at $p$) that extend holomorphically to $p$.
Then
\[
\sum_{p\in\RoneC^\times}\bigl(\dim V-h_p\bigr)\le r,\qquad\text{and}\qquad
\dim V-h_p\ge\operatorname{rk}(T_p-1),
\]
where $T_p$ is the local monodromy at $p$.
\end{lemma}

\begin{proof}
As in the Franel argument of \cite{K-006}, with $\vartheta=t\,d/dt$ the operator
$L_0=\sum_jt^{r-j}p_j(\vartheta-j)$ satisfies $L_0F=R$ with $R$ a polynomial, so
$L=\partial_t^{\,1+\deg R}L_0$ annihilates $F$. Write $L=\sum_{k\le D'}a_k\partial_t^k$ with
polynomial coefficients; its leading coefficient is $t^{D}H(t)$ with $\deg H\le r$ (the
leading coefficient of $L_0$ is $t^DH(t)$, $H=\sum_jc_jt^{r-j}$, $c_j$ the top coefficient of
$p_j$, and left multiplication by $\partial_t^e$ does not change it). Fix
$p\in\RoneC^\times$ and let $\mu_p=\operatorname{ord}_pH$. By Malgrange's index theorem
\cite{R1-Malgrange}, the index of $L$ acting on convergent power series at $p$ equals
$D'-\operatorname{ord}_pa_{D'}=D'-\mu_p$ (the index on formal power series is
$\max_k(k-\operatorname{ord}_pa_k)$; the difference is the irregularity, which we do not need);
since the dimension of the kernel is at least the index, $L$ has at least $D'-\mu_p$ linearly
independent solutions holomorphic at $p$. The minimal operator divides $L$ on the right, so
near $p$ its solution space $V$ is contained in the $D'$-dimensional solution space of $L$;
therefore $h_p\ge\dim V+(D'-\mu_p)-D'$, i.e.\ $\dim V-h_p\le\mu_p$. Summing over $p$ gives
$\sum_p(\dim V-h_p)\le\deg H\le r$. Holomorphic solutions are fixed by $T_p$, which gives the
second inequality.
\end{proof}

For the rook diagonal $F_d(t)=\sum_na^{(d)}_nt^n$, Lemma~\ref{R1:lem:index} reduces sharpness
to local monodromy. Put $t_m=\bigl(\frac{m+1}{d+1-m}\bigr)^{d-2m}$ for $0\le m<d/2$; these
are the critical values of $\prod x_i$ on the polar divisor at the points with $d-m$
coordinates equal to $r_m=\frac{m+1}{d+1-m}$ and $m$ equal to $1/r_m$ (Lagrange multipliers).

\begin{conjecture}[sharpness; Kauers--Zeilberger's empirical statement]\label{R1:conj:sharp}
For every $d\ge1$, $a^{(d)}_n$ satisfies no recurrence of order $<d$. More precisely, the
minimal operator of $F_d$ has nontrivial local monodromy at each $t_m$ ($0\le m<d/2$) and
$\operatorname{rk}(T_1-1)=\lfloor d/2\rfloor$ at $t=1$, so that Lemma~\ref{R1:lem:index}
gives order $\ge\lceil d/2\rceil+\lfloor d/2\rfloor=d$; and the characteristic polynomial of
the order-$d$ recurrence is $(T-1)^{\lfloor d/2\rfloor}\prod_{0\le m<d/2}(T-t_m^{-1})$.
\end{conjecture}

\begin{proposition}\label{R1:prop:small}
Conjecture~\ref{R1:conj:sharp} (first sentence) holds for $d\le3$. Hence for $d\le3$ the
minimal order is exactly $d$.
\end{proposition}

\begin{proof}
A recurrence of order $0$ is impossible since $a^{(d)}_n\neq0$; this settles $d=1$
($a^{(1)}_n=2^{n-1}$ for $n\ge1$). By Lemma~\ref{R1:lem:index} it suffices, for $d=2,3$, to
exhibit $d$ distinct points of $\RoneC^\times$ at which some continuation of $F_d$ is not
holomorphic.

$d=2$: $F_2=\frac12\bigl(1+\sqrt{(1-t)/(1-9t)}\bigr)$, which has branch points at $t=1/9$ and
$t=1$. (This is classical: $F_2(t)$ is the sum of the residues of
$R_2(x,t/x)/x=\frac{(1-x)(x-t)}{x(-2x^2+(1+3t)x-2t)}$ at its poles that tend to $0$ with $t$,
namely $\frac12$ at $x=0$ and $\frac12\sqrt{(1-t)/(1-9t)}$ at the small root of
$2x^2-(1+3t)x+2t$.)

$d=3$: by \cite{R1-BCvHP} (van Hoeij's formula),
$F_3(t)=1+\int_0^t g(w)\,dw$ with $g(w)=\frac{6\,G(z(w))}{(1-4w)(1-64w)}$,
$G={}_2F_1(\frac13,\frac23;2;\cdot)$ and $z(w)=\frac{27w(2-3w)}{(1-4w)^3}$; one checks
$1-z(w)=\frac{(1-w)^2(1-64w)}{(1-4w)^3}$. At $w=1/4$ every branch of $g$ is analytic
(the local solutions of the hypergeometric equation at $z=\infty$ are $z^{-1/3}$ and
$z^{-2/3}$ times power series in $1/z$, and $z^{-1}$ has a triple zero at $w=1/4$), so the
only possible singularities of branches of $F_3$ in $\RoneC^\times$ are $1/64$, $2/3$, $1$.
The preimage $z^{-1}([1,\infty))$ is the segment $[\frac1{64},\frac14]$ together with two arcs
from $\frac14$ to $1$ (near $w=1$, $1-z\sim\frac73(w-1)^2$, so the arcs arrive vertically),
which bound a lens containing $(\frac14,1)$; the complementary component containing $0$ is
unbounded and contains $(1,\infty)$. On it $G\circ z$ stays on the principal branch.
(i) As $w\uparrow\frac1{64}$, $z\uparrow1$ and $G(z)\to G(1)=\frac{\Gamma(2)\Gamma(1)}
{\Gamma(5/3)\Gamma(4/3)}\neq0$; $g$ has a simple pole and $F_3$ a logarithmic singularity at
$1/64$. (ii) Since $c-a-b=1$, near $z=1$ the principal branch is
$G=A(z)+B(1-z)\log(1-z)\,h(z)$ with $A,h$ analytic, $h(1)=1$ and
$B=\Gamma(2)/(\Gamma(\frac13)\Gamma(\frac23))\neq0$ \cite[15.8.10]{R1-DLMF}. Approaching
$w=1$ along $(1,1+\delta)$ inside the unbounded component, $z\uparrow1$ and $g$ contains
$\frac73 B\,(w-1)^2\log((w-1)^2)$ times a nonvanishing analytic factor, so this
continuation of $F_3$ is singular at $1$. (iii) Crossing one of the arcs once replaces $G$
by $G+\kappa y_1$, $\kappa=\pm2\pi iB\neq0$, where $y_1=(1-z)h(z)$ is the exponent-$1$
solution at $z=1$. Inside the lens we may follow the real axis to $w\to\frac23$, where $z$
has a simple zero. The solutions analytic at $z=0$ form the line $\RoneC G$ (the exponents
there are $0,-1$), and $y_1$ is not a multiple of $G$ on $(0,1)$ because $y_1(1)=0\neq G(1)$;
hence $y_1$ has a nonzero pole part $\beta z^{-1}$ at $0$, and $\beta$ does not depend on the
side from which $0$ is approached. Thus $g$ acquires a simple pole at $\frac23$ and this
continuation of $F_3$ has a logarithmic singularity there. Three points give order $\ge3$.
\end{proof}

\begin{remark}[evidence for Conjecture~\ref{R1:conj:sharp}]\label{R1:rem:evidence}
For $d\le5$ we reconstructed (by guessing, exactly over $\RoneQ$) differential operators
annihilating $F_d$ that are minimal among those of small degree; they have orders $2,3,5,8$
($d=2,\dots,5$; order $11$ for $d=6$, found modulo a prime), and for $d\le4$ they annihilate
constants. Their exact local analysis
(\texttt{explore\_\allowbreak{}t1.py}, \texttt{explore\_\allowbreak{}jordan.py}) gives: at each $t_m$ a pseudo-reflection
(exponents $0,\dots,r-2$ and one exponent in $\frac12+\mathbb{Z}$ for even $d$, a logarithm for
odd $d$); at $t=1$, $\dim\ker(T_1-1)=r-\lfloor d/2\rfloor$ (for $d=2$ the eigenvalue $-1$ occurs once),
with unipotent Jordan blocks $(2,1)$, $(3,1,1)$, $(2,2,1,1,1,1)$ for $d=3,4,5$; and the leading coefficient vanishes to order exactly $1$ at each $t_m$ and
$\lfloor d/2\rfloor$ at $t=1$. The recurrences of Proposition~\ref{R1:prop:soft} (which are
proved) have characteristic polynomial exactly $(T-1)^{\lfloor d/2\rfloor}\prod_m(T-t_m^{-1})$
up to a constant, for $2\le d\le6$ (\texttt{rook\_\allowbreak{}charpoly.py}), and no recurrence of order
$<d$ with coefficient degree up to $20$ ($d=4$), resp.\ $35$ ($d=5$), fits the data. The
differential operators are not certified as minimal, so all this is evidence, not proof. What is missing for a proof is global:
one must show that the continuations of $F_d$ realize the singular local solutions at every
$t_m$ and a $\lfloor d/2\rfloor$-dimensional image of $T_1-1$ (for odd $d$ the Jordan
blocks at $t=1$ have size $2$, so no single branch suffices). For $d=4$ it would suffice to
know that the guessed operator $M_4=L_1\partial_t$ (order $5$) annihilates $F_4$, which can be
proved like Proposition~\ref{R1:prop:soft}, and that $L_1$ is irreducible. Indeed, the
derivatives of the continuations of $F_4$ would then span the solution space of $L_1$; $L_1$
has $3,3,2$ linearly independent holomorphic solutions at $\frac1{625},\frac14,1$
(\texttt{explore\_\allowbreak{}sharp4b.py}), and the argument of Lemma~\ref{R1:lem:index} would give
order $\ge1+1+2=4$. The exponents of
$L_1$ are $\{0,0,-1,-1\}$ at $0$, $\{0,1,2,-\frac12\}$ at $\frac1{625}$ and at $\frac14$,
$\{0,1,1,3\}$ at $1$, $\{1,1,2,2\}$ at $\infty$ (solutions $\sim t^{-e}$), and $\{0,1,2,4\}$ at
four apparent singularities. The Fuchs relation therefore does not exclude right factors of
orders $1,2,3$, so the cheap irreducibility test is inconclusive, and a factorization algorithm
or a certified monodromy computation would be needed.
\end{remark}


%% ==================================================================
%% Section: Division points and HCM laws
%% ==================================================================
% R7 section: transcendence/division points, drawn curves, and HCM classification questions.
% Macros (unique prefix Rvii to avoid clashes with other sections).
\providecommand{\RviiQ}{\mathbb{Q}}
\providecommand{\RviiR}{\mathbb{R}}
\providecommand{\RviiC}{\mathbb{C}}
\providecommand{\RviiZ}{\mathbb{Z}}
\providecommand{\RviiO}{\mathcal{O}}
\providecommand{\RviiK}{\mathcal{K}}
\providecommand{\RviiIm}{\operatorname{Im}}
\providecommand{\RviiRe}{\operatorname{Re}}
\providecommand{\RviiHCM}{\mathrm{HCM}}
\providecommand{\RviiHM}{\mathrm{HM}}
\providecommand{\RviiGGC}{\mathrm{GGC}}
\providecommand{\RviiGam}{\boldsymbol{\Gamma}}
\providecommand{\RviiB}{\mathbf{B}}

\section{Division points, drawn curves, and hyperbolic complete monotonicity}
\label{R7-sec}

This section settles five explicitly posed questions by other authors.
\begin{enumerate}
\item The speculation of Mani and Rubinstein-Salzedo on individual
  constructible division points of the tricuspoid is false; we give the
  complete classification (Theorem~\ref{R7-thm-tricuspoid}).
\item Their question on the lemniscate \emph{with the curve drawn} is answered
  by combining known results: a drawn lemniscate gives exactly the origami
  field (as for any drawn non-circular conic, by Videla and Baek et al.), and
  the admissible $n$ are those of the Cox--Shurman origami theorem
  (Theorem~\ref{R7-thm-drawn}).
\item The identification $\RviiHM\cap\RviiGGC=\RviiHCM$ asked about by
  Jedidi and Simon fails (Theorem~\ref{R7-thm-JS}).
\item Bosch's conjecture that $\sqrt{\RviiGam_t\RviiGam_s}$ is always HCM is
  false; in fact it is HCM if and only if $|t-s|\le 1/2$
  (Theorem~\ref{R7-thm-bosch}); the $n$-fold version fails for every
  $n\ge2$ (Proposition~\ref{R7-prop-nfold}).
\item The belief of Bosch and Simon that $a+b\ge c$ is necessary for
  $\RviiGam_c\times\RviiB_{a,b}$ to be HCM is true
  (Theorem~\ref{R7-thm-gammabeta}).
\end{enumerate}
Items 1--2 are division-point dichotomies in the spirit of the
algebraic/transcendental division classifications of method~M8, carried out
in the constructible (two-power) setting. Items 3--5 use method~M9's
mechanism: a Bernstein (positive Laplace) representation forces sign
conditions, here in the form of a zero-free and Pick-type obstruction
(Lemma~\ref{R7-lem-pick}), which sharpens the zero-free continuation
argument used for McKay laws in \cite{K-029}.

Throughout, ``constructible'' means constructible with straightedge and
compass from the points $0$ and $1$ of $\RviiC=\RviiR^2$, unless a drawn curve
is explicitly allowed.

%=====================================================================
\subsection{Individual division points of the tricuspoid}
\label{R7-ssec-tricuspoid}

\begin{problem}[{Mani and Rubinstein-Salzedo \cite[Remark~5.7]{R7-MRS16}; Remark~4.3 of the arXiv version}]
\label{R7-prob-MRS-tri}
``One may be tempted to wonder which, if any, of the individual
$n$-division points are constructible when $n\nmid 6$. It appears that no
$n$-division points other than those contained among the $6$-division points
are constructible. The $x$-coordinates of $n$-division points are roots of
the polynomials $f_r(x)$, where $r$ is a rational number $\ge 3$ whose
numerator divides $n$. It appears that, except in the cases already
discussed, $f_r(x)$ is indeed irreducible, and further, that the field
$\RviiQ(\alpha)$ obtained by adjoining a root of $f_r(x)$ is totally ramified
at $3$.''
\end{problem}

In \cite{R7-MRS16} the tricuspoid (deltoid) is
$x=2\cos\varphi+\cos2\varphi$, $y=2\sin\varphi-\sin2\varphi$, i.e.
\[
 z(\varphi)=2e^{i\varphi}+e^{-2i\varphi},\qquad 0\le\varphi<2\pi ,
\]
the curve is not drawn, and the $n$-division points are taken relative to
the cusp $z(0)=3$: they are the points $P(k/n)$, $0\le k<n$, where
$P(r)$ denotes the point at arc length $16r$ from $3$ in the direction of
increasing $\varphi$ (the total length is $16$). Their main theorem for this
curve \cite[Theorem~5.1]{R7-MRS16} says that \emph{all} $n$-division points
are constructible if and only if $n\mid 6$. Put $T_3(X)=4X^3-3X$ and
$\omega=e^{2\pi i/3}$.

\begin{lemma}\label{R7-lem-tri-arc}
$|z'(\varphi)|=4|\sin(3\varphi/2)|$; for $0\le\varphi\le 2\pi/3$ the arc
length from $z(0)$ is $s(\varphi)=\tfrac83\bigl(1-\cos\tfrac{3\varphi}2\bigr)$,
with $s(2\pi/3)=16/3$; and $z(\varphi+2\pi/3)=\omega z(\varphi)$.
\end{lemma}

\begin{proof}
$|z'|^2=|2ie^{i\varphi}-2ie^{-2i\varphi}|^2=8-8\cos3\varphi=16\sin^2(3\varphi/2)$,
and $\sin(3\varphi/2)\ge0$ on $[0,2\pi/3]$; integrate. Finally
$e^{-2i(\varphi+2\pi/3)}=e^{-2i\varphi}e^{-4\pi i/3}=\omega e^{-2i\varphi}$.
\end{proof}

\begin{theorem}[Constructible points of the tricuspoid]
\label{R7-thm-tricuspoid}
Let $r\in\RviiQ\cap[0,1)$ and write $r=j/3+\rho$ with $j\in\{0,1,2\}$ and
$\rho\in[0,1/3)$. The following are equivalent:
\begin{enumerate}
\item[(i)] $P(r)$ is constructible;
\item[(ii)] the cubic $T_3(X)-(1-6\rho)$ has a rational root;
\item[(iii)] $\rho=\tfrac16(1-u)(1+2u)^2$ for some $u\in\RviiQ\cap[-1,1]$.
\end{enumerate}
If these fail, the abscissa of $P(r)\,\omega^{-j}$ generates a cubic field.
\end{theorem}

\begin{proof}
(ii)$\Leftrightarrow$(iii) since $1-T_3(u)=(1-u)(1+2u)^2$, and a rational
root $u$ of $T_3(X)=1-6\rho\in(-1,1]$ lies in $[-1,1]$.

By Lemma~\ref{R7-lem-tri-arc}, $P(r)=\omega^jz(\varphi)$ where
$\varphi\in[0,2\pi/3)$ is defined by $\cos(3\varphi/2)=\tau:=1-6\rho$.
Put $\theta=\varphi/2\in[0,\pi/3)$, $c=\cos\theta$, $s=\sin\theta\ge0$. Then
$T_3(c)=\cos3\theta=\tau$ and
\[
 x:=\RviiRe z(\varphi)=2\cos2\theta+\cos4\theta=8c^4-4c^2-1,\qquad
 y:=\RviiIm z(\varphi)=2\sin2\theta-\sin4\theta=8cs^3 .
\]
Multiplication by $\omega^{\pm1}$ preserves constructibility
($\cos\frac{2\pi}3=-\frac12$, $\sin\frac{2\pi}3=\frac{\sqrt3}2$), so (i) holds
iff $z(\varphi)$ is constructible.

(ii)$\Rightarrow$(i): if $g(X):=T_3(X)-\tau$ has a rational root, all its
roots lie in a field of degree $\le2$ over $\RviiQ$; in particular
$[\RviiQ(c):\RviiQ]\le2$, $s=\sqrt{1-c^2}$, and $x,y\in\RviiQ(c,s)$, a
two-step quadratic tower.

(i)$\Rightarrow$(ii): suppose $g$ is irreducible over $\RviiQ$. Then
$[\RviiQ(c):\RviiQ]=3$ and $x\in\RviiQ(c)$, so $[\RviiQ(x):\RviiQ]\in\{1,3\}$.
Modulo $g$ we have $8c^4=2c\cdot4c^3=6c^2+2\tau c$, hence
$x=2c^2+2\tau c-1$. If $x=x_0\in\RviiQ$, then $c$ would be a root of
$2X^2+2\tau X-1-x_0\in\RviiQ[X]$, contradicting $[\RviiQ(c):\RviiQ]=3$.
So $[\RviiQ(x):\RviiQ]=3$, $x$ is not constructible, and neither is
$z(\varphi)$. A reducible cubic over $\RviiQ$ has a rational root.
\end{proof}

The minimal polynomial of $x$ is the characteristic polynomial of
$2c^2+2\tau c-1$ on $\RviiQ[c]/(g)$, namely
\begin{equation}\label{R7-eq-tri-cubic}
 F_\tau(x)=4x^3-(24\tau^2+3)x+(1-20\tau^2-8\tau^4).
\end{equation}
For $\rho=1/n$ one checks $n^4F_{1-6/n}(x)=4n^4x^3-(27n^4-288n^3+864n^2)x
-(27n^4-432n^3+2448n^2-6912n+10368)$, which is exactly the cubic $f_n(x)$
printed in \cite[Section~5]{R7-MRS16} (Section~4 of the arXiv version); the polynomial $f_r$ of
Problem~\ref{R7-prob-MRS-tri} is $F_{1-6/r}$ up to a constant.

\begin{corollary}[All three speculations in Problem~\ref{R7-prob-MRS-tri} fail]
\label{R7-cor-tri}
\begin{enumerate}
\item[(1)] For $n=32$ the division points $P(k/32)$ with $k\in\{9,13,19,23\}$ are
  constructible, and none of them is a $6$-division point. For instance,
  with $u=1/4$, $T_3(1/4)=-11/16=1-6\cdot\frac{9}{32}$, and
  \[
   P\!\left(\tfrac{9}{32}\right)=\left(\frac{39-45\sqrt5}{64},\
   \sqrt{\frac{3969+1755\sqrt5}{2048}}\right).
  \]
  No $n<32$ has a constructible $n$-division point outside the
  $6$-division points, but infinitely many $n$ do: every
  $u\in\RviiQ\cap(-1,1)$ with $(1-u)(1+2u)^2\notin\RviiZ$ gives a constructible
  point $P(\rho)$, $\rho=\frac16(1-u)(1+2u)^2\in(0,\frac13)\setminus\{\frac16\}$,
  and these $\rho$ have unbounded denominators.
\item[(2)] $f_{32/9}$ is reducible: $8192\,F_{-11/16}(x)=(32x+39)(1024x^2-1248x-2151)$.
\item[(3)] For $r=28/9$ (that is $\rho=9/28$, $\tau=-13/14$), $f_r$ is irreducible
  over $\RviiQ$ and $3$ is unramified in $\RviiQ(\alpha)$.
\end{enumerate}
\end{corollary}

\begin{proof}
(1) For $k=9,13,19,23$ one has $\rho=9/32,\ 7/96,\ 25/96,\ 5/96$ and
$1-6\rho=T_3(u)$ with $u=1/4,\ -3/4,\ 3/4,\ -1/4$; apply
Theorem~\ref{R7-thm-tricuspoid}. The explicit point: the roots of
$T_3(X)+11/16$ are $1/4$ and $(-1\pm3\sqrt5)/8$; the root in $[1/2,1]$ is
$c=(3\sqrt5-1)/8$, and $x=2c^2+2\tau c-1$, $y^2=64c^2(1-c^2)^3$ give the
displayed values. Minimality of $32$ and the remaining claims of (1)--(3)
are finite computations, carried out by \texttt{R7\_\allowbreak{}verify\_\allowbreak{}tricuspoid.py}
(rational-root test, factorization of $F_\tau$, and a PARI computation of
the decomposition of $3$ in the cubic field for (3)).
\end{proof}

\begin{remark}
The script also checks Theorem~\ref{R7-thm-tricuspoid} against an
independent criterion for all $k/n$ with $n\le60$: the order of the Galois
group of the minimal polynomial of $z=x+iy$ is a power of $2$ exactly when
(ii) holds; it confirms $s(\varphi)=4.5$ for $P(9/32)$ by $10^{-40}$-accurate
quadrature, and re-derives \cite[Theorem~5.1]{R7-MRS16} for $n\le2000$
(for $n\nmid6$ the first division point is not constructible).
\end{remark}

%=====================================================================
\subsection{The lemniscate with the curve drawn}
\label{R7-ssec-drawn}

\begin{problem}[{\cite[Remark~1.4]{R7-MRS16}}]\label{R7-prob-drawn}
``For curves other than the circle, there are actually two possible problems
regarding the constructibility of $n$-division points. First, for a closed
curve $C$, we can ask whether, given the curve $C$ already drawn on the
plane, it is possible to construct the $n$-division points using a compass
and a straightedge. However, we can also ask whether it is possible to
construct the $n$-division points of $C$ without $C$ ever being drawn.
[\dots] Abel's Theorem is an answer to a problem of the second type. To the
best of our knowledge, the corresponding question with a drawn lemniscate has
not yet been answered.''
\end{problem}

Let $L=\{(x^2+y^2)^2=x^2-y^2\}$, with node $O=0$ and vertex $A=1$. Let
$\RviiK_L\subset\RviiC$ be the set of points constructible from $O,A$ when, in
addition to intersecting lines and circles, one may intersect a constructed
line or circle with the drawn curve $L$. Let $\RviiO$ be the field of origami
numbers; by Alperin \cite[Theorem~5.2]{R7-Alperin00} it is the smallest
subfield of $\RviiC$ closed under square roots, cube roots and complex
conjugation, and $z\in\RviiO$ iff $z$ lies in a tower of extensions of degrees
$2$ and $3$ over $\RviiQ$ \cite[Section~10.3]{R7-CoxGalois}.

\begin{theorem}\label{R7-thm-drawn}
Let $\RviiK_L$ be the set of points constructible with straightedge and
compass from $O$ and $A$ when intersections of constructed lines and circles
with the drawn lemniscate $L$ are also allowed. Then $\RviiK_L=\RviiO$. Consequently the $n$-division points of the lemniscate
(relative to $O$, or to any point of $L\cap\RviiK_L$) are constructible with
straightedge and compass from the drawn lemniscate if and only if
\[
 n=2^a3^bp_1\cdots p_r,\qquad a,b\ge0,
\]
where $p_1,\dots,p_r$ are distinct Pierpont primes ($p_i=2^u3^v+1$) with
$p_i=7$ or $p_i\equiv1\pmod 4$. For instance $n=7,9,13,21,37$ are possible
with the drawn curve but not without it (Abel), while $n=11,19,23,25$ remain
impossible.
\end{theorem}

\begin{proof}
\emph{Inversion.} The map $\iota(z)=1/\bar z$ is a straightedge-and-compass
operation on points, maps lines and circles to lines and circles
(constructibly), and maps $L\setminus\{0\}$ onto the hyperbola
$H=\{X^2-Y^2=1\}$: if $x+iy=\iota(X+iY)$ then $x^2+y^2=(X^2+Y^2)^{-1}$ and
$x^2-y^2=(X^2-Y^2)(X^2+Y^2)^{-2}$. Hence intersecting a constructed circle
$\Gamma$ with $H$ can be done by intersecting $\iota(\Gamma)$ with $L$ and
applying $\iota$.

\emph{Real cubics.} Let $p(w)=w^3+\beta_2w^2+\beta_1w+\beta_0$ with real
$\beta_i\in\RviiK_L$, $\beta_0\ne0$, and let $w_0$ be a real root. Put
$q(w)=(w-m)p(w)$ with $m=-\operatorname{sign}\beta_0$; then
$q=w^4+a_3w^3+a_2w^2+a_1w+a_0$ with $a_0=|\beta_0|>0$. With
$\lambda=(4a_0)^{-1/4}\in\RviiK_L$ the substitution $u=\lambda w$ turns
$q=0$ into $u^4+Du^3+Fu^2+\tfrac E2u+\tfrac14=0$, where $D=a_3\lambda$,
$F=a_2\lambda^2$, $E=2a_1\lambda^3$. In the rotated coordinates
$u=(X+Y)/\sqrt2$, $v=(X-Y)/\sqrt2$ the hyperbola is $uv=\frac12$, and
substituting $v=1/(2u)$ shows that the circle
$u^2+v^2+Du+Ev+F=0$ meets it exactly at the points $(u,1/(2u))$ with $u$ a
real root of this quartic. As $\lambda w_0$ is such a root, the circle is
real, it is constructible, and the intersection point is in $\RviiK_L$.
Hence $w_0\in\RviiK_L$.

\emph{Lower bound.} $\RviiK_L$ is a subfield of $\RviiC$ closed under
conjugation and square roots (classical). If $z=\varrho e^{i\vartheta}\in
\RviiK_L$, then $\varrho^{1/3}$ (root of $w^3-\varrho$) and
$\cos(\vartheta/3)$ (root of $4w^3-3w-\cos\vartheta$) lie in $\RviiK_L$, hence
so do $\sin(\vartheta/3)$ and $z^{1/3}$. By Alperin's characterization,
$\RviiO\subseteq\RviiK_L$.

\emph{Upper bound.} Suppose all points constructed so far have coordinates in
a subfield $K\subset\RviiO$ (with $i\in K$). Intersections of lines and circles
give coordinates in a quadratic extension of $K$. For a line
$\alpha x+\beta y+\gamma=0$ over $K$, eliminating $y$ from the equation of $L$
gives a quartic in $x$. For a circle $x^2+y^2+Dx+Ey+F=0$, points of $L$ on it
satisfy $(Dx+Ey+F)^2=x^2-y^2$, so they are intersections of the circle with a
conic, and their abscissae are roots of a polynomial of degree $\le4$ over
$K$ (Bézout); then $y$ satisfies a quadratic over $K(x)$. Roots of quartics
are obtained by Ferrari's formulas with square and cube roots, so the new
coordinates lie in $\RviiO$. By induction $\RviiK_L\subseteq\RviiO$.

\emph{Division points.} By the first part, a point is constructible from the
drawn lemniscate iff it is origami constructible. Cox and Shurman proved that
the lemniscate can be divided into $n$ arcs of equal length by origami iff $n$
has the displayed form \cite[Lemniscate Theorem, part (2)]{R7-CoxShurman05}.
The lemniscatic addition formula
$\mathrm{sl}(u+v)=(\mathrm{sl}\,u\,\mathrm{sl}'v+\mathrm{sl}\,v\,\mathrm{sl}'u)/(1+\mathrm{sl}^2u\,\mathrm{sl}^2v)$,
together with $x=r\sqrt{(1+r^2)/2}$, $y=r\sqrt{(1-r^2)/2}$ for the point of
$L$ at polar radius $r$, shows that changing the base point to another
constructible point of $L$ does not change constructibility.
\end{proof}

\begin{remark}
The arithmetic behind the theorem is that, for odd $n$, the Galois group
over $\RviiQ(i)$ of the $n$-division value $\mathrm{sl}(2\varpi/n)$ is
$(\RviiZ[i]/n)^\times$ \cite[Corollary~4.8]{R7-CoxHyde14} (powers of $2$ only
add $2$-power factors), whose order
$\prod_{p^e\| n}|(\RviiZ[i]/p^e)^\times|$ is $3$-smooth exactly for the
displayed $n$. \texttt{R7\_\allowbreak{}verify\_\allowbreak{}drawn\_\allowbreak{}lemniscate.py} checks this
equivalence for $n\le 2\cdot10^4$, confirms for odd $n\le13$ by factoring the
division polynomials of $y^2=x^3-x$ over $\RviiQ(i)$ that the division value
$\wp(2\varpi/n)=\mathrm{sl}(2\varpi/n)^{-2}$ has degree
$|(\RviiZ[i]/n)^\times|/2$, and carries out the constructions of $2^{1/3}$ and
$\cos(\pi/9)$ with the drawn lemniscate numerically (error $<10^{-59}$).
The theorem uses the cited theorems as inputs. The identification
$\RviiK_L=\RviiO$ is an instance of the known fact that straightedge, compass
and a fixed drawn non-circular conic construct exactly Alperin's field
\cite{R7-Videla,R7-BCJLS}, because inversion maps $L$ to the hyperbola
$X^2-Y^2=1$; the answer to Problem~\ref{R7-prob-drawn} is thus a combination of
known results, which we record because the question was left open in
\cite{R7-MRS16}.
\end{remark}

%=====================================================================
\subsection{A zero-free and Pick obstruction for HCM functions}
\label{R7-ssec-pick}

A function $f:(0,\infty)\to(0,\infty)$ is \emph{HCM} if for each $u>0$ the
function $H_u(w)=f(uv)f(u/v)$, $w=v+v^{-1}$, is completely monotone (CM) on
$(2,\infty)$, and \emph{HM} if each $H_u$ is nonincreasing there
\cite{R7-JS13,R7-Bondesson92}. A law on $(0,\infty)$ is GGC if it is a weak
limit of finite sums of independent gamma variables. Put
$\Omega=\RviiC\setminus(-\infty,0]$.

\begin{lemma}\label{R7-lem-pick}
Let $f>0$ on $(0,\infty)$ extend analytically to $\Omega$, and suppose $f$ is
HCM. Then
\begin{enumerate}
\item[(1)] $f$ has no zero in $\Omega$;
\item[(2)] $\psi_f(z):=-zf'(z)/f(z)$ satisfies $\RviiIm\psi_f(z)\ge0$ for
  $\RviiIm z>0$.
\end{enumerate}
\end{lemma}

\begin{proof}
Fix $u>0$. The map $v\mapsto w=v+v^{-1}$ sends $\Omega$ two-to-one onto
$D:=\RviiC\setminus(-\infty,-2]$ (it sends $(-\infty,0)$ onto
$(-\infty,-2]$). The function $G(v)=f(uv)f(u/v)$ is analytic on $\Omega$ and
invariant under $v\mapsto1/v$, so it descends to an analytic function
$H_u$ on $D$ (at $w=2$ write $v=e^t$: $G$ is even in $t$ and $w-2=t^2+O(t^4)$).

By Bernstein's theorem $H_u(w)=\int_{[0,\infty)}e^{-wt}\mu(dt)$ for $w>2$, with
$\mu\ge0$, $\mu\ne0$. Let $\sigma$ be the abscissa of convergence. By Landau's
theorem \cite[Chapter~II, Theorem~5b]{R7-Widder41} the real point $\sigma$,
if finite, is a singular point of the function defined by the integral. That
function agrees with $H_u$ on $\{\RviiRe w>\max(\sigma,-2)\}$, and $H_u$ is
analytic on $D\supset(-2,\infty)$; hence $\sigma\le-2$, and for every
$w\in(-2,\infty)$
\[
 H_u(w)=\int e^{-wt}\mu(dt)>0,\qquad H_u'(w)=-\int te^{-wt}\mu(dt)\le0 .
\]
For $\phi\in(0,\pi)$, $w=2\cos\phi\in(-2,2)$ corresponds to $v=e^{i\phi}$, and
since $f(\bar z)=\overline{f(z)}$,
\[
 H_u(2\cos\phi)=f(ue^{i\phi})f(ue^{-i\phi})=|f(ue^{i\phi})|^2 .
\]
Positivity gives (1) (every $z\in\Omega\setminus(0,\infty)$ is $ue^{\pm i\phi}$).
Differentiating in $\phi$, with $z=ue^{i\phi}$,
\[
 2|f(z)|^2\,\RviiIm\psi_f(z)=\frac{d}{d\phi}|f(ue^{i\phi})|^2
 =-2\sin\phi\,H_u'(2\cos\phi)\ \ge 0 ,
\]
because $\frac{d}{d\phi}\log|f(ue^{i\phi})|=\RviiRe\bigl(iz f'(z)/f(z)\bigr)=
\RviiIm\psi_f(z)$. This is (2).
\end{proof}

\begin{remark}
The lemma also follows from Bondesson's representation
$f(x)=Cx^{\beta-1}h_1(x)h_2(1/x)$ with $-(\log h_i)'$ Stieltjes
\cite[(5.2.3), Theorem~5.3.1]{R7-Bondesson92}, recalled as equation~(7) of
\cite{R7-BondessonSimon18}; the proof above uses only Bernstein's and Landau's
theorems. Part~(1) is the zero-free obstruction used for McKay laws in
\cite{K-029}; part~(2) is stronger and is what decides the boundary cases below.
\end{remark}

%=====================================================================
\subsection{\texorpdfstring{$\RviiHM\cap\RviiGGC\ne\RviiHCM$}{HM and GGC versus HCM}}
\label{R7-ssec-JS}

\begin{problem}[{Jedidi and Simon \cite[Remark~4.1(b)]{R7-JS13}}]
\label{R7-prob-JS}
``From the above equivalence, one can ask if the general identification
$\mathrm{HM}\cap\mathrm{GGC}=\mathrm{HCM}$ is true or not, and we could not
find any counterexample. Such an identification would show Bondesson's
conjecture in view of the main result of [22].'' ([22] refers to the source's
bibliography.)
\end{problem}

\begin{theorem}\label{R7-thm-JS}
Let $0<\lambda_1<\lambda_2$ and let $X=E_1/\lambda_1+E_2/\lambda_2$ with
$E_1,E_2$ independent standard exponential variables. Then $X$ is GGC, its
density is HM, and its density is not HCM. Hence
$\RviiHM\cap\RviiGGC\ne\RviiHCM$.
\end{theorem}

\begin{proof}
$X$ is a sum of two independent gamma variables, hence GGC (Thorin measure
$\delta_{\lambda_1}+\delta_{\lambda_2}$). Its density is
$f(x)=\frac{\lambda_1\lambda_2}{\delta}e^{-\lambda_1x}(1-e^{-\delta x})$, $\delta=\lambda_2-\lambda_1$.

HM: write $u=e^\sigma$, $v=e^t$ and $\Phi(y)=\log f(e^y)$; then
$\log H_u=\Phi(\sigma+t)+\Phi(\sigma-t)$, and $w=2\cosh t$ increases with
$t\ge0$, so $H_u$ is nonincreasing for every $u$ iff $\Phi'$ is
nonincreasing, i.e.\ iff $x\mapsto xf'(x)/f(x)$ is nonincreasing. Here
\[
 \frac{xf'(x)}{f(x)}=-\lambda_1x+\frac{\delta x}{e^{\delta x}-1},
\]
and both terms decrease, since $\frac{d}{dt}\frac{t}{e^t-1}=
-\frac{te^t-e^t+1}{(e^t-1)^2}<0$ for $t>0$.

Not HCM: $f$ is entire and vanishes at $z=2\pi i/\delta\in\Omega$,
contradicting Lemma~\ref{R7-lem-pick}(1). (Concretely, for
$\lambda_1=1,\lambda_2=2$ and $u=2\pi$ one gets $H_u(0)=|f(2\pi i)|^2=0$,
and the $16$th Taylor coefficient of $H_{2\pi}$ at $w=2.2$ has the wrong sign
for complete monotonicity; see \texttt{R7\_\allowbreak{}verify\_\allowbreak{}hcm.py}.)
\end{proof}

\begin{remark}
The example is elementary; we found no earlier statement that it answers the
question of \cite{R7-JS13}. The same zero argument shows that no sum of two
independent gamma variables with different rates is HCM (the density is
$x^{\beta_1+\beta_2-1}e^{-\lambda_2x}M(\beta_1,\beta_1+\beta_2,\delta x)$, and
Kummer's $M$ is the Laplace transform of a non-degenerate measure on $[0,1]$,
so it has infinitely many zeros, all non-real); such sums are always GGC, and
they are HM in parameter ranges that include the exponential case above.
\end{remark}

%=====================================================================
\subsection{Square roots of gamma products}
\label{R7-ssec-bosch}

Let $\RviiGam_t$ denote a gamma variable with density
$x^{t-1}e^{-x}/\Gamma(t)$; products are of independent copies.

\begin{problem}[{Bosch \cite[Remark~1]{R7-Bosch15}}]\label{R7-prob-bosch}
(a) ``It is hence natural to raise the [following] Conjecture. For all
$s,t>0$, the independent product $\sqrt{\gamma_t\gamma_s}$ is HCM.''
(b) ``[\dots] we believe that the latter should be true for all independent
products of the type'' $(\gamma_{t_1}\times\cdots\times\gamma_{t_n})^{1/n}$,
$t_i>0$.
\end{problem}

For $|t-s|<1/2$, (a) follows from the argument sketched by Bosch in
\cite[Remark~1(a)]{R7-Bosch15}; see also \cite[Section~3.3]{R7-BoschSimon15}.
For $|t-s|=1/2$ it follows from the duplication formula, which gives
$\sqrt{\RviiGam_t\RviiGam_{t+1/2}}\overset d=\RviiGam_{2t}/2$, a gamma law.

\begin{theorem}\label{R7-thm-bosch}
For $s,t>0$ the density of $\sqrt{\RviiGam_t\RviiGam_s}$ is HCM if and only if
$|t-s|\le1/2$. In particular Bosch's conjecture fails, e.g.\ for
$\sqrt{\RviiGam_1\RviiGam_2}$ and for $\sqrt{\RviiGam_3\RviiGam_{1/2}}$, whose
density is $\frac14(4x^2+6x+3)e^{-2x}$.
\end{theorem}

\begin{proof}
The density of $\RviiGam_t\RviiGam_s$ is
$\frac{2}{\Gamma(t)\Gamma(s)}y^{(t+s)/2-1}K_{t-s}(2\sqrt y)$, so that of
$Y=\sqrt{\RviiGam_t\RviiGam_s}$ is
\[
 f(x)=\frac{4}{\Gamma(t)\Gamma(s)}x^{t+s-1}K_\nu(2x),\qquad \nu=|t-s|
\]
($K_{-\nu}=K_\nu$). For $(t,s)=(3,\tfrac12)$, $K_{5/2}(\zeta)=\sqrt{\pi/(2\zeta)}\,
e^{-\zeta}(1+3\zeta^{-1}+3\zeta^{-2})$ gives the displayed density.
Sufficiency is the case discussed before the theorem; we prove necessity. Let $\nu>1/2$. The
function $f$ is analytic on $\Omega$ and positive on $(0,\infty)$. From
$K_\nu'(\zeta)=-K_{\nu-1}(\zeta)-\frac\nu\zeta K_\nu(\zeta)$
\cite[10.29.2]{R7-DLMF},
\[
 \psi_f(z)=\nu-(t+s-1)+R_\nu(2z),\qquad R_\nu(\zeta):=\frac{\zeta K_{\nu-1}(\zeta)}{K_\nu(\zeta)} .
\]
We use the convergent expansions
$K_\mu(\zeta)=\tfrac12\Gamma(\mu)(\zeta/2)^{-\mu}(1+O(|\zeta|^{\min(2,2\mu)}))$
for $\mu>0$, $\mu\notin\RviiZ$, the analogous statement with logarithmic
corrections for $\mu\in\RviiZ_{>0}$, and
$K_0(\zeta)=-\log(\zeta/2)-\gamma+O(|\zeta|^2|\log\zeta|)$
\cite[10.27.4, 10.31.1, 10.31.2]{R7-DLMF}, all uniform for
$0<\arg\zeta<\pi$ with principal branches. Hence, as $\zeta\to0$ in the upper
half-plane:
\[
 R_\nu(\zeta)=\begin{cases}
 \dfrac{2\Gamma(1-\nu)}{\Gamma(\nu)}\Bigl(\dfrac\zeta2\Bigr)^{2\nu}(1+o(1)), & \tfrac12<\nu<1,\\[2mm]
 \zeta^2\bigl(-\log(\zeta/2)-\gamma\bigr)(1+o(1)), & \nu=1,\\[1mm]
 \dfrac{\zeta^2}{2(\nu-1)}(1+o(1)), & \nu>1 .
 \end{cases}
\]
Take $\zeta=\varrho e^{i\theta}$ with $\varrho\to0$ and fixed $\theta$:
for $\frac12<\nu<1$ choose $\theta\in(\pi/(2\nu),\pi)$, so that
$\sin(2\nu\theta)<0$; for $\nu=1$ choose $\theta\in(3\pi/4,\pi)$, so that
$\sin2\theta<0<\cos2\theta$ and
$\RviiIm[\zeta^2(-\log(\zeta/2)-\gamma)]=\varrho^2[\sin2\theta\,(\log(2/\varrho)-\gamma)-\theta\cos2\theta]
\le-|\sin 2\theta|\varrho^2(\log(2/\varrho)-\gamma)$; for $\nu>1$ choose
$\theta\in(\pi/2,\pi)$. In each case $\RviiIm R_\nu(\zeta)<0$ for small
$\varrho$ (the relative errors are $o(1)$ while the main terms have imaginary
parts of the same order as their moduli). Then $\RviiIm\psi_f(\zeta/2)<0$ with
$\RviiIm(\zeta/2)>0$, contradicting Lemma~\ref{R7-lem-pick}(2).
\end{proof}

\begin{proposition}\label{R7-prop-nfold}
For every $n\ge2$ there are $t_1,\dots,t_n>0$ such that
$(\RviiGam_{t_1}\cdots\RviiGam_{t_n})^{1/n}$ is not HCM. Explicitly:
\begin{enumerate}
\item[(1)] for $n\in\{2,3,4\}$ and any $t>0$, with $t_1=t+3$ and
  $t_j=t+(j-1)/n$ ($2\le j\le n$), the density is
  $Cx^{nt}e^{-nx}Q_n(x)$ with
  $Q_n(x)=x^2+\frac{3(n-1)}{n}x+\frac{(n-1)(2n-1)}{n^2}$, which has non-real
  zeros;
\item[(2)] for every $n\ge2$, $t_1=1$ and $t_j=1+\frac{j+1}{2n}$ ($2\le j\le n$)
  violate Lemma~\ref{R7-lem-pick}(2).
\end{enumerate}
\end{proposition}

\begin{proof}
(1) Let $g(x)=x^{nt-1}e^{-nx}$ and let $\mathcal D h=-(xh)'$, so that
$\int_0^\infty x^s\mathcal Dh\,dx=s\int_0^\infty x^sh\,dx$. By Gauss's
multiplication formula the Mellin transform of
$(\RviiGam_{t+3}\prod_{j=1}^{n-1}\RviiGam_{t+j/n})^{1/n}$ is a constant times
$(t+\frac sn)(t+1+\frac sn)(t+2+\frac sn)\,\Gamma(nt+s)n^{-s}$, i.e.\ the law has
density proportional to $T_tT_{t+1}T_{t+2}\,g$ with $T_c=c+\mathcal D/n$.
Since $T_c[x^\alpha e^{-nx}]=(x+c-\frac{\alpha+1}n)x^\alpha e^{-nx}$, a direct
computation gives $T_tT_{t+1}T_{t+2}\,g=x^{nt}e^{-nx}Q_n(x)$. The discriminant
of $Q_n$ is $(n-1)(n-5)/n^2<0$ for $n\le4$, so $f$ has zeros in $\Omega$;
apply Lemma~\ref{R7-lem-pick}(1).

(2) All $t_j$ lie in $[1,2)$ and are distinct, so the poles of the Mellin
transform $M(s)=\prod_j\Gamma(t_j+s/n)/\Gamma(t_j)$ are simple and the density
of $Y$ is the convergent residue series (Meijer $G$-function expansion)
\[
 f_Y(x)=\frac{n}{\prod_i\Gamma(t_i)}\sum_{j=1}^n\sum_{k\ge0}\frac{(-1)^k}{k!}
 \prod_{i\ne j}\Gamma(t_i-t_j-k)\,x^{n(t_j+k)-1},
\]
valid on $\Omega$ with principal powers (the terms are
$O(|x|^{nk}/(k!)^n)$ in $k$, cf.\ \cite[16.17.2]{R7-DLMF}). Since
$nt_1-1=n-1$, $nt_2-1=n-1+\frac32$, $nt_j-1\ge n+1$ for $j\ge3$, and $n\ge2$,
\[
 f_Y(x)=\kappa\,x^{n-1}\Bigl(1+\frac{c_2}{c_1}x^{3/2}+O(|x|^2)\Bigr),\qquad
 \kappa>0,\quad c_j:=\prod_{i\ne j}\Gamma(t_i-t_j);
\]
here $c_1>0$ and $c_2<0$ (only the factor
$\Gamma(t_1-t_2)=\Gamma(-\frac3{2n})$ is negative). Differentiating termwise,
$\psi_{f_Y}(z)=-(n-1)-\frac32\frac{c_2}{c_1}z^{3/2}+O(|z|^2)$, so for
$z=\varrho e^{i\theta}$ with $\theta\in(2\pi/3,\pi)$,
$\RviiIm\psi_{f_Y}(z)=\frac32\bigl|\frac{c_2}{c_1}\bigr|\varrho^{3/2}\sin\frac{3\theta}2+O(\varrho^2)<0$
for small $\varrho$.
\end{proof}

The cases $t_j=j/n$ singled out by Bosch for stable laws are not affected:
there consecutive gaps are $1/n$ and the leading correction exponent is
$1$.

%=====================================================================
\subsection{Gamma--beta products}
\label{R7-ssec-gammabeta}

Let $\RviiB_{a,b}$ denote a beta variable with density
$x^{a-1}(1-x)^{b-1}/B(a,b)$ on $(0,1)$.

\begin{problem}[{Bosch and Simon \cite[after Lemma~1]{R7-BoschSimon16}}]
\label{R7-prob-BS}
``A perusal of the proof of Lemma 1 in [12] shows also that the product
$\Gamma_c\times\mathbf B_{a,b}$ has a HCM density as soon as $a+b\ge c$. We
believe that this latter condition is necessary.'' (In the source, [12] is
\cite{R7-Simon11}.)
\end{problem}

\begin{theorem}\label{R7-thm-gammabeta}
For $a,b,c>0$ the law of $\RviiGam_c\times\RviiB_{a,b}$ has an HCM density
if and only if $c\le a+b$.
\end{theorem}

\begin{proof}
\emph{Density.} With $y=1/(1+t)$ in
$f(x)=\int_0^1y^{-1}g_c(x/y)\beta_{a,b}(y)\,dy$ one finds
$f(x)=\frac{\Gamma(a+b)}{\Gamma(a)\Gamma(c)}x^{c-1}e^{-x}U(b,c-a+1,x)$, and
Kummer's transformation $U(\mathsf a,\mathsf b,x)=x^{1-\mathsf b}U(\mathsf a-\mathsf b+1,2-\mathsf b,x)$
\cite[13.2.40]{R7-DLMF} gives
\[
 f(x)=\frac{\Gamma(a+b)}{\Gamma(a)\Gamma(c)}\,x^{a-1}e^{-x}\,U(-\alpha,1-\delta,x),
 \qquad \alpha=c-a-b,\ \ \delta=c-a .
\]

\emph{Sufficiency} ($c\le a+b$). If $c\le a$ this is Lemma~1 of
\cite{R7-BoschSimon16} (Simon's lemma \cite{R7-Simon11}, extended to
$c\le\min a_i$ there by weak closure). If $a<c\le a+b$, write
$\RviiB_{a,b}\overset d=\RviiB_{a,c-a}\RviiB_{c,a+b-c}$ (independent factors;
$\RviiB_{c,0}:=1$) and use $\RviiGam_c\RviiB_{a,c-a}\overset d=\RviiGam_a$:
then $\RviiGam_c\RviiB_{a,b}\overset d=\RviiGam_a\RviiB_{c,a+b-c}$ with
$a<c$, which is HCM by the same lemma.

\emph{Necessity.} Let $c>a+b$, i.e.\ $\alpha>0$, and suppose $f$ were HCM.
Put $U(z)=U(-\alpha,1-\delta,z)$, analytic on $\Omega$; $U>0$ on $(0,\infty)$
because $f>0$. Then
\[
 \psi_f(z)=(1-a)+z-\frac{zU'(z)}{U(z)} .
\]
We use: (i) $U(z)=z^{\alpha}(1+O(|z|^{-1}))$ as $z\to\infty$ with
$0\le\arg z\le\pi$ \cite[13.2.6, 13.7.3]{R7-DLMF}; (ii) $U(z)\to\Gamma(\delta)/\Gamma(b)>0$ as $z\to0$
uniformly in $0\le\arg z\le\pi$ (the limiting forms
\cite[13.2.20--13.2.22]{R7-DLMF} for second parameter $1-\delta<1$, with
$\Gamma(1-(1-\delta))=\Gamma(\delta)$ and $-\alpha-(1-\delta)+1=b$);
(iii) $U$ continues analytically across $(-\infty,0)$ from above; write
$W(x)=U(-x+i0)$ for $x>0$.

By Lemma~\ref{R7-lem-pick}, $U$ has no zero in the open upper half-plane and
$\RviiIm\psi_f\ge0$ there. $W$ has no zero either: if $W(x_0)=0$ with order
$m\ge1$, then near $z_0=-x_0$ (from above)
$\psi_f(z)=-\frac{mz_0}{z-z_0}+O(1)$, and
$\RviiIm\psi_f(-x_0+i\eta)=-mx_0/\eta+O(1)\to-\infty$, a contradiction. Hence
$\psi_f$ is continuous up to $(-\infty,0)$ from above, and since
$\frac{d}{dx}\log W(x)=-U'/U$ at $z=-x+i0$,
\[
 0\le\RviiIm\psi_f(-x+i0)=\RviiIm\Bigl(-x\frac{W'(x)}{W(x)}\Bigr)=-x\,\frac{d}{dx}\arg W(x),
\]
so a continuous branch of $\arg W$ is nonincreasing on $(0,\infty)$.

Now apply the argument principle to $U$ on the half-annulus
$\{\varepsilon<|z|<R,\ 0<\arg z<\pi\}$, whose boundary carries no zero of $U$
for small $\varepsilon$ and large $R$, by (i), (ii) and the above. Along
$[\varepsilon,R]$ the argument of $U$ is constant; along $|z|=R$ it increases
by $\alpha\pi+O(1/R)$ by (i); along $|z|=\varepsilon$ it changes by $o(1)$ by
(ii); along the upper side of the cut, traversed from $-R$ to $-\varepsilon$,
it changes by $\arg W(\varepsilon)-\arg W(R)$. Since there are no zeros inside,
the total is $0$, whence
\[
 \arg W(R)-\arg W(\varepsilon)=\alpha\pi+o(1)\qquad(\varepsilon\to0,\ R\to\infty).
\]
This is positive for small $\varepsilon$ and large $R$, contradicting the
monotonicity of $\arg W$. Hence $f$ is not HCM.
\end{proof}

\begin{remark}
For $c=a+b+1$ the density is proportional to $x^{a-1}(x+b)e^{-x}$ and the
obstruction is visible directly:
$\psi_f(z)=(1-a)+z-1+\frac{b}{z+b}$ has
$\RviiIm\psi_f(-b+i\eta)=\eta-b/\eta<0$ for $0<\eta<\sqrt b$. The script
\texttt{R7\_\allowbreak{}verify\_\allowbreak{}hcm.py} checks the density formula by Mellin moments and,
for several $(a,b,c)$ with $c>a+b$, the identity
$\arg W(R)-\arg W(\varepsilon)\approx\pi(\alpha-2N_+)$, $N_+$ the number of zeros
of $U$ in the upper half-plane: both alternatives occur ($N_+=0$ with $\arg W$
increasing, e.g.\ $(a,b,c)=(0.7,0.4,1.6)$; $N_+=1$, e.g.\ $(2,0.5,4.3)$, where
$f$ itself vanishes in $\Omega$), and each excludes HCM. It also checks
$\RviiIm\psi_f\ge0$ on a grid when $c\le a+b$.
\end{remark}

%=====================================================================
\subsection{Verification, scope, and an unresolved target}
\label{R7-ssec-scope}

The scripts \texttt{R7\_\allowbreak{}verify\_\allowbreak{}tricuspoid.py},
\texttt{R7\_\allowbreak{}verify\_\allowbreak{}drawn\_\allowbreak{}lemniscate.py}, \texttt{R7\_\allowbreak{}verify\_\allowbreak{}hcm.py} and
\texttt{R7\_\allowbreak{}verify\_\allowbreak{}lorch\_\allowbreak{}muldoon.py} (run with \texttt{python3 -I}) print
PASS/FAIL for each check described above. They are independent of the proofs:
Galois-group orders, exact factorizations, quadrature of arc length, Mellin
moments, explicit Taylor coefficients of $H_u$, and argument tracking of $U$.

Theorems~\ref{R7-thm-drawn} and \ref{R7-thm-bosch} (sufficiency) and
\ref{R7-thm-gammabeta} (sufficiency) rest on the cited results of
Alperin, Cox--Shurman, Videla, Baek et al., Bosch, and Bosch--Simon; the
remaining statements are proved here, the minimality of $n=32$ and the
irreducibility and ramification claims of Corollary~\ref{R7-cor-tri} by exact
computer algebra.

\begin{remark}[Lorch--Muldoon, Section 2.3(b),(c): not settled]
For $f=x^{-\nu}J_\nu(x)$ and $c=2\nu+1$, Lorch and Muldoon
\cite[Section~2.3]{R7-LM95} conjectured that the resultant $r_n(\nu)$ of the
two coefficient polynomials of $f^{(n)}$ has exactly the rational roots
$-m/2$ ($1\le m\le n-1$), together with $0$ for $n=4$ and $1$ for $n=5$. Since
$r_n$ is monic in $c$ (\cite{K-047}), rational roots are half-integers.
\texttt{R7\_\allowbreak{}verify\_\allowbreak{}lorch\_\allowbreak{}muldoon.py} verifies (b),(c) for all $n\le30$
(Lorch--Muldoon: $n\le13$). Computations show that for $6\le n\le30$ the non-linear part of $r_n$ is
$Q_n(c)^2$ with $Q_n$ irreducible of degree $\binom{\lfloor n/2\rfloor}{2}$,
and (for $|c|\le41$, $n\le26$) that $Q_n(c)$ is odd for even $c$; we have no
proof for all $n$.
\end{remark}


%% ==================================================================
%% Section: Mixtures
%% ==================================================================
% ======================================================================
% R5: estimation of mixtures and populations of parameters
% LaTeX fragment (no preamble). Requires amsmath, amssymb, amsthm and the
% environments theorem, lemma, proposition, corollary, definition, remark,
% conjecture, problem.  All labels are prefixed by R5.
% ======================================================================
\providecommand{\RfE}{\mathbb{E}}
\providecommand{\RfP}{\mathbb{P}}
\providecommand{\RfR}{\mathbb{R}}
\providecommand{\RfZ}{\mathbb{Z}}
\providecommand{\Rfd}{\,\mathrm{d}}
\providecommand{\Rffold}{\operatorname{fold}}
\providecommand{\RfBin}{\operatorname{Bin}}
\providecommand{\RfLip}{\operatorname{Lip}}

\section{Mixtures and populations of parameters: \texorpdfstring{$W_p$}{Wp} recovery, self-regularization, and a moment-matching constant}
\label{R5:sec}

This section treats four questions about mixture estimation that were posed
by other authors and that we could not find settled, either in the
literature or in the earlier parts of this programme.  Their status is as
follows.

\begin{enumerate}
\item[(A)] Vinayak, Kong, Valiant and Kakade asked for Wasserstein-$p$
 guarantees for the maximum likelihood estimator of a population of coin
 biases~\cite[Sec.~5]{R5-VKVK19}.  Theorem~\ref{R5:thm:wp} gives the full
 minimax rate for every fixed $p\in[1,\infty)$, for every number of tosses
 $t$ and population size $N$, and shows that every nonparametric maximum
 likelihood estimator (NPMLE) attains it.  Relative to $p=1$
 \cite[Part~VI]{K-029}, only the parametric term changes, from $N^{-1/2}$
 to $N^{-1/(2p)}$.  \emph{Closed, conditionally} on the $W_1$ theorem and three
 auxiliary estimates of \cite[Part~VI]{K-029}, recalled as (F1)--(F4) in
 Lemma~\ref{R5:lem:recall}; that manuscript is not yet refereed.
\item[(B)] Polyanskiy and Wu conjectured that a Gompertz location mixture
 with mixing law on $[-a,a]$ has $O(a)$ critical points
 \cite[Sec.~5.2]{R5-PW20}.  Theorem~\ref{R5:thm:gompertz} shows that it can
 have at least $a(a-4\log 2)/(2\pi^2)-2$ local maxima.  \emph{Disproved.}
\item[(C)] Polyanskiy and Wu asked whether their $\Omega(k^2)$ lower bound
 for the number of modes of a mixture of $k$ analytic strongly log-concave
 densities is tight \cite[Sec.~5.5]{R5-PW20}.  Theorem~\ref{R5:thm:modes}
 shows that it is not.  Already for $k=2$, and with all second derivatives
 of the log-densities confined to $[-3/2,-1/2]$, the number of modes is
 unbounded.  \emph{Answered (negatively).}
\item[(D)] Musco, Musco, Rosenblatt and Singh asked whether the constant in
 their Chebyshev moment-matching bound $W_1\le 36/k+\Gamma$ can be improved
 ``at least to $2\pi/k+\Gamma$, and possibly further''
 \cite[Sec.~1.2 and~3.1]{R5-MMRS25}.  Theorem~\ref{R5:thm:favard} shows that
 $\pi/(k+1)+\Gamma$ holds and that the constant $\pi$ is asymptotically
 optimal.  \emph{Closed, but by classical approximation theory}: the upper
 bound is a direct application of Favard's 1937 theorem.  We include it for
 completeness and do not count it as a new mathematical result.
\end{enumerate}

Items (B) and (C) are short constructive disproofs.  Item (A) is the main
result.  It reuses the likelihood-transport, Hellinger and
Chebyshev--Bernstein estimates of \cite[Part~VI]{K-029}.  The new
observation is that smoothing the fitted and true mixing laws at polynomial
degree at most $t/2$ makes the effect of the unobservable components
negligible (at most $2^{-m}$ after polynomial approximation).  The $W_1$
distance between the smoothed laws is then purely statistical, and
$W_p^p\le W_1$ on $[0,1]$ converts its parametric size into $N^{-1/(2p)}$
(Remark~\ref{R5:rem:wp}).  Subsection~\ref{R5:sec:settled} lists candidate
questions that turned out to be already settled by others.

% ----------------------------------------------------------------------
\subsection{Wasserstein-\texorpdfstring{$p$}{p} minimax rates for the binomial NPMLE}
\label{R5:sec:wp}

\begin{problem}[Vinayak, Kong, Valiant and Kakade {\cite[Sec.~5]{R5-VKVK19}}]
\label{R5:prob:wp}
``Another direction of interest is to extend the analysis to provide
guarantees in general Wasserstein-$p$ norms.''
\end{problem}

Here ``the analysis'' is the minimax analysis, in Wasserstein-$1$ distance,
of the nonparametric MLE of the population of coin biases.  That analysis
was completed for every $(N,t)$ in \cite[Part~VI, Theorem
p06:bm:full]{K-029}.  We answer the question for every finite $p$.

\paragraph{Setting.}
Let $G$ be a probability measure on $[0,1]$.  We observe $N$ independent
counts $X_1,\dots,X_N$, where $X_i\mid\theta_i\sim\RfBin(t,\theta_i)$ and
$\theta_i\sim G$ are independent.  Put $B^t_j(x)=\binom tj x^j(1-x)^{t-j}$
and $p_H(j)=\int B^t_j\Rfd H$ for $0\le j\le t$.  Let $q$ be the empirical
law of the counts.  An NPMLE is any maximizer $\widehat G$ of
$H\mapsto\sum_j q_j\log p_H(j)$ over all probability measures on $[0,1]$.
Write $\widehat p=p_{\widehat G}$, $L_N=\log(eN)$,
\[
 r_{N,t}=\min\Bigl\{1,\frac1t+\frac1{\sqrt{tL_N}}+\frac1{\sqrt N}\Bigr\},
 \qquad
 \rho_p(N,t)=\min\Bigl\{1,\frac1t+\frac1{\sqrt{tL_N}}+N^{-1/(2p)}\Bigr\}.
\]
On the count grid we use the normalized transport distance $W_{1,t}$, with
cost $|j-l|/t$, and the Hellinger distance
$h(u,v)^2=\sum_j(\sqrt{u_j}-\sqrt{v_j})^2$.

\begin{theorem}[$W_p$-minimax optimality of the binomial NPMLE]
\label{R5:thm:wp}
Assume the facts (F1)--(F4) of Lemma~\ref{R5:lem:recall}, taken from
\cite[Part~VI]{K-029}. For every $p\in[1,\infty)$ there are constants $0<c_p\le C_p<\infty$ such
that, for all integers $N,t\ge1$,
\begin{equation}\label{R5:eq:wp}
 c_p\,\rho_p(N,t)\;\le\;\inf_{\widetilde G}\sup_G\RfE_G W_p(\widetilde G,G)
 \;\le\;\sup_G\RfE_G W_p(\widehat G,G)\;\le\;C_p\,\rho_p(N,t).
\end{equation}
The infimum is over all (possibly randomized) estimators.  The upper bound
holds for every measurable choice of an NPMLE.
\end{theorem}

For $p=1$ the theorem is \cite[Part~VI]{K-029}.  For $p>1$ the
deconvolution terms $1/t+(tL_N)^{-1/2}$ are unchanged, and the parametric
term becomes $N^{-1/(2p)}$.  This term dominates once
$t\gtrsim N^{1/p}/L_N$.  As $p\to\infty$ it tends to $1$, which is
consistent with $W_\infty$: there the minimax risk is bounded below by a
positive constant for all $N$ and $t$.  Indeed $\delta_0$ and
$(1-\varepsilon)\delta_0+\varepsilon\delta_1$ are at $W_\infty$-distance one,
and they cannot be told apart when $\varepsilon\ll1/N$.

\subsubsection{Facts recalled from Part VI of \cite{K-029}}

\begin{lemma}[recalled]\label{R5:lem:recall}
There is an absolute constant $C_0\ge1$ with the following properties.
\begin{enumerate}
\item[(F1)] $\sup_G\RfE_GW_1(\widehat G,G)\le C_0r_{N,t}$ for all $N,t\ge1$
 and every NPMLE.  Also
 \[\inf_{\widetilde G}\sup_G\RfE_GW_1(\widetilde G,G)\ge C_0^{-1}r_{N,t}.\]
 \emph{(\cite[Part VI, Theorem p06:bm:full]{K-029}.)}
\item[(F2)] If $P=\sum_{j=0}^tb_jB_j^t$ is a polynomial of degree $k\le t/2$,
 then $\max_j|b_j|\le C_0e^{2k^2/t}\|P\|_{\infty}$.  If moreover $k\ge1$, then
 $\max_{j<t}t|b_{j+1}-b_j|\le C_0e^{4k^2/t}\|P'\|_{\infty}$.  Here and below
 $\|\cdot\|_\infty$ is the supremum on $[0,1]$.
 \emph{(\cite[Part VI, Lemma p06:bm:approx]{K-029}; the second bound is the
 first one applied to $P'=t\sum_{j<t}(b_{j+1}-b_j)B^{t-1}_j$.)}
\item[(F3)] $\RfE_GW_{1,t}(\widehat p,p_G)\le C_0N^{-1/2}$ for all $N,t$ and $G$.
 \emph{(\cite[Part VI, Lemma p06:bm:transport]{K-029}.)}
\item[(F4)] If $1\le t\le N$, then
 $\RfE_Gh(p_G,\widehat p)\le C_0t^{1/4}L_N^{3/4}N^{-1/2}$.
 \emph{(\cite[Part VI, Corollary p06:bm:hellinger]{K-029}.)}
\end{enumerate}
\end{lemma}

We also use two elementary dualities on the count grid.  Let $u,v$ be
probability vectors on $\{0,\dots,t\}$ and let $b\in\RfR^{t+1}$.  Then
\begin{equation}\label{R5:eq:griddual}
 \Bigl|\sum_jb_j(u_j-v_j)\Bigr|\le\Bigl(\max_{j<t}t|b_{j+1}-b_j|\Bigr)\,W_{1,t}(u,v),
 \qquad
 \Bigl|\sum_jb_j(u_j-v_j)\Bigr|\le 2\max_j|b_j|\,h(u,v).
\end{equation}
The first inequality is Kantorovich--Rubinstein duality.  The second follows
from $\|u-v\|_1\le h(u,v)\bigl(\sum_j(\sqrt{u_j}+\sqrt{v_j})^2\bigr)^{1/2}\le2h(u,v)$.

\subsubsection{Reflected Jackson smoothings}

Let $\Rffold(z)=\operatorname{dist}(z,2\RfZ)$.  This map sends $\RfR$ onto
$[0,1]$, it is $1$-Lipschitz, and it is the identity on $[0,1]$.  For a real
random variable $U$ and a probability measure $H$ on $[0,1]$, let $S_UH$ be
the law of $\Rffold(Y+U)$, where $Y\sim H$ is independent of $U$.  Then
$S_U$ is a Markov operator, and $|\Rffold(Y+U)-Y|\le|U|$.  Hence
\begin{equation}\label{R5:eq:smoothcost}
 W_q(H,S_UH)\le(\RfE|U|^q)^{1/q}\qquad(q\ge1).
\end{equation}

For integers $n\ge1$ and $\rho\ge2$ let
\[
 J_{n,\rho}(u)=c_{n,\rho}\Bigl(\frac{\sin(n\pi u/2)}{\sin(\pi u/2)}\Bigr)^{2\rho},
 \qquad u\in[-1,1],
\]
normalized to be a probability density on $[-1,1]$, and let $U_{n,\rho}$
have this density.  Write $S_{n,\rho}=S_{U_{n,\rho}}$ and
$\Lambda=\rho(n-1)$.

\begin{lemma}\label{R5:lem:kernel}
Let $U=U_{n,\rho}$ and $S=S_{n,\rho}$.
\begin{enumerate}
\item[(i)] For $0<v\le1$, $\RfP(|U|>v)\le(\pi/2)^{2\rho}(nv)^{1-2\rho}$.
 Consequently, for $1\le q<2\rho-1$,
 \[
 \RfE|U|^q\le\Bigl(1+\frac{q(\pi/2)^{2\rho}}{2\rho-1-q}\Bigr)n^{-q}.
 \]
\item[(ii)] For bounded measurable $f$ on $[0,1]$, the adjoint
 $S^*f(y)=\RfE f(\Rffold(y+U))$ is the restriction to $[0,1]$ of a
 trigonometric polynomial $\sum_{|l|\le\Lambda}c_le^{i\pi ly}$ with
 $|c_l|\le\sup|f|$.  If $f$ is $1$-Lipschitz, then so is $S^*f$, and
 $\|S^*f-f\|_\infty\le\RfE|U|$.
\item[(iii)] Let $g=\sum_{|l|\le\Lambda}c_le^{i\pi ly}$ and $D\ge9\Lambda+1$.
 Let $\mathcal T_D[g]$ be the polynomial obtained by replacing each
 $e^{i\pi ly}$ by its Taylor polynomial of degree $D$ at $y=1/2$.  Then on
 $[0,1]$,
 \[
 \|g-\mathcal T_D[g]\|_\infty\le(2\Lambda+1)\max_l|c_l|\,2^{-D-1},\qquad
 \|g'-\mathcal T_D[g]'\|_\infty\le(2\Lambda+1)\pi\Lambda\max_l|c_l|\,2^{-D}.
 \]
 If $g$ is real-valued, so is $\mathcal T_D[g]$.
\end{enumerate}
\end{lemma}

\begin{proof}
(i) For $|u|\le1$ we have $|\sin(\pi u/2)|\ge|u|$, so
$J_{n,\rho}(u)\le c_{n,\rho}|u|^{-2\rho}$.  For $|u|\le1/n$, concavity of
$\sin$ on $[0,\pi/2]$ gives $|\sin(n\pi u/2)|\ge n|u|$.  Also
$|\sin(\pi u/2)|\le\pi|u|/2$, so the ratio is at least $2n/\pi$.  Hence
$c_{n,\rho}^{-1}\ge(2/n)(2n/\pi)^{2\rho}$.  Therefore
$\RfP(|U|>v)\le2c_{n,\rho}\int_v^\infty u^{-2\rho}\Rfd u\le(\pi/2)^{2\rho}(nv)^{1-2\rho}/(2\rho-1)$.
Then
\[
 \RfE|U|^q=\int_0^1qv^{q-1}\RfP(|U|>v)\Rfd v
 \le n^{-q}+q(\pi/2)^{2\rho}n^{1-2\rho}\int_{1/n}^\infty v^{q-2\rho}\Rfd v ,
\]
which is the stated bound.

(ii) Fej\'er's identity
$(\sin(n\pi u/2)/\sin(\pi u/2))^2=\sum_{|l|<n}(n-|l|)e^{i\pi lu}$ shows that
$J_{n,\rho}$ is a nonnegative even trigonometric polynomial of degree
$\Lambda$ in $\pi u$.  Its coefficients
$\hat\jmath_l=\frac12\int_{-1}^1J_{n,\rho}(u)e^{-i\pi lu}\Rfd u$ satisfy
$|\hat\jmath_l|\le\frac12$.  With $\tilde f=f\circ\Rffold$, which is even
and $2$-periodic, we get
$S^*f(y)=\int_{-1}^1\tilde f(y+u)J_{n,\rho}(u)\Rfd u
=\sum_l\hat\jmath_le^{i\pi ly}\int_{-1}^1\tilde f(s)e^{-i\pi ls}\Rfd s$,
using that $J_{n,\rho}$ is even.  This proves the first claim.  The map
$y\mapsto f(\Rffold(y+u))$ is $1$-Lipschitz for each $u$, and so is its
average.  Finally $|S^*f(y)-f(y)|\le\RfE|\Rffold(y+U)-y|\le\RfE|U|$.

(iii) For real $z$, $|e^{iz}-\sum_{d\le D}(iz)^d/d!|\le|z|^{D+1}/(D+1)!$.
Here $|z|=\pi|l||y-\frac12|\le\pi\Lambda/2$.  Also
$(D+1)!\ge((D+1)/e)^{D+1}$ and $D+1\ge e\pi\Lambda$, so the remainder is at
most $2^{-D-1}$.  The derivative of the remainder is $i\pi l$ times the
remainder of order $D-1$, which gives the second bound.  If $g$ is real,
then $c_{-l}=\overline{c_l}$, the terms of indices $\pm l$ are complex
conjugates, and so $\mathcal T_D[g]$ is real.
\end{proof}

\subsubsection{Smoothed \texorpdfstring{$W_1$}{W1} errors of the NPMLE}

The next proposition is the two-scale inversion of
\cite[Part~VI, Proposition p06:bm:twoscale]{K-029}, applied to smoothed
measures.  For these, the test functions are polynomials up to an
exponentially small error, so the Jackson term $C/k$ of that proposition
disappears.  This is the only place where the binomial data enter.

\begin{proposition}\label{R5:prop:smooth}
Let $1\le t\le N$ and let $\rho\ge2$ and $n\ge1$, $n_s\ge1$, $m\ge1$ be
integers.  Put $\Lambda=\rho(n-1)$, $\Lambda_s=\rho(n_s-1)$,
$D=\lceil9\Lambda\rceil+m+1$ and $D_s=\lceil9\Lambda_s\rceil+m+1$, and assume
$\max(D,D_s)\le t/2$.  Then for every $G$ and every NPMLE,
\[
 W_1(S_{n,\rho}\widehat G,S_{n,\rho}G)\le
 2C_0e^{4D_s^2/t}\,W_{1,t}(\widehat p,p_G)
 +2C_0e^{2\max(D,D_s)^2/t}\bigl(\RfE|U_{n_s,\rho}|+2^{-m}\bigr)h(\widehat p,p_G)
 +2^{-m}.
\]
\end{proposition}

\begin{proof}
Write $S=S_{n,\rho}$ and $S_s=S_{n_s,\rho}$.  By Kantorovich--Rubinstein
duality, $W_1(S\widehat G,SG)=\sup_f\int S^*f\Rfd(\widehat G-G)$ over
$1$-Lipschitz $f$ on $[0,1]$.  Adding a constant to $f$ does not change the
integral, so we may take $f(0)=0$; then $|f|\le1$.

Put $g=S^*f$ and $g_s=S_s^*g$.  By Lemma~\ref{R5:lem:kernel}(ii), $g$ is
$1$-Lipschitz, and it is a real trigonometric polynomial of degree $\Lambda$
with coefficients of modulus at most $1$.  Similarly $g_s$ is $1$-Lipschitz,
and it is a real trigonometric polynomial of degree at most $\Lambda_s$ with
coefficients of modulus at most $1$.  Moreover
$\|g-g_s\|_\infty\le\RfE|U_{n_s,\rho}|$.  Let $P=\mathcal T_D[g]$ and
$Q=\mathcal T_{D_s}[g_s]$.  These are real polynomials of degree at most $D$
and $D_s$.  Since $9\Lambda+m+1\le D$, Lemma~\ref{R5:lem:kernel}(iii) gives
$\|g-P\|_\infty\le(2\Lambda+1)2^{-9\Lambda-m-2}\le2^{-m-2}$.  In the same way
$\|g_s-Q\|_\infty\le2^{-m-2}$ and $\|Q'\|_\infty\le1+2^{-m}\le2$.

Every polynomial of degree at most $t$ is a combination of the $B^t_j$, and
$\int B^t_j\Rfd(\widehat G-G)=\widehat p_j-p_G(j)$.  We write
\[
 \int g\Rfd(\widehat G-G)=\int(g-P)\Rfd(\widehat G-G)+\int Q\Rfd(\widehat G-G)
 +\int(P-Q)\Rfd(\widehat G-G).
\]
The first term is at most $2\|g-P\|_\infty\le2^{-m-1}$.  By (F2) and the
first inequality of \eqref{R5:eq:griddual}, the second term is at most
$C_0e^{4D_s^2/t}\cdot2\cdot W_{1,t}(\widehat p,p_G)$.  For the third term,
$\|P-Q\|_\infty\le\|g-g_s\|_\infty+2^{-m-1}$.  By (F2) and the second
inequality of \eqref{R5:eq:griddual}, the third term is therefore at most
$2C_0e^{2\max(D,D_s)^2/t}(\RfE|U_{n_s,\rho}|+2^{-m})h(\widehat p,p_G)$.
These bounds hold uniformly in $f$, and also for $-f$.
\end{proof}

\subsubsection{Proof of Theorem~\ref{R5:thm:wp}}

\begin{proof}[Proof of the lower bound]
Since $W_p\ge W_1$, (F1) gives
$\inf\sup\RfE W_p\ge C_0^{-1}r_{N,t}$.  For the parametric term, let
$\delta=1/(4\sqrt N)$ and
$G_\pm=(\tfrac12\pm\delta)\delta_0+(\tfrac12\mp\delta)\delta_1$.  Under
$G_\pm$ every count equals $0$ or $t$.  So the sample is a vector of $N$
independent $\{0,t\}$-valued labels with probabilities $\frac12\pm\delta$.
The Bhattacharyya coefficient of one label pair is at least $1-4\delta^2$.
For the product laws it is therefore at least $(1-4\delta^2)^N\ge1-4N\delta^2=\frac34$,
and their total variation is at most $\sqrt{1-(3/4)^2}<0.7$.  Moreover
$W_p(G_+,G_-)=(2\delta)^{1/p}$.  Le Cam's two-point bound gives
$\inf\sup\RfE W_p\ge\frac12(1-0.7)(2\delta)^{1/p}\ge c\,N^{-1/(2p)}$.
Finally $\rho_p\le r_{N,t}+N^{-1/(2p)}$, and the maximum of two numbers is
at least half their sum.
\end{proof}

\begin{proof}[Proof of the upper bound]
Fix $p>1$; the case $p=1$ is (F1).  Constants $C$ may depend on $p$ and
change from line to line.  We use two elementary facts.  First,
\eqref{R5:eq:smoothcost} holds.  Second, $W_p^p(\mu,\nu)\le W_1(\mu,\nu)$
for probability measures on $[0,1]$: use a $W_1$-optimal coupling and
$|x-y|^p\le|x-y|$.  Since $W_p\le1$ and $\rho_p$ is bounded below by a
positive constant when $t$ is bounded, we may assume $t\ge T_0(p)$ for a
suitable $T_0(p)$.

\emph{Case $t>N$.}  For a law $H$ on $[0,1]$, let $H^\natural$ be the law of
$X/t$, where $X\mid\theta\sim\RfBin(t,\theta)$ and $\theta\sim H$.  Coupling
$\theta$ with $X/t$, and using $\RfE|X/t-\theta|^{q}\le C_qt^{-q/2}$ for an
even $q\ge p$, gives $W_p(H,H^\natural)\le Ct^{-1/2}$.  Moreover
$W_p^p(\widehat G^\natural,G^\natural)\le W_1(\widehat G^\natural,G^\natural)=W_{1,t}(\widehat p,p_G)$.
By Jensen's inequality and (F3),
\[
 \RfE W_p(\widehat G,G)\le2Ct^{-1/2}+(\RfE W_{1,t}(\widehat p,p_G))^{1/p}
 \le2CN^{-1/2}+C_0^{1/p}N^{-1/(2p)}\le C\rho_p(N,t).
\]

\emph{Case $T_0(p)\le t\le N$.}  Fix
\[
 \rho=\lceil p\rceil+1,\qquad \gamma=\frac1{4(2p+1)},\qquad
 c_1=\min\Bigl\{\frac1{36\rho},\frac{\sqrt\gamma}{18\rho}\Bigr\},
\]
and set $k_0=\lfloor c_1\min\{t,\sqrt{tL_N}\}\rfloor$.  For $t\ge T_0(p)$
we have $k_0\ge c_1\sqrt t/2\ge1$.  Let $m=\lceil p\log_2(k_0+1)\rceil+2$
and $n_s=\lfloor\sqrt t/(40\rho)\rfloor$.  Let $S=S_{k_0,\rho}$, so that
$\Lambda=\rho(k_0-1)$ and $D=\lceil9\Lambda\rceil+m+1\le9\rho k_0+m+2$.  By
the triangle inequality, \eqref{R5:eq:smoothcost}, $W_p^p\le W_1$ for the
smoothed laws, and Jensen's inequality,
\begin{equation}\label{R5:eq:wpchain}
 \RfE W_p(\widehat G,G)\le2\bigl(\RfE|U_{k_0,\rho}|^p\bigr)^{1/p}
 +\bigl(\RfE W_1(S\widehat G,SG)\bigr)^{1/p}.
\end{equation}

\emph{The smoothing cost.}  Since $2\rho-1\ge2p+1>p$,
Lemma~\ref{R5:lem:kernel}(i) gives
$(\RfE|U_{k_0,\rho}|^p)^{1/p}\le C/k_0$.

\emph{Parameters.}  We have $D_s\le9\rho n_s+m+2\le\sqrt t/4$ for
$t\ge T_0(p)$, so $e^{4D_s^2/t}\le e^{1/4}$.  Also $D\le t/2$, and
\[
 \frac{2D^2}{t}\le\frac{4(9\rho c_1)^2\min\{t,\sqrt{tL_N}\}^2}{t}
 +\frac{4(m+2)^2}{t}\le\gamma L_N+1 .
\]
So $e^{2\max(D,D_s)^2/t}\le e\,(eN)^{\gamma}$.  By
Lemma~\ref{R5:lem:kernel}(i), $\RfE|U_{n_s,\rho}|\le C/\sqrt t$.  Moreover
$2^{-m}\le(k_0+1)^{-p}/4\le k_0^{-1}\le C/\sqrt t$, because
$k_0\ge c_1\sqrt t/2$.

\emph{The smoothed $W_1$ error.}  Proposition~\ref{R5:prop:smooth}, (F3)
and (F4) give
\begin{align*}
 \RfE W_1(S\widehat G,SG)&\le C\Bigl[N^{-1/2}+(eN)^{\gamma}\frac{1}{\sqrt t}\cdot
 \frac{t^{1/4}L_N^{3/4}}{\sqrt N}+k_0^{-p}\Bigr]\\
 &=C\bigl[N^{-1/2}+(eN)^{\gamma}L_N^{3/4}t^{-1/4}N^{-1/2}+k_0^{-p}\bigr].
\end{align*}
If $t\ge(eN)^{4\gamma}L_N^3$, the middle term is at most $N^{-1/2}$.
Otherwise $k_0\le c_1\sqrt{tL_N}\le c_1(eN)^{2\gamma}L_N^2$.  Since
$\gamma(1+2p)=1/4$, the middle term divided by $k_0^{-p}$ is then at most
\[
 e^{1/2}c_1^p\,(eN)^{\gamma(1+2p)-1/2}L_N^{2p+3/4}=e^{1/2}c_1^p(eN)^{-1/4}L_N^{2p+3/4}\le C .
\]
Hence $\RfE W_1(S\widehat G,SG)\le C(k_0^{-p}+N^{-1/2})$.

Inserting both bounds into \eqref{R5:eq:wpchain} gives
\[
 \RfE W_p(\widehat G,G)\le C\Bigl(\frac1{k_0}+N^{-1/(2p)}\Bigr),
 \qquad
 \frac1{k_0}\le\frac2{c_1}\max\Bigl\{\frac1t,\frac1{\sqrt{tL_N}}\Bigr\}.
\]
Together with $W_p\le1$, this proves the upper bound in \eqref{R5:eq:wp}.
\end{proof}

\begin{remark}\label{R5:rem:wp}
(1) \emph{Mechanism.}  The moment-limited error of a mixture estimate
consists of components of polynomial degree above $t$.  These are
invisible to the data, and after smoothing at degree at most $t/2$ their
effect is at most $2^{-m}$.  So
$W_1(S\widehat G,SG)$ is purely statistical, and the inequality
$W_p^p\le W_1$ on the unit interval turns its parametric size $N^{-1/2}$
into $N^{-1/(2p)}$.  The smoothing itself costs only
$(\RfE|U|^p)^{1/p}\asymp1/k_0$ in $W_p$.  Applying $W_p^p\le W_1$ directly
to $(\widehat G,G)$ would only give the weaker rate
$(1/k_0+N^{-1/2})^{1/p}$.

(2) The proof uses only (F1)--(F4) and deterministic facts, so it applies
to every estimator for which analogues of (F1), (F3) and (F4) hold.  This
includes approximate NPMLEs within $N^{-1}$ of the optimal log-likelihood
\cite[Part~VI]{K-029}.

(3) A Monte Carlo check (script \texttt{R5\_\allowbreak{}wp\_\allowbreak{}npmle\_\allowbreak{}sim.py}, with the
NPMLE computed by EM on a grid of $2001$ points) is consistent with
\eqref{R5:eq:wp}.  For $G=\frac12(\delta_{0.2}+\delta_{0.8})$, $t=4000$ and
$N=100,\dots,6400$, the fitted log--log slopes of $\RfE W_p$ in $N$ are
$-0.54$, $-0.28$ and $-0.18$ for $p=1,2,3$.  The predictions are $-1/2$,
$-1/4$ and $-1/6$, and the dependence of the slope on $p$ is the new
parametric term $N^{-1/(2p)}$.  For uniform $G$, $N=20000$ and
$t=10,40,160,640$, the ratios $\RfE W_p/\RfE W_1$ lie between $1.2$ and
$1.45$ for $p=2,3$.  So for this smooth $G$ the dependence on $t$ is that
of the case $p=1$, and the worst-case term $N^{-1/(2p)}$ is not visible.
The products $\RfE W_p\cdot\sqrt{tL_N}$ lie between $0.18$ and $0.50$.
They rise at $t=640$, where $N^{-1/2}$ is already about
$0.6\,(tL_N)^{-1/2}$.  These experiments use an EM approximation of the
NPMLE (KKT gaps up to $1.8\cdot10^{-2}$) and only illustrate the theorem.

(4) As in \cite{R5-VKVK19} and \cite{K-029}, the estimation target is the
mixing law $G$.  If the target is instead the empirical law of
$\theta_1,\dots,\theta_N$, the upper bound is unchanged, by the triangle
inequality and $\RfE W_p(G,N^{-1}\sum_i\delta_{\theta_i})\le CN^{-1/(2p)}$.
For that target, however, the two-point lower bound of order $N^{-1/(2p)}$
need not hold.
\end{remark}

% ----------------------------------------------------------------------
\subsection{The Gompertz critical-point conjecture of Polyanskiy and Wu is false}
\label{R5:sec:gompertz}

Polyanskiy and Wu \cite{R5-PW20} prove that the NPMLE of a Gaussian
location mixture self-regularizes: for subgaussian data it has $O(\log n)$
atoms.  For a mixture of exponential distributions they pose the open
question of whether the NPMLE has $\Theta(\log n)$ atoms with high
probability when the mixing law is supported in $[1,2]$
\cite[Sec.~5.2, eq.~(37)]{R5-PW20}.  They reduce the upper bound to an
analytic conjecture.

\begin{conjecture}[Polyanskiy and Wu {\cite[Sec.~5.2]{R5-PW20}}]
\label{R5:conj:gompertz}
``The crucial $O(\log n)$ upper bound would follow from the following
analytic \emph{conjecture}: For any distribution $\pi$ on $[-a,a]$ the
convolution $(\pi*h)(x)\triangleq\int h(x-y)\pi(dy)$ has at most $O(a)$
critical points, where $h(x)=e^{-e^x+x}$ is the density of a Gompertz
distribution.''
\end{conjecture}

\begin{theorem}\label{R5:thm:gompertz}
Let $h(x)=e^{x-e^x}$.  For every $a\ge4$ let
$\omega_a=(a-4\log2)/\pi$ and let $\pi_a$ be the probability measure on
$[-a,a]$ with density proportional to $1+\sin(\omega_ay)$.  Then $\pi_a*h$
has at least
\[
 \Bigl\lfloor\frac{a(a-4\log2)}{2\pi^2}\Bigr\rfloor-1
\]
strict local maxima in $(-a/2,a/2)$, and each of them is a critical point.
In particular Conjecture~\ref{R5:conj:gompertz} is false: no bound of the
form $Ca$ can hold for all large $a$.
\end{theorem}

For the Gaussian kernel the analogous question had a similar history.  A
conjectured $O(a)$ bound was refuted by Kashyap and Krishnapur \cite{R5-KK20},
and \cite[Sec.~5.4]{R5-PW20} gives a short proof of an $\Omega(a^2)$ lower
bound.  The disproof is a direct corollary of the lemma Polyanskiy and Wu use to
show that the $O(a^2)$ bound for Gaussian kernels is sharp
\cite[Lemma~9]{R5-PW20}, which already applies to arbitrary continuous
densities, combined with $|\Gamma(1-i\omega)|^2=\pi\omega/\sinh(\pi\omega)$: the
Gompertz kernel has an exponentially decaying Fourier transform and an
exponential left tail, and these balance to give quadratic growth.  For
completeness we include a short proof of the lemma in the form we use.

\begin{lemma}[Sinusoidal windows]\label{R5:lem:window}
Let $h$ be a bounded continuous probability density on $\RfR$ with a bounded
continuous derivative.  Let $H$ be its distribution function and
$\hat h(\omega)=\int e^{-i\omega s}h(s)\Rfd s$.  Let $a,\omega>0$ and suppose
\[
 A:=|\hat h(\omega)|>2T,\qquad T:=H(-a/2)+1-H(a/2).
\]
Let $\pi$ have density proportional to $(1+\sin(\omega y))\mathbf 1_{[-a,a]}(y)$.
Then $\pi*h$ is $C^1$ and has at least $\lfloor a\omega/(2\pi)\rfloor-1$ local
maxima in $(-a/2,a/2)$, each of them a critical point.
\end{lemma}

\begin{proof}
Let $Z$ be the normalizing constant, so that $Z\,(\pi*h)(x)=g(x)-\Delta(x)$
with $g(x)=\int_{\RfR}(1+\sin\omega y)h(x-y)\Rfd y$ and
$\Delta(x)=\int_{|y|>a}(1+\sin\omega y)h(x-y)\Rfd y$.  Write
$\hat h(\omega)=Ae^{i\beta}$.  Then
$\int\sin(\omega y)h(x-y)\Rfd y=\operatorname{Im}\bigl(e^{i\omega x}\hat h(\omega)\bigr)$,
so $g(x)=1+A\sin(\omega x+\beta)$.  For $|x|\le a/2$,
\[
 0\le\Delta(x)\le2\RfP_{S\sim h}(|x-S|>a)=2\bigl[H(x-a)+1-H(x+a)\bigr]\le2T .
\]
Let $u_m=(\pi/2-\beta+2\pi m)/\omega$ and $v_m=u_m+\pi/\omega$.  If
$u_m\in[-a/2,a/2]$, then $Z(\pi*h)(u_m)\ge1+A-2T>1$.  If
$v_m\in[-a/2,a/2]$, then $Z(\pi*h)(v_m)\le1-A<1$.  Consider an interval
$[v_{m-1},v_m]\subset[-a/2,a/2]$.  It contains $u_m$, and the continuous
function $\pi*h$ is larger at $u_m$ than at both endpoints.  So its maximum
over the interval is attained at an interior point, which is a local maximum
on $\RfR$.  Distinct $m$ give disjoint interiors.  The points $v_m$ have
spacing $2\pi/\omega$, so at least $\lfloor a\omega/(2\pi)\rfloor-1$ such
intervals lie in $[-a/2,a/2]$.  Differentiation under the integral shows
that $\pi*h\in C^1$, so local maxima are critical points.
\end{proof}

\begin{proof}[Proof of Theorem~\ref{R5:thm:gompertz}]
The density $h(x)=e^{x-e^x}$ is that of $\log E$ with $E\sim\mathrm{Exp}(1)$.
It is smooth with bounded derivatives, and its distribution function is
$H(x)=1-\exp(-e^x)$.  Substituting $u=e^s$ gives
$\hat h(\omega)=\int_0^\infty u^{-i\omega}e^{-u}\Rfd u=\Gamma(1-i\omega)$.
From $\Gamma(1+z)\Gamma(1-z)=\pi z/\sin(\pi z)$ with $z=i\omega$, and
$\Gamma(1+i\omega)=\overline{\Gamma(1-i\omega)}$, we get
\[
 |\hat h(\omega)|^2=\frac{\pi\omega}{\sinh(\pi\omega)}\ge2\pi\omega\,e^{-\pi\omega}.
\]
Next, $H(-a/2)=1-\exp(-e^{-a/2})\le e^{-a/2}$ and
$1-H(a/2)=\exp(-e^{a/2})\le e^{-a/2}$, so $T\le2e^{-a/2}$.  With
$\omega=\omega_a$ we have $e^{-\pi\omega/2}=4e^{-a/2}$.  Hence
$|\hat h(\omega_a)|\ge4\sqrt{2\pi\omega_a}\,e^{-a/2}>4e^{-a/2}\ge2T$
whenever $2\pi\omega_a>1$, which holds for $a\ge4$.
Lemma~\ref{R5:lem:window} gives at least
$\lfloor a\omega_a/(2\pi)\rfloor-1$ local maxima, as claimed.  The function
$\pi_a*h$ is real-analytic and nonconstant, so these local maxima are
strict.  If the number of critical points were at most $Ca$ for all large
$a$, we would need $a(a-4\log2)/(2\pi^2)-2\le Ca$ for all large $a$, which
is impossible.
\end{proof}

\begin{remark}\label{R5:rem:gompertz}
(1) \emph{Numerical confirmation.}  The script
\texttt{R5\_\allowbreak{}gompertz\_\allowbreak{}check.py} evaluates $\pi_a*h$ exactly through
incomplete Gamma functions, in $60$-digit arithmetic, and cross-checks the
formula against direct quadrature.  Counting local maxima in $(-a/2,a/2)$ on a fine grid
gives $2,6,11,26,47,110$ for $a=8,12,16,24,32,48$.  The proved bounds are
$1,4,9,24,46,108$, and $a^2/(2\pi^2)$ takes the values $3,7,13,29,52,117$.

(2) \emph{A better constant.}  The window can also be placed
asymmetrically, using the double-exponential right tail of $h$.  With
$\omega\approx 2a/\pi$ the numerics give about $a^2/\pi^2$ local maxima
(for example $137$ at $a=36$).  We do not prove this refinement.

(3) \emph{What remains of the open question.}  Riemann sums converge
uniformly to $\pi_a*h$, and the local maxima produced by
Lemma~\ref{R5:lem:window} have a quantitative margin $A-2T$.  Hence the
theorem also holds for finite positive combinations $\sum_iw_ih(x-y_i)$ with
$y_i\in[-a,a]$.  Through $\theta=e^{-x}$ and $x_i=e^{y_i}$, these are, up to
reflection of $h$, exactly the gradient functions
$\theta\mapsto\sum_iw_i\theta e^{-\theta x_i}$ whose critical points carry
the atoms of the exponential-mixture NPMLE \cite[eq.~(35)]{R5-PW20}.  So
the deterministic route proposed in \cite{R5-PW20} cannot give
$O(\log n)$: for suitable positive weights it allows
$\Omega((\log(x_{\max}/x_{\min}))^2)$ critical points.  The open question
itself is not resolved by this.  Its upper bound would have to use the
special NPMLE weights $w_i\propto1/\widehat f(x_i)$ or the randomness of the
data.  We also do not know whether $O(a^2)$ is the correct upper bound for
Gompertz mixtures.
\end{remark}

% ----------------------------------------------------------------------
\subsection{Mixtures of two analytic strongly log-concave densities can have arbitrarily many modes}
\label{R5:sec:modes}

\begin{problem}[Polyanskiy and Wu {\cite[Sec.~5.5]{R5-PW20}}]
\label{R5:prob:modes}
The authors observe that mixtures of two unimodal, or even log-concave,
densities can have infinitely many modes, and that $C^\infty$ strongly
log-concave examples with flat pieces exist.  They then ask: given a convex
combination of $k$ analytic densities that are strongly log-concave, how
many modes can the mixture have?  Their Corollary~11 gives $\Omega(k^2)$
modes, and they write: ``Whether this is tight is an open question.''
\end{problem}

A density $f$ is $c$-strongly log-concave if $(\log f)''\le-c<0$.

\begin{theorem}\label{R5:thm:modes}
Let $\varepsilon\in(0,1)$ and $\omega>0$, and define
\[
 f_\pm(x)=Z^{-1}\exp\Bigl(-\frac{(x\mp1)^2}2+\frac{\varepsilon}{\omega^2}\cos(\omega x)\Bigr),
 \qquad m=\tfrac12(f_++f_-).
\]
Both functions use the same constant $Z$.  Then:
\begin{enumerate}
\item[(i)] $f_+$ and $f_-$ are entire probability densities with
 $-1-\varepsilon\le(\log f_\pm)''\le-1+\varepsilon$.  So they are
 $(1-\varepsilon)$-strongly log-concave and $(1+\varepsilon)$-log-smooth,
 uniformly in $\omega$.
\item[(ii)] The mixture $m$ has at least
 $\frac1\pi\bigl(\frac{3\varepsilon}2\bigr)^{1/3}\omega^{2/3}-\frac52$ strict
 local maxima.
\end{enumerate}
Consequently, for every $k\ge2$ and every $M$ there is a mixture of $k$
distinct analytic, $\tfrac12$-strongly log-concave densities with at least
$M$ modes.  The $\Omega(k^2)$ bound is therefore not tight: no function of
$k$ bounds the number of modes.
\end{theorem}

\begin{proof}
(i) The substitution $x\mapsto-x$ and the evenness of $\cos$ show that the
two normalizing integrals coincide.  The exponents are entire functions of
$x$ and are bounded above by $-(x\mp1)^2/2+\varepsilon/\omega^2$, so
$f_\pm$ are integrable entire functions.  Directly,
$(\log f_\pm)''=-1-\varepsilon\cos(\omega x)$.

(ii) Because the factor $e^{(\varepsilon/\omega^2)\cos\omega x}$ is common to
both components,
\[
 \log m(x)=\mathrm{const}+\frac{\varepsilon}{\omega^2}\cos(\omega x)-\frac{x^2}2+\log\cosh x,
 \qquad
 (\log m)'(x)=-\frac\varepsilon\omega\sin(\omega x)-(x-\tanh x).
\]
For $x\ge0$ we have $0\le x-\tanh x\le x^3/3$.  Indeed, the function
$\tanh x-x+x^3/3$ vanishes at $0$ and has derivative
$x^2-\tanh^2x\ge0$.  Let $X=(3\varepsilon/(2\omega))^{1/3}$, so that
$X^3/3=\varepsilon/(2\omega)$, and let $x_j=(\pi/2+j\pi)/\omega$ for
$j\ge0$.  If $x_j\le X$, then $\sin(\omega x_j)=(-1)^j$ and
$0\le x_j-\tanh x_j\le\varepsilon/(2\omega)$.  Hence $(\log m)'(x_j)$ is
negative for even $j$ and positive for odd $j$.  On each interval
$[x_{2i+1},x_{2i+2}]\subset(0,X]$, the function $m$ increases at the left
end and decreases at the right end, so it has an interior local maximum.
These intervals are disjoint.  Let $J=\#\{j\ge0:x_j\le X\}$.  Then
$J\ge\omega X/\pi-\frac12$, and $(0,X]$ contains at least
$\lfloor(J-1)/2\rfloor$ such maxima.  Since $m$ is even, the same holds on
$[-X,0)$.  Hence there are at least $J-2\ge\omega X/\pi-\frac52$ local
maxima.  The function $m$ is real-analytic and nonconstant, so its critical
points are isolated and its local maxima are strict.

For $k=2$ the mixture $m$ itself, with $\varepsilon=\frac12$ and $\omega$
large, already has at least $M$ modes.  For $k\ge3$, take
$\varepsilon=\frac12$, choose $\omega$ so large that $m$ has at least $M$
strict local maxima in $[-1,1]$, and let $\delta\in(0,1)$.  Consider
$(1-\delta)m+\frac{\delta}{k-2}\sum_{i=3}^k\phi(\cdot-10i)$, where $\phi$ is
the standard normal density.  This is a mixture of $k$ distinct analytic
$\frac12$-strongly log-concave densities.  On $[-1,1]$, the function $m$ is
bounded below and the added term is $O(\delta)$ in $C^1$.  At the finitely
many points $x_j$ used above, the strict sign conditions on $m'$ therefore
persist for small $\delta$, and so do the $M$ local maxima.
\end{proof}

\begin{remark}\label{R5:rem:modes}
(1) The proof locates the extra modes near the degenerate critical point
$x=0$ of the equal mixture of $N(\pm1,1)$, at which
$(\log m)'(x)\sim-x^3/3$.  Analytic wiggles of amplitude
$\varepsilon/\omega^2$ in the log-densities change the second derivatives
by at most $\varepsilon$, but they create about $\omega^{2/3}$ modes.  The
script \texttt{R5\_\allowbreak{}modes\_\allowbreak{}check.py} and an independent high-precision recount give
$9,37,169,785$ strict local maxima for $\omega=10^2,10^3,10^4,10^5$ with
$\varepsilon=\frac12$ (an odd number, as it must be: $m$ is even and $(\log m)''(0)=-\varepsilon<0$).

(2) The construction needs $|(\log f)'''|$ of order $\varepsilon\omega$.  If
one also imposes a uniform bound on the third derivative of $\log f_i$, the
question becomes different, and we leave it open.  The components in the
construction behind \cite[Cor.~11]{R5-PW20} are of a special kind: Gaussian
convolutions of laws on consecutive intervals of length $1/4$.  If all
components are Gaussian convolutions of laws supported in a common interval
of length $O(k)$, the mixture is a Gaussian location mixture, and the
$O(a^2)$ bound of \cite[Sec.~5.4]{R5-PW20} shows that $\Theta(k^2)$ is the
correct order in that structured class.
\end{remark}

% ----------------------------------------------------------------------
\subsection{The optimal constant in Chebyshev moment matching}
\label{R5:sec:favard}

\begin{problem}[Musco, Musco, Rosenblatt and Singh {\cite{R5-MMRS25}}]
\label{R5:prob:favard}
For distributions $p,q$ on $[-1,1]$ whose Chebyshev moments satisfy
$\sum_{j=1}^kj^{-2}(\RfE_pT_j-\RfE_qT_j)^2\le\Gamma^2$, the authors prove
$W_1(p,q)\le36/k+\Gamma$ and write: ``we believe the constants can be
improved, at least to $\frac{2\pi}{k}+\Gamma$, and possibly further''
\cite[footnote~3 and Sec.~3.1]{R5-MMRS25}.  The $36$ is twice the constant
$18$ in their Jackson-type Fact~12.  Citing \cite[App.~C.2]{R5-BKM22}, they
add that ``strong numerical evidence suggests that this 18 can be improved
to $\pi$''.
\end{problem}

\begin{theorem}\label{R5:thm:favard}
Let $k\ge1$.
\begin{enumerate}
\item[(a)] (Favard's damping.)  Let
 $\lambda_j=\frac{j\pi}{2(k+1)}\cot\frac{j\pi}{2(k+1)}$ for $1\le j\le k$ and
 $\lambda_0=1$.  Then $1=\lambda_0>\lambda_1>\dots>\lambda_k>0$.  For every
 $\ell$-Lipschitz $f$ on $[-1,1]$ with Chebyshev series $\sum_ja_jT_j$, the
 damped truncation $f_k=\sum_{j\le k}\lambda_ja_jT_j$ satisfies
 $\|f-f_k\|_\infty\le\pi\ell/(2(k+1))$.
\item[(b)] Under the hypothesis of Problem~\ref{R5:prob:favard},
 $W_1(p,q)\le\pi/(k+1)+\Gamma$.
\item[(c)] The constant $\pi$ in (b) cannot be lowered.  For every $k\ge4$
 there are distributions $p,q$ on $[-1,1]$ with the same first $k$ Chebyshev
 moments (so $\Gamma=0$) and $W_1(p,q)\ge\frac\pi k(1-5k^{-1/2})$.
\item[(d)] (The damping of \cite{R5-BKM22}.)  For the Jackson coefficients
 of \cite{R5-BKM22} ($N=4z$, coefficients $(g*g)*(g*g)$ with
 $g=\mathbf 1_{2z+1}$, normalized), every $1$-Lipschitz $f$ on $[-1,1]$
 satisfies $\|f-\bar f_N\|_\infty\le\pi/N$, for every $z\ge1$.  The best
 constant $c_N=N\sup_f\|f-\bar f_N\|_\infty$ tends to
 $12\log2/\pi\approx2.648$.
\end{enumerate}
\end{theorem}

\begin{remark}[Status]\label{R5:rem:favardstatus}
Part (a) is Favard's theorem for $W^1_\infty$ \cite{R5-Favard37,R5-AK37}
(see also \cite[Ch.~4]{R5-Korneichuk}), transported to $[-1,1]$ by
$x=\cos\theta$.  Part (c) is the matching lower bound, which reduces to the
classical asymptotics of best approximation of Lipschitz functions at the
centre of the interval.  So the answer, ``the optimal constant is $\pi$'',
follows from known results, and we count this item as \emph{settled by
classical theory} rather than as a new theorem.  The only new content is
part (d), which confirms rigorously the numerical observation of
\cite{R5-BKM22}, and the self-contained proofs below.
\end{remark}

\begin{proof}
(a) Let $n=k$ and $N'=n+1$.  For $s\in(0,\pi/2)$ the function $s\cot s$ is
decreasing with values in $(0,1)$, which gives the ordering of the
$\lambda_j$.  Let $B(\theta)=\pi-\theta=2\sum_{j\ge1}\sin(j\theta)/j$ on
$(0,2\pi)$, and let $\sigma(\theta)=2\sum_{j=1}^n\lambda_j\sin(j\theta)/j$.
A direct computation with the discrete sine transform shows that $\sigma$
is the unique sine polynomial of degree $n$ with
$\sigma(j\pi/N')=B(j\pi/N')$ for $1\le j\le2n+1$.  Indeed
$\frac2{N'}\sum_{i=1}^{N'-1}(\pi-\frac{i\pi}{N'})\sin\frac{ij\pi}{N'}
=\frac\pi{N'}\cot\frac{j\pi}{2N'}=2\lambda_j/j$.  Put $\Theta=B-\sigma$ on
$(0,2\pi)$.  It vanishes at the $2n+1$ nodes, while $\Theta(0^+)=\pi$ and
$\Theta(2\pi^-)=-\pi$.  Its derivative $-1-\sigma'$ is a trigonometric
polynomial of degree $n$, so it has at most $2n$ zeros per period.  By
Rolle's theorem these are exhausted between consecutive nodes.  Hence
$\Theta$ has no other zeros, they are all simple, and
$\operatorname{sgn}\Theta=\operatorname{sgn}\sin(N'\theta)$.  Since
$\operatorname{sgn}\sin(N'\theta)=\frac4\pi\sum_{m\ \mathrm{odd}}\sin(mN'\theta)/m$
is orthogonal to $\sigma$,
\[
 \int_0^{2\pi}|\Theta|=\int_0^{2\pi}B(\theta)\operatorname{sgn}\sin(N'\theta)\Rfd\theta
 =\sum_{m\ \mathrm{odd}}\frac{2}{mN'}\cdot\frac{4}{\pi m}\cdot\pi=\frac{\pi^2}{N'} .
\]
Now let $g$ be $2\pi$-periodic with $|g'|\le\ell$ almost everywhere, and let
$F_ng=\sum_j\lambda_j(a_j\cos j\theta+b_j\sin j\theta)$.  Integration by parts
gives $(g-F_ng)(0)=-\frac1{2\pi}\int_0^{2\pi}g'(s)\Theta(s)\Rfd s$, so
$\|g-F_ng\|_\infty\le\ell\pi/(2N')$ by translation invariance.  For $f$
Lipschitz on $[-1,1]$, the function $g(\theta)=f(\cos\theta)$ satisfies
$|g'|\le\ell$, and $F_ng(\theta)=f_k(\cos\theta)$.

(b) This is the proof of \cite[Thm.~1]{R5-MMRS25} with Fact~12 replaced by
(a).  That proof uses Fact~12 only through $|\lambda_j|\le1$ and the uniform
error bound; together with their Lemma~13 it gives
$W_1\le\Gamma+2\|f-f_k\|_\infty$.

(c) Let $D=\lceil\sqrt k\rceil$, $N'=k+1+2D$, $\theta_j=j\pi/N'$,
$x_j=\cos\theta_j$ ($0\le j\le N'$), and let $w_j$ be the Chebyshev--Lobatto
weights ($1/N'$ inside, $1/(2N')$ at the ends).  The Lobatto rule is exact
for polynomials of degree at most $2N'-1$ against the arcsine law, and
$T_{N'}(x_j)=(-1)^j$.  Hence the signed measure
$\mu=\sum_jw_j(-1)^j(1-x_j^2)^D\delta_{x_j}$ annihilates every polynomial of
degree at most $N'-1-2D=k$.  So $\mu=\frac{\|\mu\|}2(p-q)$ for probability
measures $p,q$ with the same first $k$ moments.

Let $s_j=\frac12\min(x_{j-1}-x_j,x_j-x_{j+1})$ for $0<j<N'$, and let $f$ be
the piecewise linear function with $f(x_j)=(-1)^js_j$ and
$f(x_0)=f(x_{N'})=0$.  This $f$ is $1$-Lipschitz, and
\[
 W_1(p,q)\ge\frac{2}{\|\mu\|}\int f\Rfd\mu=\frac{2\sum_jw_j\sin^{2D}\theta_j\,s_j}{\sum_jw_j\sin^{2D}\theta_j}.
\]
Here
\[
 s_j=\sin\frac\pi{2N'}\Bigl(\sin\theta_j\cos\frac\pi{2N'}-|\cos\theta_j|\sin\frac\pi{2N'}\Bigr).
\]
Lobatto exactness gives
$\sum w_j\sin^{2D+2}\theta_j/\sum w_j\sin^{2D}\theta_j=(2D+1)/(2D+2)$, and
$\sin\theta\ge\sin^2\theta$.  Therefore
$W_1(p,q)\ge\sin\frac\pi{N'}\cdot\frac{2D+1}{2D+2}-2\sin^2\frac\pi{2N'}$.
For $k\ge4$ this is at least $\frac\pi k(1-5k^{-1/2})$, by elementary
estimates with $N'\le k+3+2\sqrt k$.

(d) The BKM kernel is $K(\theta)=C_\nu^{-1}(\sin(\nu\theta/2)/\sin(\theta/2))^4$
with $\nu=2z+1$ and $C_\nu=\nu(2\nu^2+1)/3$.  It has degree $N=2\nu-2$ and is
nonnegative.  For a nonnegative kernel the worst $1$-Lipschitz $f$ at
$x_0=\cos\theta_0$ is $|x-x_0|$.  Moreover
$|\cos\theta_0-\cos(\theta_0-s)|\le2|\sin(s/2)|$.  Hence
\[
 \sup_f\|f-\bar f_N\|_\infty\le\frac1{2\pi}\int_{-\pi}^\pi2|\sin\tfrac s2|K(s)\Rfd s
 =\frac{4}{\pi C_\nu}I(\nu),\qquad I(\nu)=\int_0^{\pi/2}\frac{\sin^4(\nu u)}{\sin^3u}\Rfd u .
\]
On $(0,1]$ we have $\sin^{-3}u\le u^{-3}+u^{-1}$, because
$\sin u\ge u(1-u^2/6)$ and $(1-u^2/6)^{-3}\le(1-u^2/2)^{-1}\le1+u^2$.  On
$[1,\pi/2]$ we have $\sin^{-3}u\le1.678$.  Two integrations by parts and
Frullani's integral give
$\int_0^\infty\sin^4v\,v^{-3}\Rfd v=\int_0^\infty(\cos2v-\cos4v)v^{-1}\Rfd v=\log2$,
since $(\sin^4)''(v)=2\cos2v-2\cos4v$.  Also
$\int_0^1\sin^4(\nu u)u^{-1}\Rfd u\le\frac14+\log\nu$.  Together these give
$I(\nu)\le\nu^2\log2+\log\nu+1.21$.  Since
$(2\nu-2)/(\nu(2\nu^2+1))\le\nu^{-2}$,
\[
 N\sup_f\|f-\bar f_N\|_\infty\le\frac{12}{\pi}\Bigl(\log2+\frac{\log\nu+1.21}{\nu^2}\Bigr)
 \le3.08<\pi\qquad(\nu\ge5).
\]
For $\nu=3$ ($N=4$), exact evaluation gives $1.807$.

The displayed bound also gives $\limsup c_N\le12\log2/\pi$.  For the
reverse inequality, test $f=|x|$ at $\theta_0=\pi/2$.  Since
$|\sin s|\ge2|\sin\frac s2|-|s|^3/8$, we get
$c_N\ge N\bigl(\frac4{\pi C_\nu}I(\nu)-\frac1{16\pi}\int_{-\pi}^{\pi}|s|^3K(s)\Rfd s\bigr)$.
Moreover $I(\nu)\ge\nu^2\log2-2/\pi^2$.  The kernel satisfies
$K(s)\le C\nu(\nu|s|)^{-4}$, so $\int|s|^3K=O(\nu^{-3}\log\nu)$.  Hence
$\liminf c_N\ge12\log2/\pi$.
The script \texttt{R5\_\allowbreak{}jackson\_\allowbreak{}check.py} evaluates the sharp constant for
$N\le192$: it increases from $1.635$ to $2.620$.
\end{proof}

% ----------------------------------------------------------------------
\subsection{Candidates examined and found settled elsewhere}
\label{R5:sec:settled}

The following explicitly posed questions in this area were checked and
found to be already answered by other authors, so they were not pursued.
\begin{itemize}
\item Polyanskiy--Wu conjecture that the statistical degree of Gaussian
 mixtures with compactly supported mixing law is $\asymp\log n/\log\log n$
 \cite[Sec.~5.3]{R5-PW20}: resolved by Ma, Wu and Yang \cite{R5-MWY24}.
\item Polyanskiy--Wu conjecture that the multivariate Gaussian NPMLE is
 unique \cite[Sec.~5.6]{R5-PW20}: non-uniqueness for an equilateral
 triangle in $\RfR^2$ \cite[Lemma~2]{R5-SGS25}, and uniqueness for almost
 every data set \cite{R5-Wang26}.
\item Minimax rates for smooth binomial mixing densities with heterogeneous
 trials, stated as open in \cite[Sec.~2.3, Sec.~6]{R5-LBK25}: determined
 by Lee and Kennedy \cite{R5-LK26}.  Adaptivity to the smoothness remains
 listed there as future work.
\item The Vinayak et al.\ Chebyshev--Bernstein coefficient conjecture and
 the binomial NPMLE $W_1$ range, as well as Poisson-mixture rates, are
 treated in earlier parts of this programme \cite{K-029}.  The
 fixed-test ``fingerprint'' conjecture of \cite[Sec.~5.1]{R5-VKVK19} remains
 open and is not addressed here.
\end{itemize}


%% ==================================================================
%% Section: Polygamma Wronskians
%% ==================================================================
% ===================== Section R2 =====================
% Area R2: completely monotonic degrees and related conjectures on polygamma functions.
% Requires: amsmath, amssymb, amsthm; environments theorem, lemma, proposition,
% corollary, definition, remark, conjecture, problem, proof.
\providecommand{\RtwoL}{\mathcal{L}}
\providecommand{\Rtwodeg}{\operatorname{deg}_{\mathrm{cm}}}
\providecommand{\RtwoB}{\operatorname{B}}

\section{Wronskians of polygamma products and two degree questions}
\label{R2:sec}

This section settles several explicitly posed guesses and questions of
F.~Qi and coauthors on complete monotonicity (CM) of functions built from
polygamma functions, answers one CM-degree question from the survey
\cite{R2-QiAgarwal2019} and partly answers another. All proofs follow the pattern of the degree
theorem for Stirling remainders (\cite[Part~VIII]{K-033}): write the
function as a Laplace transform, reduce CM to the sign of an explicit density
(Bernstein's theorem plus uniqueness), and decide that sign. Monotonicity of
ratios of Laplace transforms is standard; what we add for the guesses is a
likelihood-ratio argument for convolution densities, driven by the
log-concavity of the Bernoulli kernel $c(t)=t/(1-e^{-t})$ and of its
derivative, and by the concavity of $tc'(t)/c(t)$.

\medskip\noindent\textbf{Summary of outcomes.}
\begin{itemize}
\item Qi, C.~R.~Math.\ 2022, Remark~17 \cite{R2-QiCR2022}: guess \emph{true};
 the attached necessary-and-sufficient-condition problem is solved with
 explicit constants (Theorem~\ref{R2:thm:A}). We point out
 (Remark~\ref{R2:rem:A-split}) that the whole remark also follows from Qi's own
 Theorem~1 in \cite{R2-QiCR2022} by a splitting identity, so this closure is
 easy rather than deep.
\item Qi, C.~R.~Math.\ 2022, Remark~18 \cite{R2-QiCR2022}: guess \emph{true}
 (Theorem~\ref{R2:thm:B}). Attached problem: one half solved explicitly
 ($\Omega\ge1$); the other half is reduced to an explicit one-variable
 minimum $\Omega_*\in(0,1)$, computed numerically (\emph{partial} for that
 half: no closed form).
\item Yin--Han--Qi, Demonstr.\ Math.\ 2024, Remark~4 \cite{R2-YHQ2024}:
 guess \emph{true} (Theorem~\ref{R2:thm:C}). Attached problem: one half
 solved ($\Lambda\le1$), the other half reduced to an explicit maximum
 $\Lambda^*_{m,n}>1$ (\emph{partial}: no closed form). As for Remark~17, the
 guess itself can also be obtained by splitting from a published majorization
 theorem (Remark~\ref{R2:rem:C}(2)); our proof is independent of it.
\item Yin--Han--Qi, Remark~2 \cite{R2-YHQ2024}: both sign guesses
 \emph{true}, with the complete sign pattern (Proposition~\ref{R2:prop:D};
 elementary).
\item Qi--Agarwal 2019, Section~5.4 \cite{R2-QiAgarwal2019}:
 $\Rtwodeg H_2=0$ and $\Rtwodeg H_3=1$ (Theorem~\ref{R2:thm:E}).
\item Qi--Agarwal 2019, Section~5.5 \cite{R2-QiAgarwal2019}: the negative
 has degree exactly $1$ for every $\mu\ge4$; for $\lambda=0$ the degree is
 $0$; for every $\lambda<0$ the degree is a non-integer in $(0,1)$,
 nonincreasing in $\lambda$, tending to $1$ and to $0$ at the two ends
 (Theorem~\ref{R2:thm:F}). For $\lambda<0$ the degree is therefore
 located in $(0,1)$ but not determined, so this question is answered only
 \emph{partially}.
\item Qi, Math.\ Inequal.\ Appl.\ 2021, Remark~4.1 \cite{R2-QiMIA2021}:
 \emph{partial} only (Proposition~\ref{R2:prop:G}).
\end{itemize}

% ------------------------------------------------------------------
\subsection{Notation and general tools}
\label{R2:subsec:tools}

A function $f\colon(0,\infty)\to\mathbb R$ is \emph{completely monotonic}
(CM) if $f\in C^\infty$ and $(-1)^kf^{(k)}\ge0$ for all $k\ge0$. Products,
sums, nonnegative multiples and pointwise limits (with all derivatives) of
CM functions are CM. Following \cite{R2-QiAgarwal2019}, for a CM function
$f\not\equiv0$ with $f(\infty)=0$ (the only case needed here; in general
\cite{R2-QiAgarwal2019} subtracts $f(\infty)$) the \emph{completely
monotonic degree} is
\[
 \Rtwodeg f=\sup\{\alpha\in\mathbb R:\ x^\alpha f(x)\text{ is CM on }(0,\infty)\}.
\]

\begin{lemma}\label{R2:lem:degset}
Let $f\not\equiv0$ be CM. The set $A_f=\{\alpha:\ x^\alpha f\text{ is CM}\}$
is of the form $(-\infty,d]$ or $(-\infty,\infty)$, and $d=\Rtwodeg f$ when
it is finite. If $x^{\alpha}f(x)\to0$ as $x\to0^+$, then $\alpha\notin A_f$.
\end{lemma}
\begin{proof}
If $\alpha\in A_f$ and $\beta<\alpha$, then $x^\beta f=x^{\beta-\alpha}\cdot
x^\alpha f$ is a product of CM functions. If $\alpha_k\uparrow\alpha$ with
$\alpha_k\in A_f$, then $x^{\alpha_k}f\to x^\alpha f$ locally uniformly on
$(0,\infty)$ together with all derivatives, so the sign conditions pass to
the limit. Finally, a CM function is nonnegative and nonincreasing; if it
tends to $0$ at $0^+$ it vanishes identically, which is impossible for
$x^\alpha f$.
\end{proof}

\begin{lemma}[Laplace criterion]\label{R2:lem:bern}
Let $\varphi\colon(0,\infty)\to\mathbb R$ be continuous with
$\int_0^\infty|\varphi(t)|e^{-xt}\,dt<\infty$ for every $x>0$, and let
$F(x)=\RtwoL\varphi(x)=\int_0^\infty\varphi(t)e^{-xt}dt$. Then $F$ is CM
on $(0,\infty)$ if and only if $\varphi\ge0$ on $(0,\infty)$. If
$\varphi\ge0$ and $\varphi\not\equiv0$, then $(-1)^kF^{(k)}>0$ for all $k$.
\end{lemma}
\begin{proof}
If $\varphi\ge0$, differentiation under the integral sign (dominated
convergence) gives $(-1)^kF^{(k)}(x)=\int_0^\infty t^k\varphi(t)e^{-xt}dt$.
Conversely, by Bernstein's theorem $F(x)=\int_{[0,\infty)}e^{-xt}d\mu(t)$
with a positive measure $\mu$. For fixed $a>0$ the finite signed measures
$e^{-at}\varphi(t)dt$ and $e^{-at}d\mu(t)$ have equal Laplace transforms on
$[0,\infty)$, hence coincide \cite[Ch.~II, \S6]{R2-Widder}. Thus
$\varphi\,dt=\mu$ on $(0,\infty)$, so $\varphi\ge0$ almost everywhere, and
everywhere by continuity.
\end{proof}

\begin{lemma}[Wronskian identity]\label{R2:lem:wronski}
Let $\alpha,\gamma\colon(0,\infty)\to(0,\infty)$ be continuous with
$\int_0^\infty(1+t)(\alpha+\gamma)(t)e^{-xt}dt<\infty$ for all $x>0$, and
let $A=\RtwoL\alpha$, $C=\RtwoL\gamma$. Put
\[
 M_1(t)=\int_0^t\alpha(u)\gamma(t-u)\,u\,du,\qquad
 M_2(t)=\int_0^t\alpha(u)\gamma(t-u)(t-u)\,du .
\]
Then $AC'=-\RtwoL M_2$, $A'C=-\RtwoL M_1$, so $AC'-A'C=\RtwoL[M_1-M_2]$, and
\begin{equation}\label{R2:eq:symm}
 M_1(t)-M_2(t)=\int_0^{t/2}(t-2u)\bigl[\alpha(t-u)\gamma(u)-\alpha(u)\gamma(t-u)\bigr]du .
\end{equation}
If $\alpha/\gamma$ is nondecreasing, then $M_1\ge M_2$ and $AC'-A'C$ is CM;
if $\alpha/\gamma$ is strictly increasing, then $M_1>M_2$ on $(0,\infty)$.
\end{lemma}
\begin{proof}
$A'=-\RtwoL[t\alpha]$ and $C'=-\RtwoL[t\gamma]$, and the convolution theorem
(Fubini, absolute convergence) gives $AC'=-\RtwoL[\alpha*(t\gamma)]=-\RtwoL M_2$
and $A'C=-\RtwoL[(t\alpha)*\gamma]=-\RtwoL M_1$. Splitting $[0,t]$ at $t/2$ and
substituting $u\mapsto t-u$ on $[t/2,t]$ gives \eqref{R2:eq:symm}. For
$0<u<t/2$ we have $t-u>u$, so monotonicity of $\alpha/\gamma$ gives
$\alpha(t-u)\gamma(u)\ge\alpha(u)\gamma(t-u)$ (strictly in the strict case).
Lemma~\ref{R2:lem:bern} finishes the proof.
\end{proof}

\begin{lemma}[likelihood ratio]\label{R2:lem:lr}
Let $\nu_1,\nu_2$ be finite positive measures on an interval $J$, with
$\nu_2=\rho\,\nu_1$ for a nondecreasing $\rho\ge0$, and let $h$ be bounded
and nondecreasing on $J$. Then
$\int h\,d\nu_2/\nu_2(J)\ge\int h\,d\nu_1/\nu_1(J)$. If $\nu_1$ has a
positive density and $\rho$, $h$ are strictly increasing, the inequality is
strict.
\end{lemma}
\begin{proof}
$\iint[h(v)-h(w)][\rho(v)-\rho(w)]\,d\nu_1(v)\,d\nu_1(w)\ge0$
(Chebyshev), which expands to
$\nu_1(J)\int h\rho\,d\nu_1\ge\int h\,d\nu_1\int\rho\,d\nu_1$.
\end{proof}

\noindent\textbf{The kernels.} For $t>0$ put
\begin{gather*}
 c(t)=\frac{t}{1-e^{-t}},\qquad
 b(t)=c'(t)=\frac{e^t(e^t-1-t)}{(e^t-1)^2},\\
 \ell(t)=\frac{tc'(t)}{c(t)}=1-\frac{t}{e^t-1}.
\end{gather*}
For $k\ge1$ \cite[5.15.1]{R2-DLMF}
\begin{equation}\label{R2:eq:psirep}
 |\psi^{(k)}(x)|=(-1)^{k+1}\psi^{(k)}(x)=\int_0^\infty t^{k-1}c(t)e^{-xt}dt ,
\end{equation}
and, integrating by parts in $x\psi'(x)=x\RtwoL[c](x)$ (note $c(0^+)=1$),
\begin{equation}\label{R2:eq:Phirep}
 \Phi(x):=x\psi'(x)-1=\RtwoL[b](x),\qquad (-1)^k\Phi^{(k)}=\RtwoL[t^kb].
\end{equation}
Moreover $|\psi^{(k)}(x)|=k!\,Z_k(x)$ with the Hurwitz function
$Z_k(x)=\zeta(k+1,x)=\sum_{j\ge0}(x+j)^{-k-1}$; here $Z_k'=-(k+1)Z_{k+1}$ and
\cite[25.11.43]{R2-DLMF}
\begin{equation}\label{R2:eq:Zasy}
 Z_k(x)=x^{-k-1}+O(1)\ \ (x\to0^+),\qquad
 Z_k(x)=\frac{x^{-k}}{k}\bigl(1+O(x^{-1})\bigr)\ \ (x\to\infty).
\end{equation}

\begin{lemma}\label{R2:lem:kernels}
On $(0,\infty)$:
\begin{enumerate}
\item[(a)] $c$ is strictly log-concave:
 $(\log c)''(t)=\dfrac{1}{4\sinh^2(t/2)}-\dfrac1{t^2}<0$;
\item[(b)] $\ell$ is strictly concave:
 $\ell''(t)=-\dfrac{e^t[(t-2)e^t+t+2]}{(e^t-1)^3}$ and
 $(t-2)e^t+t+2=\sum_{k\ge3}\frac{(k-2)t^k}{k!}>0$;
\item[(c)] $b$ is strictly log-concave, increasing, and $\tfrac12\le b<1$.
\end{enumerate}
\end{lemma}
\begin{proof}
(a) $\log c=\log t-\log(1-e^{-t})$ and $\sinh y>y$ for $y>0$.
(b) $\ell=1-f$ with $f(t)=t/(e^t-1)$; differentiation gives the displayed
$f''$, and the series follows from $e^t=\sum t^k/k!$.
(c) Since $\log b=t+\log(e^t-1-t)-2\log(e^t-1)$,
\[
 (\log b)''=-\frac{1+(t-1)e^t}{(e^t-1-t)^2}+\frac{2e^t}{(e^t-1)^2}
 =-\frac{P(t)}{(e^t-1-t)^2(e^t-1)^2},
\]
where
\begin{gather*}
 P(t)=(t-3)e^{3t}+(2t+7)e^{2t}-(2t^2+3t+5)e^t+1=\sum_{n\ge1}c_n\frac{t^n}{n!},\\
 c_n=3^{n-1}(n-9)+2^n(n+7)-(2n^2+n+5).
\end{gather*}
One computes $c_1=\dots=c_5=0$, $c_6=20$, $c_7=224$, $c_8=1512$; for
$n\ge9$, $3^{n-1}(n-9)\ge0$ and $2^n(n+7)\ge n^2(n+7)>2n^2+n+5$ (because
$n^3+5n^2-n-5=(n+5)(n^2-1)>0$). Hence $P>0$ on $(0,\infty)$. Finally
$c=t+f$, so $b'=c''=f''>0$ by (b): $b$ increases from $b(0^+)=\frac12$ to
$b(\infty)=1$; directly, $b<1\iff e^t(1-t)<1$, which holds for $t>0$.
\end{proof}

\begin{lemma}[symmetrized Beta weights]\label{R2:lem:cosh}
For integers $a,b\ge1$ let
$\sigma_{a,b}(v)=v^{a-1}(1-v)^{b-1}+v^{b-1}(1-v)^{a-1}$, $0<v<1$.
If $m,n,p,q\ge1$, $m+n=p+q$ and $\max(p,q)>\max(m,n)$, then
$\sigma_{m,n}/\sigma_{p,q}$ is strictly increasing on $(0,\tfrac12]$, with
value $1$ at $v=\tfrac12$.
\end{lemma}
\begin{proof}
With $z=v/(1-v)\in(0,1]$ one has
$\sigma_{a,b}(v)=(1-v)^{a+b-2}z^{\min(a,b)-1}(1+z^{|a-b|})$. Put $d=|m-n|$
and $k=\min(m,n)-\min(p,q)\ge1$; then $|p-q|=d+2k$ and, with $z=e^{-\tau}$,
\[
 \frac{\sigma_{m,n}}{\sigma_{p,q}}=\frac{z^k(1+z^d)}{1+z^{d+2k}}
 =\frac{\cosh(d\tau/2)}{\cosh((d/2+k)\tau)} .
\]
Its logarithmic $\tau$-derivative is
$\frac d2\tanh\frac{d\tau}2-(\frac d2+k)\tanh((\frac d2+k)\tau)<0$ for
$\tau>0$, because $y\mapsto y\tanh(y\tau)$ is strictly increasing on
$[0,\infty)$. As $v$ increases to $\frac12$, $\tau=\log\frac{1-v}{v}$
decreases to $0$.
\end{proof}

\begin{lemma}\label{R2:lem:hurwitz}
For integers $r>s\ge1$ the function $Z_{r+1}Z_s-Z_rZ_{s+1}$ is CM on
$(0,\infty)$.
\end{lemma}
\begin{proof}
By absolute convergence,
\[
 Z_{r+1}Z_s-Z_rZ_{s+1}=\sum_{i,j\ge0}\frac{(x+j)-(x+i)}{(x+i)^{r+2}(x+j)^{s+2}} ,
\]
and symmetrizing in $(i,j)$ gives
\[
 Z_{r+1}Z_s-Z_rZ_{s+1}=\sum_{0\le i<j}\frac{j-i}{(x+i)^{s+2}(x+j)^{s+2}}
 \Bigl[(x+i)^{-(r-s)}-(x+j)^{-(r-s)}\Bigr].
\]
For $i<j$ the bracket equals
$\RtwoL\bigl[t^{r-s-1}(e^{-it}-e^{-jt})\bigr]/(r-s-1)!$, so it is CM; each
summand is a product of CM functions, and the series converges with all
derivatives locally uniformly.
\end{proof}

% ------------------------------------------------------------------
\subsection{Qi's Remark~17: the three-polygamma Wronskian}
\label{R2:subsec:A}

\begin{problem}[{Qi \cite[Remark~17]{R2-QiCR2022}}]\label{R2:prob:A}
``The decreasing property of $Q_{m,n}(x)$ in Theorem~11 implies that the inequality
\[
 \psi^{(m+n)}(x)[\psi^{(m)}(x)\psi^{(n)}(x)]'-\psi^{(m+n+1)}(x)\psi^{(m)}(x)\psi^{(n)}(x)>0
\]
\dots is valid on $(0,\infty)$ for $m,n\in\mathbb N$. We guess that, for
$m,n\in\mathbb N$, the function
\[
 \psi^{(m+n)}(x)[\psi^{(m)}(x)\psi^{(n)}(x)]'-\psi^{(m+n+1)}(x)\psi^{(m)}(x)\psi^{(n)}(x)
\]
should be completely monotonic in $x\in(0,\infty)$. Generally, one can
discuss necessary and sufficient conditions on $\Omega_{m,n}\in\mathbb R$
such that the function
\[
 \psi^{(m+n)}(x)[\psi^{(m)}(x)\psi^{(n)}(x)]'-\Omega_{m,n}\psi^{(m+n+1)}(x)\psi^{(m)}(x)\psi^{(n)}(x)
\]
and its opposite are respectively completely monotonic on $(0,\infty)$.''
Here $Q_{m,n}$ denotes the quotient $\psi^{(m+n)}/(\psi^{(m)}\psi^{(n)})$.
\end{problem}

\begin{theorem}\label{R2:thm:A}
Let $m,n\in\mathbb N$, $M=m+n$, $\Omega\in\mathbb R$, and
$H_\Omega=\psi^{(M)}[\psi^{(m)}\psi^{(n)}]'-\Omega\,\psi^{(M+1)}\psi^{(m)}\psi^{(n)}$.
Then
\begin{enumerate}
\item[(i)] $H_\Omega$ is CM on $(0,\infty)$ if and only if $\Omega\le1$; in
 particular Qi's guess ($\Omega=1$) holds;
\item[(ii)] $-H_\Omega$ is CM on $(0,\infty)$ if and only if
 $\Omega\ge\Omega^\sharp:=\dfrac{M+2}{M+1}$.
\end{enumerate}
\end{theorem}

\begin{proof}
Since $|\psi^{(k)}|=k!Z_k$ and $Z_k'=-(k+1)Z_{k+1}$, a sign count gives
\begin{multline}\label{R2:eq:HZ}
 -H_\Omega=M!\,m!\,n!\Bigl[\Omega(M+1)Z_{M+1}Z_mZ_n\\
 -(m+1)Z_MZ_{m+1}Z_n-(n+1)Z_MZ_mZ_{n+1}\Bigr].
\end{multline}

\emph{Sufficiency in (i).} By \eqref{R2:eq:psirep},
$P:=|\psi^{(m)}\psi^{(n)}|=\RtwoL\alpha$ and $B:=|\psi^{(M)}|=\RtwoL\beta$
with $\alpha=a_m*a_n$, $\beta=a_M$, $a_k(t)=t^{k-1}c(t)$. Also
$\psi^{(m)}\psi^{(n)}=(-1)^MP$, $\psi^{(M)}=(-1)^{M+1}B$ and
$\psi^{(M+1)}=(-1)^{M+1}B'$, whence $H_1=PB'-P'B$. By
Lemma~\ref{R2:lem:wronski} (with $(A,C)=(P,B)$) it suffices that
$\alpha/\beta$ be increasing. Now
\begin{gather*}
 \frac{\alpha(s)}{\beta(s)}=\int_0^1v^{m-1}(1-v)^{n-1}e^{L_v(s)}dv,\\
 L_v(s)=\log c(sv)+\log c(s(1-v))-\log c(s),
\end{gather*}
and $L_v'(s)=v(\log c)'(sv)+(1-v)(\log c)'(s(1-v))-(\log c)'(s)>0$,
because $(\log c)'$ is strictly decreasing (Lemma~\ref{R2:lem:kernels}(a))
and $sv,\,s(1-v)<s$. So $H_1$ is CM (indeed strictly). For $\Omega<1$,
$H_\Omega=H_1+(1-\Omega)\,|\psi^{(M+1)}|\,|\psi^{(m)}\psi^{(n)}|$ is a sum
of CM functions.

\emph{Necessity in (i).} By \eqref{R2:eq:Zasy} at infinity, the bracket in
\eqref{R2:eq:HZ} equals
\[
 \Bigl(\frac{\Omega}{mn}-\frac{1}{Mn}-\frac{1}{Mm}\Bigr)x^{-2M-1}+O(x^{-2M-2})
 =\frac{\Omega-1}{mn}\,x^{-2M-1}+O(x^{-2M-2}).
\]
Hence $H_\Omega(x)=M!(m-1)!(n-1)!(1-\Omega)x^{-2M-1}+O(x^{-2M-2})$, which is
negative for large $x$ when $\Omega>1$.

\emph{Sufficiency in (ii).} At $\Omega^\sharp$ write $M+2=(m+1)+(n+1)$ in
\eqref{R2:eq:HZ}:
\begin{align*}
 -H_{\Omega^\sharp}&=M!\,m!\,n!\Bigl[(m+1)Z_n\bigl(Z_{M+1}Z_m-Z_MZ_{m+1}\bigr)\\
 &\qquad\qquad+(n+1)Z_m\bigl(Z_{M+1}Z_n-Z_MZ_{n+1}\bigr)\Bigr].
\end{align*}
Since $M>m$ and $M>n$, Lemma~\ref{R2:lem:hurwitz} shows that
$-H_{\Omega^\sharp}$ is CM. For $\Omega>\Omega^\sharp$ add the CM function
$(\Omega-\Omega^\sharp)(M+1)!\,m!\,n!\,Z_{M+1}Z_mZ_n$.

\emph{Necessity in (ii).} By \eqref{R2:eq:Zasy} at $0^+$, the bracket in
\eqref{R2:eq:HZ} is $[\Omega(M+1)-(M+2)]x^{-2M-4}(1+o(1))$, which is negative
near $0$ if $\Omega<\Omega^\sharp$.
\end{proof}

\begin{remark}\label{R2:rem:A-split}
Remark~17 is in fact a consequence of Theorem~1 of \cite{R2-QiCR2022}
(that is, \cite[Thm.~3.1]{R2-QiRM2021}): for $p'>m'\ge n'>q'\ge1$ with
$m'+n'=p'+q'$, the function
$|\psi^{(m')}\psi^{(n')}|-c\,|\psi^{(p')}\psi^{(q')}|$ is CM iff
$c\le\frac{(m'-1)!(n'-1)!}{(p'-1)!(q'-1)!}$, and its negative is CM iff
$c\ge\frac{m'!n'!}{p'!q'!}$. Indeed, splitting
$|\psi^{(M+1)}\psi^{(m)}\psi^{(n)}|$ with weights $\frac mM+\frac nM=1$,
\begin{align*}
 H_1&=|\psi^{(n)}|\Bigl(|\psi^{(M)}\psi^{(m+1)}|-\tfrac mM|\psi^{(M+1)}\psi^{(m)}|\Bigr)\\
 &\quad+|\psi^{(m)}|\Bigl(|\psi^{(M)}\psi^{(n+1)}|-\tfrac nM|\psi^{(M+1)}\psi^{(n)}|\Bigr),
\end{align*}
and $\frac mM=\frac{(M-1)!\,m!}{M!\,(m-1)!}$ is exactly the threshold for
the pair $(M+1,m)\succ(M,m+1)$. Likewise part~(ii) at $\Omega^\sharp$
follows from the ``negative'' half of the same theorem (equivalently, from
Lemma~\ref{R2:lem:hurwitz}). We have not seen this observation in print; the
two-term inequalities used are also contained in Z.-H.~Yang's majorization
theorem \cite[Prop.~2]{R2-Yang2025}. The proof given above is independent of
both.
\end{remark}

% ------------------------------------------------------------------
\subsection{Qi's Remark~18: the four-polygamma Wronskian}
\label{R2:subsec:B}

\begin{problem}[{Qi \cite[Remark~18]{R2-QiCR2022}}]\label{R2:prob:B}
``The decreasing property of $Q_{m,n;p,q}(x)$ in Theorem~12 implies that the
inequality
\[
 [\psi^{(m)}(x)\psi^{(n)}(x)][\psi^{(p)}(x)\psi^{(q)}(x)]'
 -[\psi^{(m)}(x)\psi^{(n)}(x)]'[\psi^{(p)}(x)\psi^{(q)}(x)]>0
\]
\dots is valid on $(0,\infty)$ for $(p,q)\succ(m,n)$. We guess that, for
$(p,q)\succ(m,n)$, the function
\[
 [\psi^{(m)}(x)\psi^{(n)}(x)][\psi^{(p)}(x)\psi^{(q)}(x)]'
 -[\psi^{(m)}(x)\psi^{(n)}(x)]'[\psi^{(p)}(x)\psi^{(q)}(x)]
\]
should be completely monotonic in $x\in(0,\infty)$. Generally, for
$(p,q)\succ(m,n)$, one can discuss necessary and sufficient conditions on
$\Omega_{m,n;p,q}\in\mathbb R$ such that the function
\[
 [\psi^{(m)}\psi^{(n)}][\psi^{(p)}\psi^{(q)}]'
 -\Omega_{m,n;p,q}[\psi^{(m)}\psi^{(n)}]'[\psi^{(p)}\psi^{(q)}]
\]
and its opposite are respectively completely monotonic on $(0,\infty)$.''
Here $Q_{m,n;p,q}=\psi^{(m)}\psi^{(n)}/(\psi^{(p)}\psi^{(q)})$,
$m,n,p,q\in\mathbb N$, and $(p,q)\succ(m,n)$ means $p+q=m+n$ and
$\max(p,q)>\max(m,n)$.
\end{problem}

\begin{theorem}\label{R2:thm:B}
Let $m,n,p,q\in\mathbb N$ with $(p,q)\succ(m,n)$, $M=m+n$, and
\[
 W_\Omega=[\psi^{(m)}\psi^{(n)}][\psi^{(p)}\psi^{(q)}]'-\Omega[\psi^{(m)}\psi^{(n)}]'[\psi^{(p)}\psi^{(q)}].
\]
Let $\alpha=a_m*a_n$, $\gamma=a_p*a_q$ ($a_k=t^{k-1}c$), and let $M_1,M_2$ be
as in Lemma~\ref{R2:lem:wronski}.
\begin{enumerate}
\item[(i)] $W_\Omega$ is CM on $(0,\infty)$ if and only if $\Omega\ge1$; in
 particular Qi's guess ($\Omega=1$) holds.
\item[(ii)] $-W_\Omega$ is CM on $(0,\infty)$ if and only if
 $\Omega\le\Omega_*$, where $\Omega_*=\Omega_*(m,n;p,q)=\min_{t>0}M_2(t)/M_1(t)$;
 the minimum exists and $0<\Omega_*<1$.
\end{enumerate}
\end{theorem}

\begin{proof}
Put $P=|\psi^{(m)}\psi^{(n)}|=\RtwoL\alpha$ and
$C=|\psi^{(p)}\psi^{(q)}|=\RtwoL\gamma$. Both products carry the sign
$(-1)^M$, so $W_\Omega=PC'-\Omega P'C=\RtwoL[\Omega M_1-M_2]$
(Lemma~\ref{R2:lem:wronski}).

\emph{Step 1: $\alpha/\gamma$ is strictly increasing.} With
$K_s(v)=c(sv)c(s(1-v))$, which is symmetric under $v\mapsto1-v$,
\[
 \frac{\alpha(s)}{\gamma(s)}
 =\frac{\int_0^1v^{m-1}(1-v)^{n-1}K_s(v)dv}{\int_0^1v^{p-1}(1-v)^{q-1}K_s(v)dv}
 =\frac{\int_0^{1/2}\sigma_{m,n}K_s\,dv}{\int_0^{1/2}\sigma_{p,q}K_s\,dv}
 =\frac{\int_J h\,d\nu_s}{\nu_s(J)},
\]
where $J=(0,\frac12]$, $d\nu_s=\sigma_{p,q}K_s\,dv$ and
$h=\sigma_{m,n}/\sigma_{p,q}$, which is bounded and strictly increasing on
$J$ by Lemma~\ref{R2:lem:cosh}. For $s<s'$,
$d\nu_{s'}/d\nu_s=K_{s'}/K_s=\exp\int_s^{s'}\partial_\sigma\log K_\sigma\,d\sigma$, and
\begin{gather*}
 \partial_\sigma\log K_\sigma(v)=\frac{\ell(\sigma v)+\ell(\sigma(1-v))}{\sigma},\\
 \partial_v\bigl[\ell(\sigma v)+\ell(\sigma(1-v))\bigr]=\sigma\bigl[\ell'(\sigma v)-\ell'(\sigma(1-v))\bigr]>0
\end{gather*}
for $0<v<\frac12$, since $\ell'$ is strictly decreasing
(Lemma~\ref{R2:lem:kernels}(b)). So $K_{s'}/K_s$ is strictly increasing on
$J$, and Lemma~\ref{R2:lem:lr} gives the claim.

\emph{Step 2: proof of (i).} By Step~1 and Lemma~\ref{R2:lem:wronski},
$M_1>M_2$, so $W_1$ is CM. For $\Omega>1$, $W_\Omega=W_1+(\Omega-1)(-P')C$,
and $-P'=\RtwoL[t\alpha]$ is CM. For $\Omega<1$: as $u\to0^+$,
$\alpha(u)=u^{M-1}(\RtwoB(m,n)+o(1))$ and
$\gamma(u)=u^{M-1}(\RtwoB(p,q)+o(1))$ (because $c(0^+)=1$), hence
\[
 \frac{M_2(t)}{M_1(t)}\to\frac{\int_0^1v^{M-1}(1-v)^{M}dv}{\int_0^1v^{M}(1-v)^{M-1}dv}=1
 \qquad(t\to0^+),
\]
so $\Omega M_1-M_2<0$ for small $t$, and Lemma~\ref{R2:lem:bern} shows that
$W_\Omega$ is not CM.

\emph{Step 3: proof of (ii).} By Lemma~\ref{R2:lem:bern}, $-W_\Omega$ is CM
iff $M_2-\Omega M_1\ge0$, i.e.\ iff $\Omega\le\inf_{t>0}r(t)$ with
$r=M_2/M_1$; $r$ is continuous and $r<1$ (Step~2), and $r(0^+)=1$. As
$t\to\infty$, write $M_i(t)=t^2\int_0^1\alpha(tv)\gamma(t(1-v))\,\omega_i(v)\,dv$
with $\omega_1(v)=v$, $\omega_2(v)=1-v$. Since $c(y)/y\to1$, one has
$\alpha(u)/u^{M+1}\to\RtwoB(m+1,n+1)$ and $\gamma(u)/u^{M+1}\to\RtwoB(p+1,q+1)$
as $u\to\infty$; and the bound $\alpha(u)\le\RtwoB(m,n)u^{M-1}(1+u)^2$ (from
$c(y)\le1+y$), with its analogue for $\gamma$, shows that for $t\ge1$ the
functions $t^{-2M-2}\alpha(tv)\gamma(t(1-v))$ are dominated by
$16\,\RtwoB(m,n)\RtwoB(p,q)\,v^{M-1}(1-v)^{M-1}$. Dominated convergence gives
\[
 r(t)\to\frac{\int_0^1v^{M+1}(1-v)^{M+2}dv}{\int_0^1v^{M+2}(1-v)^{M+1}dv}=1 .
\]
Hence $\inf r$ is attained at an interior point, and $0<\Omega_*<1$.
\end{proof}

\begin{remark}\label{R2:rem:B-num}
Numerically (Gauss--Legendre quadrature of the densities, script
\texttt{R2\_\allowbreak{}verify\_\allowbreak{}remark18.py}): $\Omega_*(2,2;3,1)\approx0.97612$,
$\Omega_*(2,3;4,1)\approx0.97174$, $\Omega_*(2,4;5,1)\approx0.97207$,
$\Omega_*(3,3;4,2)\approx0.99206$ and $\Omega_*(3,3;5,1)\approx0.96453$, each
attained near $t\approx16$--$19$. We do not know a closed form, so the
second half of the $\Omega$-problem is answered only by this
characterization. Unlike Remark~17, we found no way to reduce
Theorem~\ref{R2:thm:B}(i) to known two-term inequalities: any two of the four
products in $W_1$ share two factors, and the pairwise splittings we tried,
with the sharp two-term thresholds of \cite{R2-QiRM2021}, cannot be
balanced.
\end{remark}

% ------------------------------------------------------------------
\subsection{Yin--Han--Qi's Remark~4: the function \texorpdfstring{$\Phi(x)=x\psi'(x)-1$}{Phi(x)=x psi'(x)-1}}
\label{R2:subsec:C}

\begin{problem}[{Yin--Han--Qi \cite[Remark~4]{R2-YHQ2024}}]\label{R2:prob:C}
With $\Phi(x)=x\psi'(x)-1$: ``We guess that, for $m,n\in\{0\}\cup\mathbb N$,
the function
\[
 \Phi^{(m+n+1)}(x)\big[\Phi^{(m)}(x)\Phi^{(n)}(x)\big]'-\Phi^{(m+n+2)}(x)\big[\Phi^{(m)}(x)\Phi^{(n)}(x)\big]
\]
is completely monotonic in $x\in(0,\infty)$. One can also consider necessary
and sufficient conditions on $\Lambda_{m,n}\in\mathbb R$ for
$m,n\in\{0\}\cup\mathbb N$ such that the function
\[
 \Phi^{(m+n+1)}(x)\left[\Phi^{(m)}(x)\Phi^{(n)}(x)\right]'-\Lambda_{m,n}\Phi^{(m+n+2)}(x)\left[\Phi^{(m)}(x)\Phi^{(n)}(x)\right]
\]
and its opposite are respectively completely monotonic on $(0,\infty)$.''
\end{problem}

\begin{theorem}\label{R2:thm:C}
Let $m,n\ge0$, $M=m+n$, $\Lambda\in\mathbb R$, and
\[
 G_\Lambda=\Phi^{(M+1)}[\Phi^{(m)}\Phi^{(n)}]'-\Lambda\Phi^{(M+2)}\Phi^{(m)}\Phi^{(n)}.
\]
Let $\alpha=(t^mb)*(t^nb)$, $\beta=t^{M+1}b$, and
\[
 N_1(t)=\int_0^t\alpha(u)\beta(t-u)\,u\,du,\qquad
 N_2(t)=\int_0^t\alpha(u)\beta(t-u)(t-u)\,du .
\]
\begin{enumerate}
\item[(i)] $G_\Lambda$ is CM on $(0,\infty)$ if and only if $\Lambda\le1$;
 in particular the guess ($\Lambda=1$) holds.
\item[(ii)] $-G_\Lambda$ is CM on $(0,\infty)$ if and only if
 $\Lambda\ge\Lambda^*_{m,n}:=\max_{t>0}N_1(t)/N_2(t)$; the maximum exists and
 $1<\Lambda^*_{m,n}<\infty$.
\end{enumerate}
\end{theorem}

\begin{proof}
By \eqref{R2:eq:Phirep}, $P:=(-1)^M\Phi^{(m)}\Phi^{(n)}=\RtwoL\alpha$ and
$B:=(-1)^{M+1}\Phi^{(M+1)}=\RtwoL\beta$, with $\Phi^{(M+2)}=(-1)^{M+1}B'$.
Hence $G_\Lambda=\Lambda B'P-BP'=\RtwoL[N_1-\Lambda N_2]$
(Lemma~\ref{R2:lem:wronski} with $(A,C)=(P,B)$).

(i) For $\Lambda=1$ we need $\alpha/\beta$ increasing. Here
\[
 \frac{\alpha(s)}{\beta(s)}=\int_0^1v^m(1-v)^n\,\frac{b(sv)\,b(s(1-v))}{b(s)}\,dv,
\]
and for each $v$ the logarithmic $s$-derivative of the integrand is
$v(\log b)'(sv)+(1-v)(\log b)'(s(1-v))-(\log b)'(s)>0$ by
Lemma~\ref{R2:lem:kernels}(c). So $N_1>N_2$ and $G_1$ is CM. For
$\Lambda<1$, $G_\Lambda=G_1+(1-\Lambda)(-B')P$ with $-B'=\RtwoL[t\beta]$ CM.
For $\Lambda>1$: $b$ is continuous on $[0,\infty)$ with $b(0)=\frac12$, so
$\alpha(u)\sim\frac14\RtwoB(m+1,n+1)u^{M+1}$ and $\beta(s)\sim\frac12s^{M+1}$
at $0$; thus $N_1/N_2\to1$ as $t\to0^+$, and $N_1-\Lambda N_2<0$ for small
$t$. Lemma~\ref{R2:lem:bern} applies.

(ii) $-G_\Lambda$ is CM iff $\Lambda N_2-N_1\ge0$, i.e.\ iff
$\Lambda\ge\sup N_1/N_2$. The ratio is continuous, $>1$ by (i), and tends to
$1$ as $t\to0^+$. It also tends to $1$ as $t\to\infty$: after $u=tv$, the
normalized densities $\alpha(tv)/(tv)^{M+1}\in[\frac14\RtwoB,\RtwoB]$ and
$\beta(t(1-v))/(t(1-v))^{M+1}=b(t(1-v))\in[\frac12,1)$ converge pointwise to
$\RtwoB=\RtwoB(m+1,n+1)$ and $1$ (Lemma~\ref{R2:lem:kernels}(c)), and
dominated convergence gives the limit
$\int_0^1v^{M+2}(1-v)^{M+1}dv/\int_0^1v^{M+1}(1-v)^{M+2}dv=1$. So the
supremum is a maximum and exceeds $1$.
\end{proof}

\begin{remark}\label{R2:rem:C}
(1) Numerically (script \texttt{R2\_\allowbreak{}verify\_\allowbreak{}yhq\_\allowbreak{}remark4.py}):
$\Lambda^*_{0,0}\approx1.11422$, $\Lambda^*_{0,1}\approx1.08382$,
$\Lambda^*_{1,1}\approx1.07538$, $\Lambda^*_{0,2}\approx1.06173$,
$\Lambda^*_{1,2}\approx1.06254$, $\Lambda^*_{2,2}\approx1.05612$ and
$\Lambda^*_{2,3}\approx1.04904$, attained near $t\approx12$--$16$. In
contrast with Theorem~\ref{R2:thm:A}(ii), the endpoint asymptotics of
$G_\Lambda$ cancel at $\Lambda=1$ at \emph{both} ends ($\Phi\sim x^{-1}$ at
$0$ and $\Phi\sim(2x)^{-1}$ at $\infty$), which is why $\Lambda^*$ is
determined by an interior maximum; no closed form is known to us.

(2) As in Remark~\ref{R2:rem:A-split}, part~(i) at $\Lambda=1$ can also be
deduced by splitting,
\begin{align*}
 G_1&=|\Phi^{(n)}|\Bigl(|\Phi^{(M+1)}\Phi^{(m+1)}|-\tfrac{m+1}{M+2}|\Phi^{(M+2)}\Phi^{(m)}|\Bigr)\\
 &\quad+|\Phi^{(m)}|\Bigl(|\Phi^{(M+1)}\Phi^{(n+1)}|-\tfrac{n+1}{M+2}|\Phi^{(M+2)}\Phi^{(n)}|\Bigr),
\end{align*}
from the $\Phi$-version of Qi's majorization theorem in the form stated in
\cite[Thm.~17]{R2-QiAgarwal2024} (``$-Y_{i,j;\ell,m;\Omega}$ is CM iff
$\Omega\ge\frac{i!j!}{\ell!m!}$'', applied to $(M+2,m)\succ(M+1,m+1)$). We
could not consult the original article cited there; the proof above is
independent of it.
\end{remark}

% ------------------------------------------------------------------
\subsection{Yin--Han--Qi's Remark~2: an elementary sign guess}
\label{R2:subsec:D}

\begin{problem}[{Yin--Han--Qi \cite[Lemma~3 and Remark~2]{R2-YHQ2024}}]\label{R2:prob:D}
For $0<x<y/2$ let
\[
 F(x,y)=2\Bigl(\frac1x-\frac1{y-x}\Bigr)+\frac12\Bigl(\frac{2^{y-x}}{y-x}-\frac{2^x}{x}\Bigr)
 -\frac{2^{y-x}-2^x}{(y-x)x}.
\]
The authors guess that $F(x,y)>0$ when $2<x<y/2$, and $F(x,y)<0$ when $y>4$
and $0<x<2$, and suggest discussing the sign of $F$ on the whole region
$0<x<y/2$.
\end{problem}

\begin{proposition}\label{R2:prop:D}
For $0<x<y/2$, with $E(s)=(2^s-1)/s$ ($E(0)=\log2$),
\[
 x(y-x)\,F(x,y)=2(x-2)(y-x-2)\bigl[E(y-x-2)-E(x-2)\bigr].
\]
Consequently
$\operatorname{sign}F(x,y)=\operatorname{sign}\bigl((x-2)(y-x-2)\bigr)$, and
$F$ vanishes exactly on the lines $x=2$ and $y=x+2$. Both guesses hold.
\end{proposition}
\begin{proof}
With $z=y-x$, $u=x-2$, $v=z-2$, clearing denominators gives
$xzF=(x-2)2^{z-1}-(z-2)2^{x-1}+2(z-x)=2[u(2^v-1)-v(2^u-1)]=2uv[E(v)-E(u)]$.
Since $E(s)=\log2\int_0^12^{sr}dr$ is strictly increasing and $v-u=z-x>0$,
the bracket is positive. If $2<x<y/2$ then $u>0$ and $v>u>0$; if
$0<x<2<y-2$ then $u<0<v$.
\end{proof}

% ------------------------------------------------------------------
\subsection{Completely monotonic degrees in the Qi--Agarwal survey}
\label{R2:subsec:EF}

Let
\begin{align*}
 Q_7(x)&=6x^7+51x^6+181x^5+345x^4+399x^3+283x^2+120x+24,\\
 Q_6(x)&=24x^6+156x^5+412x^4+621x^3+523x^2+242x+48,
\end{align*}
and
\[
 H_2=\psi'(x)-\frac{Q_7(x)}{6x^2(1+x)^4(2+x)^2},\qquad
 H_3=\frac{Q_6(x)}{24x^2(x+1)^4(x+2)}-\psi'(x),
\]
\[
 f_\lambda=H_2-\lambda g,\qquad g=\frac{4+8x+5x^2}{24x(1+x)^3(2+x)^2}.
\]
Then $f_\lambda$ is the function (5.2) of \cite{R2-QiAgarwal2019}. It is
known that $H_2,H_3$ are CM, that $f_\lambda$ is CM iff $\lambda\le0$, and
that $-f_\lambda$ is CM iff $\lambda\ge4$ (see
\cite[Sections~3.3, 5.4, 5.5]{R2-QiAgarwal2019} and the references there).

\begin{problem}[{Qi--Agarwal \cite[Sections~5.4 and 5.5]{R2-QiAgarwal2019}}]\label{R2:prob:EF}
``Can one compute completely monotonic degrees of the functions $H_2(x)$ and
$H_3(x)$ on $(0,\infty)$, as one did in [32]?'' and ``What are completely
monotonic degrees of the function (5.2) for $\lambda\le0$ and its negative
for $\lambda\ge4$ on $(0,\infty)$?''
\end{problem}

Partial fractions give the Laplace densities ($\psi'=\RtwoL c$):
\begin{align*}
 H_2&=\RtwoL p_{H_2},& p_{H_2}&=c(t)-t-e^{-t}\Bigl(\tfrac{43}6-\tfrac72t+t^2-\tfrac{t^3}{12}\Bigr)+e^{-2t}\Bigl(\tfrac{37}6+\tfrac{13}6t\Bigr),\\
 H_3&=\RtwoL p_3,& p_3&=t+\tfrac{13}{24}+e^{-t}\Bigl(\tfrac{55}{24}-\tfrac{15}8t+\tfrac{35}{48}t^2-\tfrac{t^3}{12}\Bigr)-\tfrac{11}6e^{-2t}-c(t),\\
 g&=\RtwoL p_g,& p_g&=\tfrac1{24}+e^{-t}\Bigl(-\tfrac38+\tfrac t8-\tfrac{t^2}{48}\Bigr)+e^{-2t}\Bigl(\tfrac13+\tfrac t6\Bigr),
\end{align*}
with $p_{H_2}(0)=p_3(0)=p_g(0)=0$. With $T=e^t$ one finds the exponential
polynomials
\begin{align*}
 \mathcal E_3&:=p_3'(t)(T-1)^2T^2\\
 &=T^3\bigl(\tfrac{t^3}{12}-\tfrac{47}{48}t^2+\tfrac{13}3t-\tfrac{31}6\bigr)
 +T^2\bigl(-\tfrac{t^3}6+\tfrac{47}{24}t^2-\tfrac{20}3t+13\bigr)\\
 &\quad+T\bigl(\tfrac{t^3}{12}-\tfrac{47}{48}t^2+\tfrac{10}3t-\tfrac{23}2\bigr)+\tfrac{11}3,\\
 \mathcal E_{2}&:=p_{H_2}(t)(T-1)^2T^2\\
 &=T^3\bigl(\tfrac{t^3}{12}-t^2+\tfrac92t-\tfrac{43}6\bigr)
 +T^2\bigl(-\tfrac{t^3}6+2t^2-\tfrac{35}6t+\tfrac{41}2\bigr)\\
 &\quad+T\bigl(\tfrac{t^3}{12}-t^2-\tfrac56t-\tfrac{39}2\bigr)+\tfrac{13}6t+\tfrac{37}6,\\
 \mathcal E_{2}^{\,\prime}&:=p_{H_2}'(t)(T-1)^2T^2\\
 &=T^3\bigl(-\tfrac{t^3}{12}+\tfrac54t^2-\tfrac{13}2t+\tfrac{35}3\bigr)
 +T^2\bigl(\tfrac{t^3}6-\tfrac52t^2+\tfrac{20}3t-\tfrac{65}2\bigr)\\
 &\quad+T\bigl(-\tfrac{t^3}{12}+\tfrac54t^2+\tfrac{19}6t+31\bigr)-\tfrac{13}3t-\tfrac{61}6,\\
 \mathcal E_4&:=(4p_g-p_{H_2})'(t)(T-1)^2T^2\\
 &=T^3\bigl(\tfrac{t^3}{12}-\tfrac76t^2+\tfrac{35}6t-\tfrac{29}3\bigr)
 +T^2\bigl(-\tfrac{t^3}6+\tfrac73t^2-\tfrac{20}3t+\tfrac{53}2\bigr)\\
 &\quad+T\bigl(\tfrac{t^3}{12}-\tfrac76t^2-\tfrac76t-25\bigr)+3t+\tfrac{49}6,\\
 p_g'(t)\,T^2&=T\bigl(\tfrac{t^2}{48}-\tfrac t6+\tfrac12\bigr)-\tfrac t3-\tfrac12 .
\end{align*}
(The notation $\mathcal E_2'$ is a name, not a $t$-derivative of $\mathcal E_2$.)

\begin{lemma}\label{R2:lem:taylor}
The Taylor coefficients $a_N=N!\,[t^N]$ of $\mathcal E_3$, $\mathcal E_2$,
$\mathcal E_4$ and of $p_g'T^2$ are all nonnegative; the first nonzero ones
are $a_4(\mathcal E_3)=\frac{45}2$, $a_5(\mathcal E_2)=\frac{50}3$,
$a_5(\mathcal E_4)=\frac{82}3$, $a_2(p_g'T^2)=\frac5{24}$. Hence $p_3'>0$,
$p_{H_2}>0$, $(4p_g-p_{H_2})'>0$ and $p_g'>0$ on $(0,\infty)$.
\end{lemma}
\begin{proof}
For $p_g'T^2$: $a_0=a_1=0$ and $a_N=(N^2-9N+24)/48>0$ for $N\ge2$ (negative
discriminant). For $\mathcal E=\sum_{k=0}^3Q_k(t)e^{kt}$ one has, for
$N\ge2$, $a_N=3^{N-3}u_3(N)+2^{N-3}u_2(N)+u_1(N)$ with
\begin{align*}
 \mathcal E_3:\ \ & u_3=\tfrac{4N^3-153N^2+2021N-6696}{48},\quad
 u_2=-\tfrac{2N^3-53N^2+371N-1248}{12},\\
 & u_1=\tfrac{4N^3-59N^2+215N-552}{48};\\
 \mathcal E_2:\ \ & u_3=\tfrac{N^3-39N^2+524N-2322}{12},\quad
 u_2=-\tfrac{N^3-27N^2+166N-984}{6},\\
 & u_1=\tfrac{N^3-15N^2+4N-234}{12};\\
 \mathcal E_4:\ \ & u_3=\tfrac{N^3-45N^2+674N-3132}{12},\quad
 u_2=-\tfrac{N^3-31N^2+190N-1272}{6},\\
 & u_1=\tfrac{N^3-17N^2+2N-300}{12}.
\end{align*}
For $N\ge90$ each $u_3(N)\ge N^3/275$ (the polynomial $275u_3(90+s)-(90+s)^3$
has nonnegative coefficients), while $|u_2(N)|\le7N^3$ and $|u_1(N)|\le N^3$
for $N\ge4$ (bound each coefficient using $N\ge4$). Hence
\[
 a_N\ge N^3\,2^{N-3}\bigl[(3/2)^{N-3}/275-8\bigr]>0\qquad(N\ge90).
\]
For $N<90$ the coefficients are evaluated in exact rational arithmetic (the
verification script checks all $N\le200$): all are $\ge0$, and the first
nonzero ones are as stated. Since $(T-1)^2T^2>0$, the
sign claims follow.
\end{proof}

\begin{lemma}[fractional powers]\label{R2:lem:frac}
Let $p\in C^1[0,\infty)$ with $p(0)=0$ and $|p'(t)|\le Ke^{-\kappa t}$
($K,\kappa>0$). For $0<\varepsilon<1$ put
$D^\varepsilon p(t)=\Gamma(1-\varepsilon)^{-1}\int_0^t(t-u)^{-\varepsilon}p'(u)du$.
Then $D^\varepsilon p$ is continuous and bounded, and
$x^\varepsilon\RtwoL p(x)=\RtwoL[D^\varepsilon p](x)$ for $x>0$. If moreover
$p>0$ on $(0,\infty)$, $p'>0$ on some $(0,\delta)$ and $p(t)\to L>0$, then
$D^\varepsilon p>0$ on $(0,\infty)$ for all sufficiently small
$\varepsilon>0$; hence $\Rtwodeg\RtwoL p>0$.
\end{lemma}
\begin{proof}
The identity is the convolution theorem:
$\RtwoL[D^\varepsilon p]=x^{\varepsilon-1}\RtwoL[p']=x^{\varepsilon-1}(x\RtwoL p-p(0))$.
Let $J_\varepsilon(t)=\int_0^t(1-u/t)^{-\varepsilon}p'(u)du$, which has the
sign of $D^\varepsilon p(t)$. For $t\le\delta$, $J_\varepsilon(t)>0$. For
$0<\varepsilon\le\frac12$, splitting at $t/2$ and using
$(1-y)^{-\varepsilon}-1\le2y$ on $[0,\frac12]$, we get
$|J_\varepsilon(t)-p(t)|\le\frac2t\int_0^\infty u|p'(u)|du+4tKe^{-\kappa t/2}$,
so there is $T$ with $J_\varepsilon\ge L/2$ on $[T,\infty)$ for all such
$\varepsilon$. On $[\delta,T]$,
$|J_\varepsilon(t)-p(t)|\le\max_{[0,T]}|p'|\,T\,\frac{\varepsilon}{1-\varepsilon}$,
which is smaller than $\min_{[\delta,T]}p>0$ for small $\varepsilon$.
Lemma~\ref{R2:lem:bern} concludes.
\end{proof}

\begin{theorem}\label{R2:thm:E}
$\Rtwodeg H_2=0$ and $\Rtwodeg H_3=1$.
\end{theorem}
\begin{proof}
By Lemma~\ref{R2:lem:taylor}, $p_{H_2}>0$, so $H_2$ is CM
(Lemma~\ref{R2:lem:bern}). From the partial fractions,
$Q_7/(6x^2(1+x)^4(2+x)^2)=x^{-2}+\frac{37}{24}+O(x)$ and
$\psi'(x)=x^{-2}+\zeta(2)+O(x)$, so
$H_2(0^+)=\zeta(2)-\frac{37}{24}\in(0,\infty)$. By
Lemma~\ref{R2:lem:degset}, $x^\alpha H_2$ is not CM for $\alpha>0$; thus
$\Rtwodeg H_2=0$.

Since $p_3(0)=0$ and $p_3'>0$ (Lemma~\ref{R2:lem:taylor}), $p_3>0$, so $H_3$
is CM, and $xH_3(x)=\RtwoL[p_3'](x)$ (integration by parts, $p_3$ bounded) is
CM; so $\Rtwodeg H_3\ge1$. As $x\to0^+$,
$Q_6/(24x^2(x+1)^4(x+2))=x^{-2}+\frac{13}{24}x^{-1}+O(1)$, so
$xH_3(x)\to\frac{13}{24}$ and $x^\alpha H_3\to0$ for $\alpha>1$;
Lemma~\ref{R2:lem:degset} gives $\Rtwodeg H_3=1$.
\end{proof}

\begin{theorem}\label{R2:thm:F}
\begin{enumerate}
\item[(a)] $\Rtwodeg f_0=0$.
\item[(b)] For every $\lambda<0$, $0<\Rtwodeg f_\lambda<1$; in particular
 the degree is not an integer. The map $\lambda\mapsto\Rtwodeg f_\lambda$ is
 nonincreasing on $(-\infty,0)$, tends to $1$ as $\lambda\to-\infty$ and to
 $0$ as $\lambda\to0^-$.
\item[(c)] For every $\mu\ge4$, $\Rtwodeg(-f_\mu)=1$.
\end{enumerate}
\end{theorem}
\begin{proof}
(a) is Theorem~\ref{R2:thm:E}. By Lemma~\ref{R2:lem:taylor}, $p_g$ increases
from $0$ to $\frac1{24}$, so $g$ and $xg(x)=\RtwoL[p_g'](x)$ are CM; since
$xg(x)\to\frac1{24}$ at $0^+$, $\Rtwodeg g=1$.

(b) Let $\lambda<0$. The density $p=p_{H_2}+|\lambda|p_g$ of $f_\lambda$ is
positive, so $f_\lambda$ is CM; $xf_\lambda(x)\to|\lambda|/24$ at $0^+$, so
$\Rtwodeg f_\lambda\le1$. The leading terms of $\mathcal E_2'$ and of
$p_g'T^2$ give, as $t\to\infty$,
\[
 p_{H_2}'(t)=-\tfrac{t^3}{12}e^{-t}\bigl(1+O(t^{-1})\bigr),\qquad
 p_g'(t)=\tfrac{t^2}{48}e^{-t}\bigl(1+O(t^{-1})\bigr).
\]
Hence $p'(t)<0$ for large $t$, $xf_\lambda=\RtwoL[p']$ is not CM
(Lemma~\ref{R2:lem:bern}), and $\Rtwodeg f_\lambda<1$ by
Lemma~\ref{R2:lem:degset}. The first nonzero Taylor coefficients give
$p_{H_2}=\frac5{36}t^3+O(t^4)$ and $p_g=\frac5{144}t^3+O(t^4)$, so $p'>0$ near
$0$; moreover $p\to|\lambda|/24>0$ and $p'$ decays exponentially.
Lemma~\ref{R2:lem:frac} yields $\Rtwodeg f_\lambda>0$.

Monotonicity: for $\lambda_1<\lambda_2<0$ and $\alpha\le1$,
$x^\alpha f_{\lambda_1}=x^\alpha f_{\lambda_2}+(\lambda_2-\lambda_1)x^\alpha g$,
and $x^\alpha g$ is CM. Limits: fix $\alpha\in(0,1)$. Since
$\int_0^\infty p_{H_2}'=0$, splitting the Caputo integral at $t/2$ gives
$|D^\alpha p_{H_2}(t)|=O(t^{-\alpha-1})$ as $t\to\infty$, while
$D^\alpha p_g(t)\ge t^{-\alpha}\Gamma(1-\alpha)^{-1}\int_0^tp_g'\ge c_\alpha t^{-\alpha}$
for large $t$; near $0$ both are $\asymp t^{3-\alpha}$ with positive
constants. Hence $S_\alpha=\sup_t|D^\alpha p_{H_2}|/D^\alpha p_g<\infty$,
and for $|\lambda|>S_\alpha$ the density
$D^\alpha p=D^\alpha p_{H_2}+|\lambda|D^\alpha p_g$ of $x^\alpha f_\lambda$
(Lemma~\ref{R2:lem:frac}) is positive, so $\Rtwodeg f_\lambda\ge\alpha$.
Conversely, $x^\alpha H_2=\RtwoL[D^\alpha p_{H_2}]$ is not CM by (a), so
$D^\alpha p_{H_2}(t_\alpha)<0$ for some $t_\alpha$; then
$D^\alpha p(t_\alpha)<0$ whenever
$|\lambda|<-D^\alpha p_{H_2}(t_\alpha)/D^\alpha p_g(t_\alpha)$, and
$\Rtwodeg f_\lambda<\alpha$.

(c) The density of $-f_\mu$ is $q_\mu=\mu p_g-p_{H_2}$, and
$q_\mu'=(4p_g-p_{H_2})'+(\mu-4)p_g'>0$ by Lemma~\ref{R2:lem:taylor}. With
$q_\mu(0)=0$ this gives $q_\mu>0$, so $-f_\mu$ and
$-xf_\mu=\RtwoL[q_\mu']$ are CM. Since $-xf_\mu(x)\to\mu/24>0$ at $0^+$,
Lemma~\ref{R2:lem:degset} gives $\Rtwodeg(-f_\mu)=1$.
\end{proof}

\begin{remark}\label{R2:rem:EF}
(1) Theorem~\ref{R2:thm:F}(b) is the sign-change mechanism of method~M9 in a
different guise: for $\lambda<0$ the obstruction to degree $1$ is not a
vanishing moment but the exponentially small negative tail
$-\frac{t^3}{12}e^{-t}$ of $p_{H_2}'$, which beats $|\lambda|\frac{t^2}{48}e^{-t}$
for every fixed $\lambda$; numerically the first sign change of $p'$ lies in
$(2,2.25)$, $(9,9.25)$ and $(32,32.25)$ for $\lambda=-1,-10,-100$.
(2) The related conjecture of \cite[Remark~7.4]{R2-QiAIMS2020} (degree $4$ for
$h_\lambda$ and $-h_\mu$) is false for the trivial reason that
$h_\lambda(x)\sim-\lambda/(12x^3)$ at $0^+$; this has already been recorded
publicly \cite{R2-Pudim} and is not claimed here.
\end{remark}

% ------------------------------------------------------------------
\subsection{A partial answer for Qi's Remark~4.1 (MIA 2021)}
\label{R2:subsec:G}

Let $\mathcal G(x)=x(x\psi'(x)-1)-\frac12=\RtwoL w$ with
$w(t)=e^t[(t-2)e^t+t+2]/(e^t-1)^3$, $w(0)=\frac16$ (note $w=f''$ with
$f=t/(e^t-1)$). For $(i,j)\succ(p,q)$ (same sum, $\max(i,j)>\max(p,q)$,
$i,j,p,q\ge0$) Qi \cite[Thm.~3.1]{R2-QiMIA2021} proved that
$\mathcal G_\Lambda=(-1)^{i+j}\mathcal G^{(i)}\mathcal G^{(j)}-(-1)^{p+q}\Lambda\mathcal G^{(p)}\mathcal G^{(q)}$
is CM if $\Lambda\le1$, and that
$1<\mathcal G^{(i)}\mathcal G^{(j)}/(\mathcal G^{(p)}\mathcal G^{(q)})<\frac{i!j!}{p!q!}$.
\cite[Remark~4.1]{R2-QiMIA2021} asks: ``What is the necessary and sufficient
condition on $\Lambda_{i,j;p,q}$ such that'' $\mathcal G_\Lambda$ is CM, and
``What is the sharp lower bound of the left hand side inequality in (3.1)?''

\begin{proposition}\label{R2:prop:G}
Let $R(t)=\int_0^1\sigma_{i+1,j+1}(v)K_t(v)dv\big/\int_0^1\sigma_{p+1,q+1}(v)K_t(v)dv$
with $K_t(v)=w(tv)w(t(1-v))$. Then $\mathcal G_\Lambda$ is CM iff
$\Lambda\le\Lambda_*:=\inf_{t>0}R(t)$, and
\[
 1<\Lambda_*\le\frac{(i+1)!\,(j+1)!}{(p+1)!\,(q+1)!}<\frac{i!\,j!}{p!\,q!}.
\]
In particular Qi's sufficient condition $\Lambda\le1$ is never sharp.
\end{proposition}
\begin{proof}
The density of $\mathcal G_\Lambda$ is
$t^{i+j+1}\int_0^1[v^i(1-v)^j-\Lambda v^p(1-v)^q]K_t(v)dv$, so
Lemma~\ref{R2:lem:bern} gives the characterization. The pointwise inequality
$\sigma_{i+1,j+1}>\sigma_{p+1,q+1}$ on $(0,\frac12)$ (Lemma~\ref{R2:lem:cosh})
gives $R>1$; $R(0^+)=\frac{i!j!}{p!q!}$, and since $w(y)=(y-2)e^{-y}(1+o(1))$
and $0<w(y)e^{y}\le C(1+y)$, dominated convergence gives
$R(\infty)=\RtwoB(i+2,j+2)/\RtwoB(p+2,q+2)=\frac{(i+1)!(j+1)!}{(p+1)!(q+1)!}>1$.
Hence $\inf R>1$ and $\inf R\le R(\infty)$. The last inequality of the
display is $(i+1)(j+1)<(p+1)(q+1)$, i.e.\ $ij<pq$.
\end{proof}

\noindent Numerically $R$ is \emph{not} monotone: it decreases to an interior
minimum near $t\approx19$--$25$ and then increases to $R(\infty)$ (script
\texttt{R2\_\allowbreak{}verify\_\allowbreak{}mia\_\allowbreak{}partial.py}); e.g.\
$\Lambda_*(2,0;1,1)\approx1.4078<\frac32$, $\Lambda_*(3,0;2,1)\approx1.8157<2$
and $\Lambda_*(3,1;2,2)\approx1.2820<\frac43$. So the sharp constant is not
$R(\infty)$, and no closed form is offered. For the second question the
numerics support the answer
\[
 \inf_{x>0}\frac{\mathcal G^{(i)}(x)\mathcal G^{(j)}(x)}{\mathcal G^{(p)}(x)\mathcal G^{(q)}(x)}
 =\frac{\mu_i\mu_j}{\mu_p\mu_q}\qquad(\text{the value at }x=0^+),
\]
where $\mu_0=\frac12$,
$\mu_1=1$ and $\mu_k=k!(k-1)\zeta(k)$ for $k\ge2$ (e.g.\ $\pi^2/6$ for
$(2,0;1,1)$). \emph{But} the Wronskian density
$\int_0^t\alpha(u)\gamma(t-u)(t-2u)du$ changes sign (near $t\approx50$--$60$ for
$(2,0;1,1)$), so our method does not prove monotonicity of the ratio. This
question remains open.

% ------------------------------------------------------------------
\subsection{Verification and limits}
\label{R2:subsec:verif}

{\sloppy
Independent checks (scripts in the area directory; each prints PASS/FAIL):
\texttt{R2\_\allowbreak{}verify\_\allowbreak{}kernels.py} (the symbolic identities and the coefficient
argument of Lemma~\ref{R2:lem:kernels}); \texttt{R2\_\allowbreak{}verify\_\allowbreak{}remark17.py} (for
seven pairs $(m,n)$: signs of $(-1)^jH_1^{(j)}$ and
$(-1)^j(-H_{\Omega^\sharp})^{(j)}$ for $j\le8$ at eight points in $[0.02,60]$,
with exact polygamma derivatives at 60 digits; failure of CM just outside
both thresholds; the Hurwitz form of $-H_{\Omega^\sharp}$; the splitting
identity of Remark~\ref{R2:rem:A-split}; positivity of the density on
$[0.05,120]$); \texttt{R2\_\allowbreak{}verify\_\allowbreak{}remark18.py} and
\texttt{R2\_\allowbreak{}verify\_\allowbreak{}yhq\_\allowbreak{}remark4.py} (the same checks for $W_1$ and $G_1$, the
monotonicity of $h$ in Lemma~\ref{R2:lem:cosh}, and the constants
$\Omega_*$, $\Lambda^*$); \texttt{R2\_\allowbreak{}verify\_\allowbreak{}yhq\_\allowbreak{}remark2.py} (the identity of
Proposition~\ref{R2:prop:D} and 20000 random sign tests);
\texttt{R2\_\allowbreak{}verify\_\allowbreak{}qa2019\_\allowbreak{}degrees.py} (partial fractions, a numerical Laplace
check of $p_{H_2},p_3,p_g$, the exponential-polynomial identities, the exact
Taylor certificates of Lemma~\ref{R2:lem:taylor} with the tail bound, the
endpoint asymptotics); \texttt{R2\_\allowbreak{}verify\_\allowbreak{}mia\_\allowbreak{}partial.py}
(Proposition~\ref{R2:prop:G} and the numerics after it). The finite
derivative checks are consistency tests only; the proofs do not rely on
them, except for the exact finite coefficient check in
Lemma~\ref{R2:lem:taylor}.\par}

What remains open here: closed forms (if any) for $\Omega_*$ and
$\Lambda^*$ in Theorems~\ref{R2:thm:B}(ii) and~\ref{R2:thm:C}(ii); the exact
values of $\Rtwodeg f_\lambda$ for $\lambda<0$; both questions of
\cite[Remark~4.1]{R2-QiMIA2021} beyond Proposition~\ref{R2:prop:G}; and the
two-variable curvature conjecture of Le~Brigant--Preston--Puechmorel recalled
in \cite{R2-QiAgarwal2024,R2-Qi2026}, which we did not attempt.


%% ==================================================================
%% Section: Composite sequences
%% ==================================================================
% R6 section: composite integer parts of powers, composite linear recurrences
\providecommand{\Z}{\mathbb{Z}}
\providecommand{\Q}{\mathbb{Q}}
\providecommand{\R}{\mathbb{R}}
\providecommand{\C}{\mathbb{C}}
\providecommand{\F}{\mathbb{F}}
\providecommand{\N}{\mathbb{N}}
\providecommand{\OO}{\mathcal{O}}
\providecommand{\Tr}{\operatorname{Tr}}
\providecommand{\Nm}{\operatorname{N}}
\providecommand{\fl}[1]{\left\lfloor #1\right\rfloor}
\providecommand{\cl}[1]{\left\lceil #1\right\rceil}
\providecommand{\Rsixord}{\operatorname{ord}}

\section{Composite integer parts of powers and composite linear recurrences}
\label{R6:sec}

This section treats two problems posed by other authors in the area of method~M10
(trace congruences of powers in residue rings and return of finite-field powers).
\begin{itemize}
\item \S\ref{R6:subsec:biquad}: the conjecture of Hahn, Ismailescu, Kim and Kim that there are
infinitely many biquadratic Pisot numbers all of whose integer parts of powers are $\equiv 3,4\pmod 6$
\cite[\S5]{R6-HIKK}. The literal statement is easy: powers of a single example with small
conjugates of total size $<1$, such as the authors' own $z^4-64z^3+8z+1$, already give infinitely
many. The content of Theorem~\ref{R6:thm:biquad} is a field-by-field version with infinitely many
multiplicatively independent examples in every real biquadratic field, valid for all $n\ge1$.
\item \S\ref{R6:subsec:siurys}: \v{S}iurys's conjecture that for every $k>2$ some sequence obeying the $k$-step
Fibonacci recurrence has coprime positive initial values and no prime terms
\cite[Conjecture~1]{R6-Siurys}. We obtain a \emph{partial} result
(Theorem~\ref{R6:thm:siurys}): the conjecture holds for all $k$ in nine residue classes
modulo $24$ and two residue classes modulo $120$ (the latter inside the hardest class
$k\equiv11\pmod{12}$), for all $k\not\equiv 0,3,11\pmod{12}$ outside an explicit set of density zero,
and, by computer, for every $3\le k\le 130$. The general conjecture remains open.
\end{itemize}
\S\ref{R6:subsec:other} records the status of further targets that were examined.

%=====================================================================================
\subsection{Biquadratic Pisot numbers with composite integer parts of powers}
\label{R6:subsec:biquad}

A \emph{Pisot number} is a real algebraic integer $\alpha>1$ whose other conjugates lie in the open
unit disc. Following \cite{R6-HIKK}, a \emph{biquadratic Pisot number} is a Pisot number $\alpha$
with $\Q(\alpha)$ a biquadratic field $\Q(\sqrt p,\sqrt q)$ ($p,q>1$ squarefree; in
\cite{R6-HIKK} also coprime). Such fields are totally real of degree $4$.

\begin{problem}[Hahn--Ismailescu--Kim--Kim {\cite[\S5, after Table~1]{R6-HIKK}}]
\label{R6:prob:HIKK}
After listing in Table~1 (caption: ``Biquadratic Pisot numbers $\alpha<100$ with the property that
$\lfloor\alpha^n\rfloor\equiv 3,4$ for all $n\ge3$'') the examples they found, with the remark
``For all these entries $\lfloor\alpha^n\rfloor\equiv 3,4\pmod 6$ as soon as $n\ge3$'', the authors write:
``We conjecture that there are infinitely many biquadratic Pisot numbers with the property described above.''
\end{problem}

The tool in \cite{R6-HIKK} is their Theorem~6: if the minimal polynomial
$z^4-az^3+bz^2-cz+d$ of a biquadratic Pisot number satisfies $a\equiv c\equiv 4\pmod 6$,
$b\equiv0\pmod 3$, $d\equiv 1\pmod 3$, then $\fl{\alpha^n}\equiv 3,4\pmod 6$ for all
sufficiently large $n$. The difficulty addressed below is to produce infinitely many genuinely
different such numbers, with control for \emph{every} $n$.

\begin{theorem}\label{R6:thm:biquad}
Let $K$ be a real biquadratic field. There is an infinite sequence $\alpha_1,\alpha_2,\dots$ of
Pisot numbers in $K$, no two of which are multiplicatively dependent, such that for every $j$:
\begin{enumerate}
\item[(i)] $\Q(\alpha_j)=K$;
\item[(ii)] the minimal polynomial $P_j$ of $\alpha_j$ satisfies $P_j(z)\equiv (z-4)^4\pmod 6$
coefficientwise; in particular $P_j(z)=z^4-az^3+bz^2-cz+d$ with $a\equiv c\equiv d\equiv 4$ and
$b\equiv 0\pmod 6$, so the hypotheses of \cite[Theorem~6]{R6-HIKK} hold;
\item[(iii)] for every $n\ge 1$, $\fl{\alpha_j^n}\equiv 3$ or $4\pmod 6$ and $\fl{\alpha_j^n}$ is
composite.
\end{enumerate}
\end{theorem}

\begin{corollary}\label{R6:cor:HIKK}
The conjecture in Problem~\ref{R6:prob:HIKK} is true. More precisely, there are infinitely many
biquadratic Pisot numbers $\alpha$ with $\fl{\alpha^n}\equiv 3,4\pmod 6$ for \emph{all} $n\ge1$;
they can be taken in each of infinitely many fields $\Q(\sqrt p,\sqrt q)$ with $p,q$ coprime
squarefree (e.g.\ $\Q(\sqrt 2,\sqrt q)$, $q$ an odd prime), infinitely many in each such field, and
pairwise multiplicatively independent (so none is a power of another). As noted in
Remark~\ref{R6:rem:explicit}(b), the literal conjecture already follows from the powers of a single
example; the new content is the independence and the statement for all $n\ge1$ in every field.
\end{corollary}

We use two lemmas; the first is a standard consequence of Dirichlet's unit theorem.

\begin{lemma}[Pisot units]\label{R6:lem:pisotunit}
Let $K\neq\Q$ be a number field with a real embedding $\sigma_1$. There is a unit $u\in\OO_K^\times$
with $\sigma_1(u)>1$ and $|\sigma(u)|<1$ for every embedding $\sigma\ne\sigma_1$ of $K$ into $\C$.
Every such $u$ satisfies $\Q(u)=K$, and so does every power $u^e$, $e\ge1$.
\end{lemma}
\begin{proof}
Let $v_1,\dots,v_s$ be the archimedean places ($v_1$ induced by $\sigma_1$), $n_v\in\{1,2\}$ the local
degrees, and $\ell(x)=(\log|x|_{v})_v\in\R^s$. By Dirichlet's theorem $\ell(\OO_K^\times)$ is a lattice
of full rank $s-1\ge1$ in $H=\{y:\sum_v n_vy_v=0\}$. The set
$C=\{y\in H: y_{v_1}>0,\ y_v<0\ (v\ne v_1)\}$ is a nonempty open cone in $H$ (it contains
$y_{v_1}=\sum_{v\ne v_1}n_v$, $y_v=-1$), hence contains a ball $B(c,r)$ and therefore all balls
$B(tc,tr)$, $t>0$. For $t$ larger than the covering radius of the lattice divided by $r$, $B(tc,tr)$
contains some $\ell(u)$. Replacing $u$ by $u^2$ makes $\sigma_1(u)>0$. If $\Q(u^e)=F\subsetneq K$,
then $[K:F]\ge2$ embeddings of $K$ restrict to $\sigma_1|_F$, so some $\sigma\ne\sigma_1$ has
$\sigma(u^e)=\sigma_1(u^e)$, contradicting $|\sigma(u)|<1<\sigma_1(u)$.
\end{proof}

\begin{lemma}[Trace congruence]\label{R6:lem:trace}
Let $K$ be a number field of degree $D$, $c\in\Z$, $M\ge 2$, and $\alpha\in\OO_K$ with
$\alpha\equiv c\pmod{M\OO_K}$. Then the characteristic polynomial of $\alpha$ is
$\equiv(z-c)^D\pmod M$ coefficientwise, and $\Tr_{K/\Q}(\alpha^n)\equiv Dc^n\pmod M$ for all
$n\ge0$.
\end{lemma}
\begin{proof}
Write $\alpha=c+M\beta$, $\beta\in\OO_K$. With $\beta_1,\dots,\beta_D$ the conjugates of $\beta$,
the characteristic polynomial of $\alpha$ is $\prod_i(z-c-M\beta_i)=Q(z-c)$, where
$Q(w)=\prod_i(w-M\beta_i)=w^D+\sum_{t\ge1}(-M)^te_t(\beta)w^{D-t}$ and the $e_t(\beta)\in\Z$; all
non-leading coefficients of $Q$ are divisible by $M$, hence so are those of $Q(z-c)-(z-c)^D$.
Also $\alpha^n-c^n\in M\OO_K$ and $\Tr(M\OO_K)\subseteq M\Z$.
\end{proof}

\begin{proof}[Proof of Theorem~\ref{R6:thm:biquad}]
Fix a real embedding $\sigma_1$ of $K$ (all four embeddings are real) and take $u$ as in
Lemma~\ref{R6:lem:pisotunit}. Since $u$ is a unit, its class in the finite group
$(\OO_K/6\OO_K)^\times$ has some order $m\ge1$. For $j\ge1$ put $\alpha_j=4u^{jm}$, so
$\alpha_j\equiv 4\pmod{6\OO_K}$. Let $\rho=\max_{\sigma\neq\sigma_1}|\sigma(u)|<1$ and choose $j_0$ with
$12\rho^{j_0m}<1$; we use $j\ge j_0$ and relabel.

\emph{(i) and the Pisot property.} $\Q(\alpha_j)=\Q(u^{jm})=K$ by Lemma~\ref{R6:lem:pisotunit}, so the
conjugates of $\alpha_j$ are $\sigma(\alpha_j)$, $\sigma$ running over the four embeddings. Now
$\sigma_1(\alpha_j)>4$ and
\begin{equation}\label{R6:eq:smallsum}
S_j:=\sum_{\sigma\ne\sigma_1}|\sigma(\alpha_j)|\le 12\rho^{jm}<1 .
\end{equation}
Thus $\alpha_j=\sigma_1(\alpha_j)$ is a Pisot number of degree $4$ in $K$.

\emph{(ii)} is Lemma~\ref{R6:lem:trace} with $c=4$, $M=6$, $D=4$; $(z-4)^4=z^4-16z^3+96z^2-256z+256$
and $16\equiv 256\equiv 4$, $96\equiv0\pmod 6$.

\emph{(iii)} Let $T_n=\Tr(\alpha_j^n)\in\Z$ and $E_n=\sum_{\sigma\ne\sigma_1}\sigma(\alpha_j)^n\in\R$, so
$\alpha_j^n=T_n-E_n$. By Lemma~\ref{R6:lem:trace}, $T_n\equiv 4\cdot4^n\equiv 4\pmod 6$ for $n\ge1$
(as $4^t\equiv4\pmod 6$ for $t\ge1$). By \eqref{R6:eq:smallsum},
$|E_n|\le\sum_{\sigma\ne\sigma_1}|\sigma(\alpha_j)|^n\le S_j^n<1$, so
$\fl{\alpha_j^n}=T_n-\cl{E_n}$ with $\cl{E_n}\in\{0,1\}$. Hence
$\fl{\alpha_j^n}\equiv 4$ or $3\pmod 6$. Since $\fl{\alpha_j^n}\ge\fl{\alpha_j}\ge4$, it is divisible
by $2$ (residue $4$) or by $3$ (residue $3$) and exceeds that prime; so it is composite.

\emph{Independence.} If $\alpha_j^x=\alpha_{j'}^y$ with integers $(x,y)\ne(0,0)$, taking norms gives
$4^{4x}=4^{4y}$, so $x=y\neq0$, and then $u^{mx(j-j')}=1$; as $u$ is not a root of unity, $j=j'$.
\end{proof}

\begin{proof}[Proof of Corollary~\ref{R6:cor:HIKK}]
Apply Theorem~\ref{R6:thm:biquad} to $K=\Q(\sqrt2,\sqrt q)$ for each odd prime $q$; these fields are
pairwise distinct (their quadratic subfields differ) and $(2,q)$ are coprime squarefree. Each
$\alpha_j$ has $\fl{\alpha_j^n}\equiv3,4\pmod 6$ for all $n\ge1$, a fortiori for $n\ge3$.
\end{proof}

\begin{remark}[Examples and the literal conjecture]\label{R6:rem:explicit}
(a) Let $K=\Q(\sqrt2,\sqrt5)$ and $u=(1+\sqrt2)(2+\sqrt5)(3+\sqrt{10})$, so $u\approx63.02$; the
other conjugates of $u$ have absolute value $\approx0.092,\ 0.285,\ 0.603$ (a product of quadratic units is a
Pisot unit as soon as the three logarithms $x\log\varepsilon_p,\ y\log\varepsilon_q,\ z\log\varepsilon_{pq}$
of the factors satisfy the strict triangle inequalities), its order modulo $6\OO_K$ is $m=4$, and
already $\alpha=4u^4$ has $S=0.554<1$ and minimal polynomial
\[
z^4-63092752\,z^3+34947168\,z^2-884992\,z+256\equiv(z-4)^4\pmod 6 .
\]
Hence every $\alpha_j=4u^{4j}$ ($j\ge1$) has $\fl{\alpha_j^n}\equiv 3,4\pmod 6$ for all $n\ge1$.
The script \texttt{verify\_\allowbreak{}biquadratic.py} recomputes these data exactly and the same construction in
the fields $\Q(\sqrt p,\sqrt q)$ for $(p,q)=(2,3)$, $(2,5)$, $(3,5)$, $(2,7)$, $(5,13)$, $(6,10)$,
$(3,7)$, $(11,13)$ (the field $\Q(\sqrt6,\sqrt{10})$ has no coprime pair of radicands), checking minimal polynomials,
the congruence mod $6$, $S<1$, and $\fl{\alpha^n}\bmod 6$ for $1\le n\le 120$.

(b) The literal conjecture also follows from a single example: for the largest root $\alpha_\ast$ of
$z^4-64z^3+8z+1$ (Table~1 of \cite{R6-HIKK}) the three small conjugates have
$\sum|\cdot|\approx0.809<1$, and the polynomial is $\equiv(z+1)^4\pmod 2$ and $\equiv(z-1)^4\pmod 3$, so
all reduced roots equal $1$, $\Tr(\alpha_\ast^n)\equiv4\pmod 6$, and the argument above gives
$\fl{\alpha_\ast^n}\equiv3,4\pmod 6$ for all $n\ge1$; the powers $\alpha_\ast^e$ ($e\ge1$) are then infinitely many biquadratic Pisot numbers with
the property. Theorem~\ref{R6:thm:biquad} is the non-degenerate form: infinitely many multiplicatively
independent examples in every real biquadratic field. We also checked all ten entries of Table~1 of
\cite{R6-HIKK} that are given by explicit polynomials (Galois group $V_4$, Pisot, congruences, and
$\fl{\alpha^n}\bmod 6\in\{3,4\}$ for $3\le n\le120$; for three of them $n=2$, and for one of these
also $n=1$, fall outside $\{3,4\}$, consistent with the restriction $n\ge3$ in \cite{R6-HIKK}).

(c) The proof gives more: if $K\ne\Q$ is any number field with a real embedding and
$[K:\Q]\not\equiv 0,2\pmod 6$, then $K$ contains infinitely many multiplicatively independent Pisot
numbers $\alpha$ with $\Q(\alpha)=K$ and all $\fl{\alpha^n}$ ($n\ge1$) composite: take
$\alpha=c\,u^{jm}$ with $c=4$ if $[K:\Q]\equiv1\pmod 3$ (then $\Tr(\alpha^n)\equiv 4$) and $c=3$ if
$[K:\Q]$ is odd (then $\Tr(\alpha^n)\equiv3\pmod 6$, and $\fl{\alpha^n}\equiv 3$ or $2\pmod 6$),
with $j\ge j_0$ large enough that the conjugates other than $\alpha$ have total size $<1$ and
$\fl{\alpha^n}>3$.
\end{remark}

%=====================================================================================
\subsection{\texorpdfstring{$k$}{k}-step Fibonacci-like sequences without primes}
\label{R6:subsec:siurys}

For $k\ge2$ and integers $a_0,\dots,a_{k-1}$ let $S_k(a_0,\dots,a_{k-1})=(x_n)_{n\ge0}$ be given by
$x_n=a_n$ ($n<k$) and $x_n=x_{n-1}+\dots+x_{n-k}$ ($n\ge k$). Graham \cite{R6-Graham} found coprime
positive $a_0,a_1$ with $S_2(a_0,a_1)$ free of primes; \v{S}iurys \cite{R6-Siurys} proved the analogue for
$2\le k\le10$ and for $k\equiv79\pmod{120}$, and stated:

\begin{problem}[\v{S}iurys {\cite[Conjecture~1]{R6-Siurys}}; also {\cite{R6-SiurysThesis}}]
\label{R6:prob:siurys}
``Let $k>2$ be some fixed positive integer. Then there exist positive integers
$a_0,a_1,\dots,a_{k-1}$ such that $\gcd(a_0,a_1,\dots,a_{k-1})=1$ and the sequence
$S_k(a_0,a_1,\dots,a_{k-1})$ contains no prime numbers.''
\end{problem}

To our knowledge (searches in October 2026) the only later progress is Lunev's smaller initial
values for $k=3$ \cite{R6-Lunev}. Throughout, $k\ge3$,
\[
f_k(X)=X^k-X^{k-1}-\dots-X-1,\quad g_k(X)=(X-1)f_k(X)=X^{k+1}-2X^k+1 .
\]
For a prime $p$, a \emph{solution mod $p$} is a sequence $(y_n)_{n\ge0}$ in $\F_p$ with
$y_n=y_{n-1}+\dots+y_{n-k}$ for $n\ge k$; it is \emph{nonzero} if $(y_0,\dots,y_{k-1})\ne0$. We say
that it \emph{vanishes on} $c\bmod d$ if $y_n=0$ for all $n\ge0$ with $n\equiv c\pmod d$.

\begin{lemma}[Assembly]\label{R6:lem:assembly}
Let $p_1,\dots,p_r$ be distinct primes and $y^{(i)}$ a nonzero solution mod $p_i$ ($1\le i\le r$)
such that every $n\ge0$ has $y^{(i)}_n=0$ for some $i$. Then there are positive integers
$a_0,\dots,a_{k-1}$ with $\gcd=1$ such that every term of $S_k(a_0,\dots,a_{k-1})$ is composite.
\end{lemma}
\begin{proof}
Let $\Pi=p_1\cdots p_r$ and, by the Chinese remainder theorem, $b_j\in[1,\Pi]$ with
$b_j\equiv y^{(i)}_j\pmod{p_i}$ for all $i,j$. Put $a_j=b_j+\Pi$ for $1\le j\le k-1$ and
$a_0=b_0+\Pi t$ with $t\ge1$ chosen below. By induction on $n$, $x_n\equiv y^{(i)}_n\pmod{p_i}$ for all
$n,i$, so every $x_n$ has some $p_i$ as a divisor; moreover $x_n\ge\min_ja_j>\Pi\ge p_i$ (for $n\ge k$
the term $x_n$ is a sum of $k$ earlier positive terms), so $x_n$ is composite. If a prime $q$ divides
all $a_j$, then $q\ne p_i$ (since some $y^{(i)}_j\ne0$) and $q\mid G:=\gcd(a_1,\dots,a_{k-1})$, a fixed
positive integer. Choosing $t$ (Chinese remainder theorem) with $b_0+\Pi t\not\equiv0\pmod q$ for the
finitely many primes $q\mid G$, $q\nmid\Pi$, gives $\gcd(a_0,\dots,a_{k-1})=1$.
\end{proof}

\begin{lemma}[The prime $2$]\label{R6:lem:two}
For $0\le a<b\le k$ the sequence with $y_n=1$ if $n\equiv a$ or $b\pmod{k+1}$ and $y_n=0$
otherwise is a nonzero solution mod $2$.
\end{lemma}
\begin{proof}
For $n\ge k$, $y_n-\sum_{i=1}^ky_{n-i}\equiv\sum_{i=0}^{k}y_{n-i}\pmod 2$, and any $k+1$ consecutive
terms of this $(k+1)$-periodic sequence contain exactly two $1$'s. Also $y_a=1$ with $a<k$.
\end{proof}

Thus the prime $2$ alone annihilates all residues mod $k+1$ except two freely chosen ones. (In fact the
solutions mod $2$ are exactly the $(k+1)$-periodic sequences of even weight, because
$g_k\equiv X^{k+1}+1\pmod 2$.) Odd primes are used through pairs of roots of $g_k$ with
root-of-unity ratio:

\begin{lemma}[Pair patterns]\label{R6:lem:pair}
Let $p$ be an odd prime, $d\ge3$ with $p\nmid d$, $\zeta\in\overline{\F}_p$ of multiplicative order
$d$, and $\theta\in\overline{\F}_p$ with $\theta\ne1$, $\zeta\theta\ne1$ and
$g_k(\theta)=g_k(\zeta\theta)=0$. Then for every $c\in\Z$ there is a nonzero solution mod $p$
vanishing on $c\bmod d$.
\end{lemma}
\begin{proof}
Let $F=\F_p(\zeta,\theta)$. Every Galois conjugate of $\theta$ or $\zeta\theta$ is a root of $g_k$
different from $1$, hence a root of $f_k$, so $n\mapsto\rho^n$ is a solution (over $F$) for each such
$\rho$. For $\lambda\in F$ put
$y_n(\lambda)=\Tr_{F/\F_p}\bigl(\lambda\,\theta^n(1-\zeta^{\,n-c})\bigr)
=\Tr(\lambda\theta^n)-\Tr(\lambda\zeta^{-c}(\zeta\theta)^n)$, a solution mod $p$ with values in
$\F_p$. If $n\equiv c\pmod d$ then $\zeta^{n-c}=1$ and $y_n(\lambda)=0$. Finally, taking the representative $0\le c<d$, the element
$\theta^{c+1}(1-\zeta)$ is nonzero and the trace form of $F/\F_p$ is nondegenerate, so
$y_{c+1}(\lambda)\neq0$ for some $\lambda$; this solution is nonzero.
\end{proof}

For $\zeta\in\overline{\Q}$ of order $d$ with $\zeta^k\ne1\ne\zeta^{k+1}$, the identity
$\zeta^k(\zeta\theta-2)=\theta-2$ for $\theta=2(\zeta^k-1)/(\zeta^{k+1}-1)$ gives
$g_k(\zeta\theta)=g_k(\theta)$; reducing modulo primes above $p$ yields pairs as in
Lemma~\ref{R6:lem:pair}. For $d\in\{3,4,6\}$ the relevant numbers are rational integers of binomial
shape, and Zsigmondy's theorem controls their prime factors:

\begin{lemma}[Four explicit families]\label{R6:lem:families}
Let $p$ be an odd prime and $\iota,\omega\in\overline{\F}_p$ roots of $X^2+1$, $X^2+X+1$. Put
\begin{align*}
M_k&=2^{(k+1)/2}-(-1)^{(k-1)/4} &&(k\equiv1\bmod4),\\
A_k&=1+(-1)^{(k+2)/4}\,2^{(3k+2)/2} &&(k\equiv2\bmod4),\\
B_k&=(-2)^{k+1}-1 &&(k\equiv1\bmod3),\\
C_k&=(-3)^{k/2}+1 &&(k\equiv2\bmod6).
\end{align*}
Lemma~\ref{R6:lem:pair} applies in each of the following cases:
\begin{enumerate}
\item[(I)] $k\equiv1\pmod4$, $p\mid M_k$: $d=4$, $\zeta=\iota$, $\theta=1-\iota$;
\item[(II)] $k\equiv2\pmod4$, $p\mid A_k$, $p\ne5$: $d=4$, $\zeta=\iota$, $\theta=2-2\iota$;
\item[(III)] $k\equiv1\pmod3$, $p\mid B_k$, $p\ne3$: $d=3$, $\zeta=\omega$, $\theta=-2\omega$;
\item[(IV)] $k\equiv2\pmod6$, $p\mid C_k$, $p\ne3$: $d=6$, $\zeta=-\omega^2$, $\theta=1-\omega$.
\end{enumerate}
\end{lemma}
\begin{proof}
In each case $\zeta\theta$ is the image of $\theta$ under $\iota\mapsto-\iota$ (resp.
$\omega\mapsto\omega^2$), namely $1+\iota$, $2+2\iota$, $-2\omega^2$, $1-\omega^2$, and we show that
$g_k(\theta)$ equals $\mp M_k$, $A_k$, $-B_k$, $C_k$ as an identity in $\Z[\iota]=\Z[X]/(X^2+1)$
(resp.\ $\Z[\omega]$); applying the automorphism, $g_k(\zeta\theta)$ is the same integer, so both
vanish mod $p$.
(I) $k=4t+1$: $(1-\iota)^2=-2\iota$, $(1-\iota)^4=-4$, and
$g_k(1-\iota)=-(1-\iota)^{k}(1+\iota)+1=1-2(1-\iota)^{4t}=1-(-1)^t2^{2t+1}$.
(II) $k=2h$, $h$ odd: $(2-2\iota)^2=-8\iota$ and
$g_k(2-2\iota)=(-8\iota)^h(-2\iota)+1=1+(-1)^{(h+1)/2}2^{3h+1}=A_k$.
(III) $\omega^k=\omega$ and $\omega(\omega+1)=-1$, so
$g_k(-2\omega)=(-2\omega)^k(-2)(\omega+1)+1=1-(-2)^{k+1}$.
(IV) $k=2h$, $h\equiv1\pmod3$: $(1-\omega)^2=-3\omega$ and $1-\omega-2=\omega^2$, so
$g_k(1-\omega)=(-3)^h\omega^h\omega^2+1=(-3)^h+1$.
Side conditions: the orders of $\iota,\omega,-\omega^2$ are $4,3,6$ for odd $p$ ($p\ne3$ for $\omega$ and $-\omega^2$); $\theta-1$ and
$\zeta\theta-1$ are $-\iota,\iota$ in (I) and $-\omega,-\omega^2$ in (IV); in (II) they are
$1\mp2\iota$, which vanish only if $\iota=\pm1/2$, i.e.\ $5=0$; in (III) they are $-1-2\omega$,
$-1-2\omega^2$, which vanish only if $3=0$. Finally $p\nmid d$: $p$ is odd, and $p\ne3$ in (III),
(IV) (note $3\nmid C_k$).
\end{proof}

We use the theorem of Bang and Zsigmondy \cite{R6-Zsigmondy}: for $a\in\{2,3\}$ and $\delta\ge2$,
$a^\delta-1$ has a prime divisor $q$ with $\Rsixord_q(a)=\delta$ (a \emph{primitive divisor}
$q_\delta(a)$), except for $(a,\delta)=(2,6)$ and $(3,2)$. Primitive divisors for distinct $\delta$ are
distinct, $q_2(2)=3$ and $q_4(2)=5$.

\begin{theorem}[\v{S}iurys's conjecture: partial resolution]\label{R6:thm:siurys}
Let $k\ge3$. In each of the following cases there are positive integers $a_0,\dots,a_{k-1}$
with greatest common divisor $1$ such that \emph{every} term of $S_k(a_0,\dots,a_{k-1})$ is
composite:
\begin{enumerate}
\item[(a)] $k\equiv 2,5,7,10,13,17,18,19,21\pmod{24}$, or $k\equiv 11,71\pmod{120}$;
\item[(b)] $k\not\equiv0,3,11\pmod{12}$ and $k\notin\mathcal E\cup\{5,7,10\}$, where
$\mathcal E=\mathcal E_1\cup\dots\cup\mathcal E_5$ with
$\mathcal E_1=\{k\equiv1\ (8):\ 2^{(k+1)/2}-1\text{ prime}\}$,
$\mathcal E_2=\{k\equiv14\ (24):\ k/2\text{ prime}\}$,
$\mathcal E_3=\{k\equiv4,10\ (12):\ k+1\text{ prime}\}$,
$\mathcal E_4=\{k\equiv6\ (24):\ \tau((3k+2)/4)\le4\}$,
$\mathcal E_5=\{2^j:j\ge3\}$; $\mathcal E$ has natural density $0$, so the conjecture holds on a set
of $k$ of lower density at least $3/4$;
\item[(c)] $3\le k\le 130$ (computer-certified; see Remark~\ref{R6:rem:certs}).
\end{enumerate}
\end{theorem}

\begin{proof}
By Lemma~\ref{R6:lem:assembly} it suffices to exhibit coverings by nonzero solutions mod distinct
primes. We use four constructions.

\emph{Construction I} ($k\equiv1\pmod4$, $M_k$ has two distinct prime factors $p_1,p_2$, necessarily
odd). Take the solution mod $2$ of Lemma~\ref{R6:lem:two} with $(a,b)=(1,3)$ (distinct mod $k+1\ge6$),
and by Lemma~\ref{R6:lem:families}(I) solutions mod $p_1,p_2$ vanishing on $1\bmod4$ and $3\bmod4$.
Since $k+1$ is even, every $n\equiv1$ or $3\pmod{k+1}$ is odd; so even $n$ are covered by $2$ and
odd $n$ by $p_1$ or $p_2$.

\emph{Construction II} ($k\equiv2\pmod4$, $A_k$ has four distinct prime factors $\ne5$, necessarily
odd): by Lemma~\ref{R6:lem:families}(II), solutions mod these primes vanishing on $0,1,2,3\bmod4$.

\emph{Construction III} ($k\equiv1\pmod3$, $B_k$ has three distinct prime factors $\ne3$, necessarily
odd): solutions vanishing on $0,1,2\bmod3$ (Lemma~\ref{R6:lem:families}(III)).

\emph{Construction IV} ($k\equiv2\pmod6$, $k\ge8$, $C_k$ has two distinct odd prime factors,
necessarily $\ge5$ since $3\nmid C_k$). Here $3\mid k+1$; take the solution mod $2$ with $(a,b)=(0,3)$ and solutions mod $p_1,p_2$
vanishing on $0\bmod6$ and $3\bmod6$. Every $n\equiv0,3\pmod{k+1}$ is $\equiv0\pmod3$, hence covered.

\emph{Construction V (class-wide coverings).} If $y$ is a $d$-periodic sequence over $\F_p$
with zero period sum and $y_n=y_{n-1}+\dots+y_{n-s}$ (indices mod $d$), then $n\mapsto y_{n\bmod d}$
is a solution mod $p$ for \emph{every} $k\equiv s\pmod d$, because
$\sum_{i=1}^{k}y_{n-i}=\frac{k-s}{d}\sum_{t<d}y_t+\sum_{i=1}^{s}y_{n-i}$ (this is the mechanism of
\cite[Lemma~4.1]{R6-Siurys}). Likewise, if $l\mid k+1$, every $l$-periodic $0/1$ pattern with an even
number of ones per period is a solution mod $2$ (Lemma~\ref{R6:lem:two}'s argument). For
$k\equiv11\pmod{120}$ we use: the prime $2$ with $x_n$ odd exactly for $n\equiv0,6\pmod{12}$; and
sequences with $s=11$ vanishing on $n\equiv4\pmod5$ ($p=19$, $d=15$), $n\equiv0\pmod5$ ($p=29$,
$d=20$), $n\equiv3\pmod5$ ($p=31$, $d=30$), $n\equiv2\pmod{10}$ ($p=521$, $d=40$) and
$n\equiv6\pmod{10}$ ($p=49921$, $d=120$). A term not divisible by $2$ has $n\equiv0,6\pmod{12}$, so $n$ is even; it is
covered by $29,31,19$ when $n\equiv0,3,4\pmod5$, and by $521$ or $49921$ when $n\equiv2$ or
$6\pmod{10}$. For
$k\equiv71\pmod{120}$ we use: $2$ with $x_n$ odd exactly for $n\equiv12,20\pmod{24}$ (so $n\equiv4\pmod
8$); $p=31$ ($d=30$, $s=11$) on $n\equiv0\pmod5$; $p=181$ ($d=60$, $s=11$) on $n\equiv4\pmod5$; and a
$40$-periodic sequence mod $3$ with $s=31$ vanishing at $n\equiv12,28,36\pmod{40}$, which are the
remaining classes $n\equiv4\pmod8$, $n\equiv1,2,3\pmod5$. The explicit periodic sequences are stored
in \texttt{siurys/fixed\_\allowbreak{}D120\_\allowbreak{}0\_\allowbreak{}119.json} and checked by \texttt{verify\_\allowbreak{}siurys\_\allowbreak{}fixed.py} (zero sum,
cyclic $s$-step relation, zeros, covering, and composite initial segments for three values of $k$ in
each class).

It remains to count prime factors with the Bang--Zsigmondy theorem.

\emph{I, $k\equiv5\pmod8$, $k\ge13$.} $M_k=2^m+1$ with $m=(k+1)/2$ odd, $m\ge7$; $3\mid M_k$, and
$q_{2m}(2)\ne3$ divides $2^{2m}-1$ but not $2^m-1$, hence divides $M_k$.
\emph{I, $k\equiv17\pmod{24}$.} $M_k=2^m-1$ with $m\equiv9\pmod{12}$: $7=q_3(2)$ and $q_m(2)\ne7$
divide $M_k$.
\emph{I, $k\equiv1\pmod 8$.} $M_k=2^m-1$; if it is not prime it has two distinct prime factors, as
$2^m-1$ is not a proper power ($\equiv3\pmod4$ excludes squares; $y^j+1=(y+1)(y^{j-1}-\dots+1)$ with
the second factor odd and $>1$ excludes odd $j\ge3$). The exceptions form $\mathcal E_1$.

\emph{II, $k\equiv2\pmod8$, $k\ge18$.} Write $k=8t+2$, $t\ge2$; then $A_k=1-2^M$ with
$M=12t+4=4(3t+1)$, $3\nmid M$. The primes $q_\delta(2)$, $\delta\mid M$, $\delta\notin\{1,4\}$, are
distinct, odd, $\ne5$, and number $\tau(M)-2$. If $u=3t+1$ is odd, $\tau(M)=3\tau(u)\ge6$; if
$u=2^bw$ ($w$ odd), $\tau(M)=(b+3)\tau(w)\ge6$ unless $w=1$, $b\le2$, i.e.\ $u=4$, $t=1$. So
$\tau(M)-2\ge4$.
\emph{II, $k\equiv6\pmod8$.} $A_k=2^{2u}+1$ with $u=(3k+2)/4$ odd; for $\delta'\mid u$, $\delta'>1$,
$q_{4\delta'}(2)$ divides $2^{4u}-1$ but not $2^{2u}-1$, hence divides $A_k$, and $\ne5$. This gives
$\tau(u)-1\ge4$ primes unless $\tau(u)\le4$ (the set $\mathcal E_4$ for $k\equiv6\pmod{24}$).

\emph{III, $k\equiv7\pmod{12}$, $k\ge19$.} $B_k=2^N-1$, $N=k+1\equiv8\pmod{12}$: the primes
$q_4(2)=5$, $q_{N/2}(2)$, $q_N(2)$ (with $N/2\ge10$) are distinct and $\ne3$.
\emph{III, $k\equiv4,10\pmod{12}$.} $N=k+1$ is odd with $3\nmid N$ and $-B_k=2^N+1$; for $\delta'\mid N$,
$\delta'>1$, $q_{2\delta'}(2)\mid 2^N+1$ and $\ne3$, giving $\tau(N)-1\ge3$ primes unless $N$ is
prime ($\mathcal E_3$); note $N=q^2$ is impossible, since $N\equiv5,11\pmod{12}$ while $q^2\equiv1\pmod{12}$
for primes $q\ge5$.
\emph{IV, $k\equiv8\pmod{12}$.} $h=k/2=2^vh'$ with $v\ge1$, $h'$ odd, and $C_k=3^h+1$; for
$\delta''\mid h'$ the prime $q_{2^{v+1}\delta''}(3)$ ($\ge5$) divides $3^{2h}-1$ but not $3^h-1$, so
divides $C_k$; this gives $\tau(h')\ge2$ primes unless $h$ is a power of $2$ ($\mathcal E_5$).
\emph{IV, $k\equiv2\pmod{12}$.} $h=k/2$ is odd and $-C_k=3^h-1$; the primes $q_\delta(3)$,
$\delta\mid h$, $\delta>1$, number $\tau(h)-1\ge2$ unless $h$ is prime ($\mathcal E_2$ when
$k\equiv14\pmod{24}$).

\emph{(a)} By the above, Construction I covers $k\equiv5,13,21,17\pmod{24}$ ($k\ge13$), II covers
$k\equiv2,10,18\pmod{24}$ ($k\ge18$), III covers $k\equiv7,19\pmod{24}$ ($k\ge19$); the remaining
$k\in\{5,7,10\}$ are in (c) (and in \cite{R6-Siurys}); Construction V covers $k\equiv11,71\pmod{120}$.

\emph{(b)} Residues mod $24$ not in (a) and not $\equiv0,3,11\pmod{12}$: $k\equiv1,9$ (I, exceptions
$\mathcal E_1$), $k\equiv14$ (IV, exceptions $\mathcal E_2$), $k\equiv4,16,22$ (III, exceptions in
$\mathcal E_3$), $k\equiv6$ (II, exceptions $\mathcal E_4$), $k\equiv8,20$ (IV, exceptions
$\mathcal E_5$). Each $\mathcal E_i$ has density $0$: $\mathcal E_1,\mathcal E_2,\mathcal E_3$ lie in
images of the set of primes under affine maps; for $\mathcal E_4$, $\tau(u)\ge2^{\omega(u)}$ and
$\omega(u)\to\infty$ on a set of density one (Hardy--Ramanujan \cite{R6-HR}), and $k\mapsto(3k+2)/4$ maps the class
$6\bmod24$ bijectively onto $5\bmod18$. (The decay for $\mathcal E_4$ is slow: integers with
$\tau(u)\le4$ have counting function $\asymp x\log\log x/\log x$, and about $64\%$ of the $u=18j+5$,
$j<10^5$, still have $\tau(u)\le4$; the other $\mathcal E_i$ are very sparse.)

\emph{(c)} For each $k$ in (c) an explicit covering was found by computer and is certified as in
Remark~\ref{R6:rem:certs}.
\end{proof}

\begin{remark}[Certificates for (c)]\label{R6:rem:certs}
For each $k$ in Theorem~\ref{R6:thm:siurys}(c) the certificate file for $k$ (directory
\texttt{siurys/certs} of the R6 area) lists positive integers $a_0,\dots,a_{k-1}$ and covering data: optionally the prime $2$ with $b$ (pattern of Lemma~\ref{R6:lem:two} with $a=0$), and
odd primes $p$ with either one class $c\bmod d$ (type P) or several classes mod $d$
(type R). The verifier script (\texttt{verify\_\allowbreak{}siurys\_\allowbreak{}certs.py}) trusts only the integers
$a_i$ and checks:
$\gcd(a_i)=1$; $\min a_i>$ every covering prime; for type P that $x_{c+dj}\equiv0\pmod p$ for
$0\le j<k$, which forces $x_{c+dj}\equiv0$ for all $j\ge0$ because $j\mapsto x_{c+dj}\bmod p$ satisfies
the order-$k$ recurrence given by the characteristic polynomial of $C^d$ ($C$ the companion matrix;
Cayley--Hamilton); for type R that $x_{n+d}\equiv x_n\pmod p$ for $n<k$ (hence for all $n$) and
$x_c\equiv0$ at the listed classes; for $2$ the $(k+1)$-periodicity and the parity pattern; and that the
union of the annihilated classes is $\Z/L\Z$. The primes come from two sources: pair patterns of
Lemma~\ref{R6:lem:pair} (prime factors below $10^7$ of
$\mathrm{N}_{\Q(\zeta_d)/\Q}\bigl(2^{k+1}(\zeta^s-1)^k\zeta^s(1-\zeta)+(\zeta^{s+1}-1)^{k+1}\bigr)$,
$s=k\bmod d$, the numerator of $g_k(\theta)$, for $d\le96$, and for $d=(k+1)m$ with $2\le m\le28$ when
$k+1$ is prime), and patterns from several roots of $f_k$ that are $d$-th roots of unity mod $p$
(type R); the coverings were found by depth-first search and integer programming. Of the $k$ in (c),
the cases $k\le10$ were known \cite{R6-Siurys}. The coverings use between $3$ and $33$ primes
(e.g.\ $k=13$: the primes $2,3,43$; $k=12$: eleven primes with moduli dividing $780$; $k=108$:
$33$ odd primes with moduli dividing $1680$). The hardest $k$ are those with $k+1$ prime and $k$
highly divisible. Then the prime $2$ leaves two classes modulo $k+1$; such a class meets every
residue class modulo any $d$ coprime to $k+1$, so the odd primes must either cover all of $\Z$ by
themselves or use moduli divisible by $k+1$; and by Remark~\ref{R6:rem:obstruction}(i) no odd prime
can use a modulus dividing $k$. For $k=72,108$ coverings by odd primes alone were found; for
$k=60,96,100$ the coverings use pair patterns with moduli divisible by $k+1$ (for $k=60$, where
every modulus dividing $60$ is unavailable to odd primes: $23$ primes, moduli $7,8,9,14,16,18,24,28,
36,63,84$ and $122,488,976$, $L=61\cdot1008$).
\end{remark}

\begin{remark}[Why $k\equiv0,3,11\pmod{12}$ remain open]\label{R6:rem:obstruction}
(i) \emph{For an odd prime $p$ and a divisor $d$ of $k$ or of $k+1$, no nonzero solution mod $p$
vanishes on a residue class modulo $d$.} Let $C$ be the companion matrix of $f_k$ over
$\overline{\F}_p$, so a solution is $y_n=e_1C^nv$ and vanishing on $c\bmod d$ means
$e_1C^c(C^d)^jv=0$ for all $j\ge0$, i.e.\ $v$ lies in the largest $C^d$-invariant subspace $U$ of
$\ker(e_1C^c)$. Every root $\theta$ of $f_k$ satisfies $\theta^k=1/(2-\theta)$ and
$\theta^{k+1}=\theta/(2-\theta)$ (from $g_k(\theta)=0$, and $g_k(2)=1$), so $\theta\mapsto\theta^d$ is
injective on the roots. If $p\mid d$ then $p\mid k(k+1)$ and $g_k'\equiv X^k$ resp.\ $2X^{k-1}$ shows that
the roots are simple, so $C^d$ is diagonalizable with distinct eigenvalues; if $p\nmid d$, each Jordan
block $J_m(\theta)$ of the non-derogatory matrix $C$ has $J_m(\theta)^d\sim J_m(\theta^d)$. In both cases
$C^d$ is non-derogatory, so $C$ (which commutes with $C^d$) is a polynomial in $C^d$ and $U$ is
$C$-invariant; then $e_1C^{c+i}U=0$ for all $i\ge0$, i.e.\ the solutions with initial vector in $U$
vanish from index $c$ on, whence $U=0$ ($C$ is invertible). In particular modulus $2$ is available only
to the prime $2$, and for $k\equiv0$ or $11\pmod{12}$ none of the moduli $3,4,6,12$ is available to odd
primes, so none of the families of Lemma~\ref{R6:lem:families} applies.
(ii) For $k\equiv3\pmod{12}$ the moduli $3\mid k$ and $4\mid k+1$ are excluded by (i), so among the
moduli $d\in\{3,4,6\}$ (those with $\theta$ quadratic) only $d=6$ remains: $\zeta=-\omega^2$,
$\theta=4(1-\omega)/3$, for which $g_k(\theta)=1-2^{2k+1}/3^{(k+1)/2}$ is rational (and the norm computed by our search equals
$(2^{2k+1}-3^{(k+1)/2})^2$ for all $k\equiv3\pmod{12}$, $k\le63$). The integer
$2^{2k+1}-3^{(k+1)/2}$ is not a difference of like powers, and we have no primitive-divisor input.
(iii) For other moduli the relevant integers are norms of
$2^{k+1}(\zeta^k-1)^k\zeta^k(1-\zeta)+(\zeta^{k+1}-1)^{k+1}$ from fields $\Q(\zeta_d)$ with units of
infinite order, for which no Zsigmondy-type theorem is available. For example, the pair patterns with
ratio of order dividing $2(k+1)$ but not $k+1$ (then $\theta=1+\zeta^{-1}$, $\zeta\theta=1+\zeta$)
come exactly from the prime factors of $R_k=\mathrm{Res}(f_k,(X-1)^{k+1}+1)$; for $3\le k\le60$ one
finds $R_k=\pm W_k^2$ with $W_k\in\Z$, and $W_k$ is prime exactly for $k=3,4,7,12,16$ in this range
(e.g.\ $W_3=7$, $W_4=29$, $W_7=4241$), so for these $k$ this source supplies a single prime.
Thus the full conjecture needs either new uniform families for $k\equiv0,3,11\pmod{12}$ or control of
prime factors of such norms; the exceptional sets $\mathcal E_i$ would also have to be handled.
\end{remark}

%=====================================================================================
\subsection{Further targets examined}
\label{R6:subsec:other}

\begin{remark}[A conjecture that is already a classical theorem]\label{R6:rem:IKLP}
Ismailescu, Ko, Lee and Park \cite[Conjecture~6.2]{R6-IKLP} state: ``Let $(u_n)_{n\ge0}$ be the Lucas
sequence associated to some $|a|\ge3$ and $b=-1$. Then for any two different primes $p$ and $q$ we
have that $u_p$ and $u_q$ are relatively prime.'' Here $u_0=0$, $u_1=1$, $u_{n+1}=au_n-u_{n-1}$, i.e.\
$u_n=U_n(a,1)$ in Lucas's notation. Since $\gcd(a,1)=1$, $(U_n)$ is a strong divisibility sequence,
$\gcd(U_m,U_n)=|U_{\gcd(m,n)}|$ (Lucas \cite{R6-Lucas}), so $\gcd(u_p,u_q)=|u_1|=1$ for all distinct
primes $p,q$ (indeed for all coprime indices, and for every integer $a$). We record this only to note
that the statement is not open; no new mathematics is involved.
\end{remark}

\begin{remark}[Varona's lacunary prime-representing series]\label{R6:rem:varona}
Varona \cite[\S1]{R6-Varona} discusses $\beta=\sum_{k\ge1}p_k10^{-2^k}$ and
$\beta'=\sum_{k\ge1}p_kr^{-k^2}$ ($p_k$ the $k$-th prime) and writes that ``it is possible that these
$\beta$ are transcendental numbers'', the defining series being unsuitable for Liouville's or Roth's
theorems. For $\beta$ this follows at once from Ridout's theorem \cite{R6-Ridout} (for algebraic
$\xi$ and a finite set $S$ of primes, $|\xi-P/Q|<Q^{-1-\varepsilon}$ has finitely many solutions with $Q$
an $S$-unit): with $Q_K=10^{2^K}$ and $P_K=Q_K\sum_{k\le K}p_k10^{-2^k}\in\Z$, using $p_k\le2^k$,
\[
0<\beta-\frac{P_K}{Q_K}\le\sum_{k>K}2^k10^{-2^k}\le 2^{K+2}Q_K^{-2}<Q_K^{-3/2}\qquad(K\ge2),
\]
infinitely many distinct approximations, while $\beta\notin\Q$. So $\beta$ is transcendental. The
second number $\beta'$ is not covered (the approximation exponent tends to $1$); its transcendence remains
open. We list this as a minor, partial item.
\end{remark}

\begin{remark}[Open questions not closed here]\label{R6:rem:notclosed}
(1) \cite[\S7]{R6-HIKK} asks: ``Can one find a rational number or a transcendental number that
satisfies this condition?'' (that $\fl{\alpha^n}$ is composite for all but finitely many $n$). The
existence of uncountably many such $\alpha$ (hence transcendental ones) is attributed in
\cite{R6-HIKK} to Baker and Harman, without an explicit example; an explicit example is not produced
here. (2) The second question of \cite[\S7]{R6-HIKK} (an $\alpha>1$ with $\fl{\alpha^n}$ prime for all
large $n$), the totally real case of Saito's Problem~1.1 \cite{R6-Saito}, and the irrationality of
$\xi(2^k)$ were not attacked; the latter two are partly treated in \cite{K-031,K-042} and the first is
far beyond the methods available.
\end{remark}


%% ==================================================================
%% Section: Floquet isospectrality
%% ==================================================================
% R3: Floquet isospectrality to the zero potential (area R3).  LaTeX fragment, no preamble.
\providecommand{\Z}{\mathbb{Z}}
\providecommand{\C}{\mathbb{C}}
\providecommand{\R}{\mathbb{R}}
\providecommand{\Q}{\mathbb{Q}}
\providecommand{\T}{\mathbb{T}}
\providecommand{\Hdiag}[1]{H_{D}^{#1}}
\providecommand{\Hcirc}{H^{\circ}}
\providecommand{\supp}{\operatorname{supp}}
\providecommand{\diag}{\operatorname{diag}}

\section{Floquet isospectrality to the zero potential: conjugate-symmetric potentials and period lattices}
\label{R3:sec}

This section treats the two items left open in the paper of Faust, Friedman,
O'Malley, Ramos and Sharma~\cite{R3-FFORS} after the results of the earlier
manuscript~\cite{K-034}. Part~XII of~\cite{K-034} completed the classification of
rectangular periods. Part~XIII proved Conjecture~1, and resolved Conjecture~3 by
proving its multiplicity formula and disproving its singularity patterns. The two
remaining items are Conjecture~2 (existence of conjugate-reflection-symmetric
potentials, open for odd periods) and Problem~1 (classification of all period
lattices). The tools are those of method~M4 (coordinate averaging, induction on
the dimension, eigenvalue perturbation), together with a degeneration argument on
diagonal fibres, interval-Newton certificates and exact computer algebra.
\emph{Neither target is closed here; both results are partial.}

For Conjecture~2:
\begin{itemize}
\item we give exact explicit solutions for $q=5$ and $q=9$;
\item we give computer-assisted (interval-arithmetic) proofs for every odd period
$5\le q\le69$;
\item it follows that Conjecture~2 holds for every period that is even (already
known) or has a divisor in $[5,69]$.
\end{itemize}
Primes larger than $69$ remain open.

For Problem~1:
\begin{itemize}
\item we prove a first-order structure theorem, valid for every lattice in every
dimension;
\item we deduce rigidity for an infinite class of non-rectangular lattices. It
contains the families treated in~\cite{R3-FFORS} and the $(3\Z)^d$ theorem
of~\cite{K-034};
\item we reduce the expected classification to a single ``core statement'';
\item using the structure theorem and exact computer algebra, we verify the
expected planar classification for all lattices of index at most $12$.
The previously reported check was index at most $8$.
\end{itemize}

\subsection{Setting and the plane-wave form of the Floquet matrices}
\label{R3:ss:setting}

Let $\Gamma\subset\Z^d$ be a full-rank sublattice, $G=\Z^d/\Gamma$, $N=|G|$,
and let $V:\Z^d\to\C$ be $\Gamma$-periodic. Put
$(\Delta f)(n)=\sum_{\|m-n\|_1=1}f(m)$. For $z\in\T^d$ let $L_V(z)$ be the
Floquet matrix of $\Delta+V$ (acting on functions on a fundamental domain with
the Bloch conditions $f(n+\gamma^{(i)})=z_if(n)$ for a basis $\gamma^{(i)}$ of
$\Gamma$). As in~\cite{R3-FFORS,R3-CFK}, $V$ is \emph{Floquet isospectral to
$0$} if $\det(\lambda-L_V(z))=\det(\lambda-L_0(z))$ for all $z\in\T^d$ and
$\lambda\in\C$.

For $\xi\in(\C^*)^d$ and $n\in\Z^d$ write $\xi^n=\prod_j\xi_j^{n_j}$. Let
\[
 H=\{\xi\in(\C^*)^d:\ \xi^\gamma=1\ \text{for all }\gamma\in\Gamma\},\qquad
 b_h=\frac1N\sum_{n\in G}V(n)h^{-n}\quad(h\in H),
\]
so that $|H|=N$, $H$ consists of roots of unity, and $V(n)=\sum_{h\in H}b_hh^n$.
Let $g_j=g_j(\Gamma)\ge1$ be the positive generator of the $j$-th coordinate
projection of $\Gamma$, and define the subsets
\begin{gather*}
 K_j=\{h\in H: h_i=1\ (i\ne j)\}\cong\mu_{g_j},\qquad
 \Hcirc=\{h\in H: h_j\neq1\ \text{for all }j\},\\
 \Hdiag{\varepsilon}=\{h\in H:\ h_1^{\varepsilon_1}=\dots=h_d^{\varepsilon_d}\}
\end{gather*}
for $\varepsilon\in\{\pm1\}^d$; $\Hdiag{\varepsilon}$ is a subgroup.

\begin{lemma}[plane-wave form]\label{R3:lem:planewave}
For $x\in(\C^*)^d$ put $D_\xi(x)=\sum_{j=1}^d\bigl(x_j\xi_j+(x_j\xi_j)^{-1}\bigr)$,
$D(x)=\diag(D_\xi(x))_{\xi\in H}$ and $B=(b_{\zeta\xi^{-1}})_{\zeta,\xi\in H}$.
Then $V$ is Floquet isospectral to $0$ if and only if
$\det(\lambda-D(x)-B)=\det(\lambda-D(x))$ for all $x\in(\C^*)^d$, $\lambda\in\C$.
\end{lemma}

\begin{proof}
Fix $x$ and let $W_x$ be the space of $f:\Z^d\to\C$ with $f(n+\gamma)=x^\gamma f(n)$
($\gamma\in\Gamma$). The functions $e_\xi(n)=(x\xi)^n$, $\xi\in H$, lie in $W_x$
and are linearly independent (distinct characters), hence form a basis
($\dim W_x=N$). One has $\Delta e_\xi=D_\xi(x)e_\xi$ and
$Ve_\xi=\sum_hb_he_{\xi h}$, so $\Delta+V$ acts on $W_x$ by $D(x)+B$. Restriction
to a fundamental domain identifies $W_x$ with the domain of $L_V(z)$ for
$z_i=x^{\gamma^{(i)}}$. The homomorphism $x\mapsto(x^{\gamma^{(i)}})_i$ of
$(\C^*)^d$ is surjective, because the exponent matrix is nonsingular. Both
characteristic polynomials are Laurent polynomials in $z$, so equality on $\T^d$ is
equivalent to equality on $(\C^*)^d$.
\end{proof}

Translations of $V$ multiply $b_h$ by $h^{n_0}$, and $\Gamma\subset\Gamma'$ with
$V$ $\Gamma'$-periodic makes $D(x)+B$ block diagonal along the cosets of
$\Gamma'^{\perp}\subset H$; hence translates of isospectral potentials are
isospectral, and isospectrality for $\Gamma'$ implies isospectrality for every
sublattice $\Gamma\subset\Gamma'$.

\subsection{Target 1: conjugate-reflection-symmetric potentials}
\label{R3:ss:conj2}

\begin{problem}[{\cite[Section~4, Conjecture~2]{R3-FFORS}}]
\label{R3:prob:conj2}
``There exist nonzero $q\Z$-periodic potentials $V$ of the form
$(v_1,v_2,\ldots,\overline{v_2},\overline{v_1})$ for $q>3$ such that $\Delta$ and
$\Delta+V$ are Floquet isospectral.''
\end{problem}

The authors prove the conjecture for even $q$, using the support-four potentials
of~\cite{R3-FLMPRTTZ}. Table~1 of~\cite{R3-FFORS} lists numerical counts of
solutions of this form for $q\le7$: $4$ for $q=5$ and $120$ for $q=7$. The odd case was open, and is not covered
by~\cite{K-034}. For odd $q=2m+1$ we index sites by $\Z/q$ and place the
middle entry $v_{m+1}$ at $0$. A cyclic shift, which preserves isospectrality,
turns the required form into
\begin{equation}\label{R3:eq:tau}
 V_{-j}=\overline{V_j}\qquad(j\in\Z/q),\qquad\text{in particular }V_0\in\R .
\end{equation}

\begin{lemma}[one-dimensional reduction]\label{R3:lem:1d}
Let $q\ge3$ and let $C_q$ be the adjacency matrix of the $q$-cycle. A $q$-periodic
$V$ is Floquet isospectral to $0$ if and only if
$\det(\lambda-C_q-\diag V)=\det(\lambda-C_q)$.
\end{lemma}

\begin{proof}
In $L_V(z)$ the parameter enters only through the two corner entries $z^{\pm1}$
(up to the choice of convention). In the permutation expansion of
$\det(\lambda-L_V(z))$, a term containing a corner entry but not its transpose
must use all $q$ edges of the cycle in one direction. These are the two cyclic
permutations, which contain no diagonal entry and are therefore independent of
$V$. Terms using a corner entry together with its transpose contain
$z\cdot z^{-1}=1$. Hence $\det(\lambda-L_V(z))-\det(\lambda-L_0(z))$ does not
depend on $z$, and it suffices to take $z=1$.
\end{proof}

\begin{lemma}[repetition]\label{R3:lem:repeat}
Let $q_0\mid q$ and let $W:\Z/q_0\to\C$ satisfy $W_{-j}=\overline{W_j}$ and be Floquet
isospectral to $0$ (period $q_0$). Then $V_j:=W_{j\bmod q_0}$ ($j\in\Z/q$)
satisfies~\eqref{R3:eq:tau} and is Floquet isospectral to $0$ for the period $q$.
\end{lemma}

\begin{proof}
Reduction modulo $q_0$ commutes with $j\mapsto-j$, which gives~\eqref{R3:eq:tau}.
The potential is $q_0\Z$-periodic and $q_0\Z$-isospectral to $0$, hence
$q\Z$-isospectral by the sublattice remark after Lemma~\ref{R3:lem:planewave}.
\end{proof}

\begin{theorem}[explicit odd solutions]\label{R3:thm:explicit}
\textup{(a)} Let $q=5$, $\omega=e^{2\pi i/5}$, $t=24^{1/10}>0$, $u=5^{1/4}>0$, and
\begin{gather*}
 b_{1}=-\frac{t^4}{2u^3},\quad b_{-1}=-\frac{t^6}{2u^3},\quad
 b_{2}=\frac{6}{u^3t^2},\quad b_{-2}=-\frac{t^2}{u^3},\\
 V_j=\sum_{h\in\{\pm1,\pm2\}}b_h\,\omega^{hj}\quad(j\in\Z/5).
\end{gather*}
Then $V\ne0$ satisfies~\eqref{R3:eq:tau} and is Floquet isospectral to $0$.
Numerically,
\begin{gather*}
 V_0=-1.1541016\ldots,\qquad V_{\pm1}=-0.7878062\ldots\pm1.3408501\ldots i,\\
 V_{\pm2}=1.3648570\ldots\mp1.1625702\ldots i.
\end{gather*}

\textup{(b)} Let $q=9$ and $V_0=V_{\pm4}=0$, $V_{\pm1}=e^{\mp i\pi/6}$,
$V_{\pm2}=-e^{\mp i\pi/6}$, $V_{\pm3}=\mp i$. Then $V\neq0$
satisfies~\eqref{R3:eq:tau} and is Floquet isospectral to $0$.
\end{theorem}

\begin{proof}
(a) All $b_h$ are real, so $\overline{V_j}=\sum_hb_h\omega^{-hj}=V_{-j}$. By
Lemma~\ref{R3:lem:1d} it suffices to prove
$\det(\lambda-C_5-\diag V)=\det(\lambda-C_5)$. Multiplying by a power of $t$ and
$u$ clears denominators, and the difference becomes a polynomial in
$\Q[t,u,\omega,\lambda]$. It lies in the ideal generated by $t^{10}-24$,
$\Phi_5(\omega)$ and $u^2-(1+2\omega+2\omega^4)$; we verified this by reducing
it to zero with an exact lex Gr\"obner basis. All three generators vanish at
$t=24^{1/10}$, $\omega=e^{2\pi i/5}$, $u=5^{1/4}$, because
$1+2(\omega+\omega^4)=1+4\cos\frac{2\pi}5=\sqrt5=u^2$. Hence the identity holds
at the stated values.

(b) Symmetry~\eqref{R3:eq:tau} is immediate. All entries lie in
$\Q(\zeta_{12})$ with $\zeta_{12}=e^{i\pi/6}$. The difference
$\det(\lambda-C_9-\diag V)-\det(\lambda-C_9)$, computed exactly as a polynomial in
$\zeta_{12}$ and $\lambda$, is divisible by $\Phi_{12}(\zeta_{12})=\zeta_{12}^4-\zeta_{12}^2+1$.
It therefore vanishes, and Lemma~\ref{R3:lem:1d} applies.
\end{proof}

\begin{remark}
The potential in (a) was found numerically and then identified exactly. The
products $b_hb_{-h}=\pm6/(5\sqrt5)$ and the ratios $(b_1/b_{-1})^5=1/24$,
$(b_2/b_{-2})^5=-27/2$ were recognised by integer-relation detection. In
the normalisation $b_h=t^{-h}s_h$ one has
$(s_{-2},s_{-1},s_1,s_2)=5^{-3/4}(-1,-\sqrt6,-\sqrt6,6)$. The solution in (b) has
support six and unimodular entries. A search for sparse solutions of
form~\eqref{R3:eq:tau} found no analogue for $q=7,11,13$.
\end{remark}

\begin{theorem}[computer-assisted, odd periods]\label{R3:thm:certified}
For every odd $q$ with $5\le q\le69$ there is a nonzero $V:\Z/q\to\C$
satisfying~\eqref{R3:eq:tau} that is Floquet isospectral to $0$.
\end{theorem}

\begin{proof}
Write $V_0=r$, $V_{\pm j}=x_j\pm iy_j$ ($1\le j\le m$), and let
$F:\R^q\to\R^q$ be the vector of real parts of the coefficients of
$\lambda^{q-1},\dots,\lambda^0$ in $\det(\lambda-C_q-\diag V)-\det(\lambda-C_q)$.
These coefficients are real on the slice~\eqref{R3:eq:tau}: if $R$ is the
permutation matrix of $j\mapsto-j$, then $RC_qR=C_q$ and $R\diag(V)R=\diag(\overline V)$,
so $R(C_q+\diag V)R=\overline{C_q+\diag V}$, and the characteristic polynomial of
$C_q+\diag V$ equals its own complex conjugate. Hence a zero of $F$ is a zero of all
$q$ complex coefficients, and by Lemma~\ref{R3:lem:1d} a zero of $F$ with
$(r,x,y)\ne0$ is the required potential.
For $q\ge29$ the coefficient system is badly conditioned in floating point, so we
use an equivalent system instead. Put $A=C_q+\diag V$,
$G(\lambda)=\det(\lambda-A)/\det(\lambda-C_q)-1$, $z_k=e^{(1+2\pi ik)/q}$ and
$\lambda_k=z_k+z_k^{-1}$ ($0\le k<q$). These are $q$ distinct points on a Bernstein
ellipse around $[-2,2]$, and
$\det(\lambda_k-C_q)=z_k^q+z_k^{-q}-2=e+e^{-1}-2\neq0$. The choice $|z_k|^q=e$
balances the Chebyshev coefficients of the difference polynomial and makes the
system well conditioned. Let
\[
E(w)=\bigl(\operatorname{Re}G(\lambda_0),\ \operatorname{Re}G(\lambda_k),\ \operatorname{Im}G(\lambda_k)\ (1\le k\le m)\bigr)\in\R^q .
\]
The polynomial $\det(\lambda-A)-\det(\lambda-C_q)$ has degree $<q$ and, on the slice,
real coefficients. Hence $G(\lambda_0)$ is real, the values at
$\lambda_{q-k}=\overline{\lambda_k}$ are conjugate, and $E(w)=0$ forces the polynomial
to vanish at all $q$ points $\lambda_k$, hence identically. So the zeros of $E$ and
$F$ on the slice coincide.

For each $q$ the script \texttt{R3\_\allowbreak{}conj2\_\allowbreak{}krawczyk.py} finds a floating-point zero
by damped Newton iteration from random starts. It refines the zero by Newton steps
at $400$-bit precision to a centre $w_0$ with exactly representable coordinates,
and fixes $\rho=10^{-30}$. In Arb ball arithmetic~\cite{R3-Arb} it then evaluates
$F(w_0)$ (resp.\ $E(w_0)$) and an enclosure $J(X)$ of the Jacobian over the box
$X=w_0+[-\rho,\rho]^q$. The Jacobian uses
$\partial_{V_j}\det(\lambda-A)=-\det(\lambda-A_{\hat j\hat j})$
(resp.\ $\partial_{V_j}G=-(G+1)[(\lambda-A)^{-1}]_{jj}$) and the chain rule. With
$Y$ the midpoint matrix of the inverse of the midpoint Jacobian, it checks, with
rigorous ball comparisons, the Krawczyk inclusion
$K(X)=w_0-YF(w_0)+(I-YJ(X))(X-w_0)\subset w_0+(-0.99\rho,0.99\rho)^q\subset\operatorname{int}X$
(with $E$ in place of $F$ when $q\ge29$). By the Krawczyk--Moore
theorem~\cite{R3-Krawczyk,R3-Moore}, the system then has a unique zero in $X$.
Since $|w_0|-\sqrt q\,\rho>0$, this zero is nonzero.
\end{proof}

As a consistency check, for $q=5$ the certified zero coincides, to the printed
digits, with the reflection $V_j\mapsto V_{-j}$ of the exact solution of
Theorem~\ref{R3:thm:explicit}(a). The residuals at the centres are below
$10^{-100}$. For $29\le q\le45$ the log also records the residual of the coefficient
system $F$ at the same centre (below $10^{-90}$; this is not needed for the proof).

\begin{table}[ht]
\centering\small
\begin{tabular}{ll}
\hline
period $q$ & certificate\\
\hline
$5$, $9$ & exact (Theorem~\ref{R3:thm:explicit}); also Krawczyk with $F$\\
odd $7\le q\le27$ & Krawczyk with $F$, $\rho=10^{-30}$, $400$-bit balls\\
odd $29\le q\le69$ & Krawczyk with $E$, $\rho=10^{-30}$, $400$-bit balls\\
\hline
\end{tabular}
\caption{Certificates for Theorem~\ref{R3:thm:certified}
(logs \texttt{R3\_\allowbreak{}conj2\_\allowbreak{}krawczyk.log} for $q\le27$ and
\texttt{R3\_\allowbreak{}conj2\_\allowbreak{}krawczyk\_\allowbreak{}ellipse.log} for $q\ge29$).}
\label{R3:tab:cert}
\end{table}

\begin{corollary}[partial resolution of Problem~\ref{R3:prob:conj2}]\label{R3:cor:conj2}
Conjecture~2 of~\cite{R3-FFORS} holds for every $q>3$ that is even or has a divisor
in $[5,69]$. In particular it holds for every $q>3$ divisible by $5$, $7$ or
$9$, hence for every power of $3$ that is at least $9$. The only odd $q>3$ not
covered are $q=n$ and $q=3n$ with $n>1$ odd and all prime factors of $n$ larger
than $69$, for instance every prime $q>69$.
\end{corollary}

\begin{proof}
Even $q$: \cite{R3-FFORS}. Otherwise apply Lemma~\ref{R3:lem:repeat} to the
solution for a divisor $q_0\in[5,69]$ given by Theorems~\ref{R3:thm:explicit}
and~\ref{R3:thm:certified}. If an odd $q>3$ has no divisor in $[5,69]$, then
$9\nmid q$ and no prime $5\le p\le69$ divides $q$, which gives the stated form.
\end{proof}

\begin{remark}[why degree theory does not settle all odd $q$]\label{R3:rem:degree}
Let $q=2m+1$ and $F$ be as in the proof of Theorem~\ref{R3:thm:certified}.
The leading forms of $F$ are $\pm e_k(V)$, with $V=V(w)$. Their only common real
zero is $0$, so $F$ is proper. The degree of $F$ on large spheres equals that of the
leading map, which is $0$. Indeed, on the slice the roots of $\prod_j(\lambda-V_j)$
are $V_0$ together with pairs $V_j,\overline{V_j}$, so this polynomial has at most
one simple real root. The leading map therefore misses the nonempty open set of
coefficient vectors of polynomials with two simple real roots. The real zero set
of $F$ is finite (the complex system has finitely many solutions), so the indices
of the real zeros sum to $0$. The bipartite-type gauge $V\mapsto-V$ (valid for odd $q$
together with $z\mapsto-z$) gives $F(-w)=\Sigma F(w)$ with
$\Sigma=\diag((-1)^k)_{k=1}^q$. Comparing degrees on a small sphere around $0$
gives $\operatorname{ind}_0F=(-1)^m\operatorname{ind}_0F$. Thus for
$q\equiv3\pmod4$ the local index at the degenerate zero $0$ vanishes and no
Brouwer-degree count can detect the nonzero real solutions, although Table~1
of~\cite{R3-FFORS} reports $120$ of them for $q=7$. For $q\equiv1\pmod 4$ the
symmetry imposes nothing, and $\operatorname{ind}_0F\neq0$ would force nonzero real
zeros. For $q=5$ a certified real-root isolation of the full system (msolve~\cite{R3-msolve})
finds exactly four nonzero real zeros. They form one orbit under $V\mapsto V(-\,\cdot)$
and $V\mapsto-V$, both of which preserve the local index when $q=5$, and they are
nondegenerate. Hence $|\operatorname{ind}_0F|=4$. We do not know how to compute $\operatorname{ind}_0F$ in
general. A proof of Conjecture~2 for all odd $q$ therefore seems to require an
explicit infinite family, a computation of the local index for $q\equiv1\pmod4$, or
an equivariant refinement of the degree for $q\equiv3\pmod4$.
\end{remark}

\subsection{Target 2: classification of period lattices}
\label{R3:ss:problem1}

\begin{problem}[{\cite[Section~6, Problem~1]{R3-FFORS}}]\label{R3:prob:classify}
``For the $\Z^d$-periodic grid graph, classify all periods $\Gamma$ such that
there exist nonzero $\Gamma$-periodic potentials that are Floquet isospectral to
the discrete Laplacian.'' The authors add: ``Problem~1 has been solved in the
case where $d=1$ [CFK], but it is still open for $d\geq2$.'' For $d=2$ they
observe that the undecided lattices (those not contained in some $q\Z\oplus\Z$ or
$\Z\oplus q\Z$ with $q>3$) have Hermite normal form $(a,0)\Z\oplus(b,c)\Z$ with
$0\le b<a$ and $\gcd(a,b),c\in\{1,2,3\}$. They write: ``We symbolically checked all
possible combinations of $a,b,c$ where $ac\leq8$'' and ``found no examples with
nonzero potentials isospectral to $\mathbf 0$''.
\end{problem}

The rectangular case is settled. Existence when some period exceeds $3$, and
rigidity when all periods are at most $3$ and at most two equal $3$, are
\cite[Theorem~C]{R3-CFK}. The remaining cases, with at least three periods equal
to $3$, are covered by \cite[Part~XII]{K-034}, through the $(3\Z)^d$ theorem and a
reduction. The natural expected answer is the following.

\begin{conjecture}[expected classification]\label{R3:conj:classification}
Let $\Gamma\subset\Z^d$ be a full-rank sublattice. A nonzero $\Gamma$-periodic
potential Floquet isospectral to $0$ exists if and only if $\max_jg_j(\Gamma)\ge4$.
\end{conjecture}

\begin{proposition}[the ``if'' direction]\label{R3:prop:if}
If $g_j(\Gamma)\ge4$ for some $j$, a nonzero $\Gamma$-periodic potential Floquet
isospectral to $0$ exists.
\end{proposition}

\begin{proof}
By \cite[Theorem~4.3]{R3-CFK} (see also \cite{R3-FLMPRTTZ} for even periods) there
is a nonzero $W:\Z/g_j\to\C$ that is Floquet isospectral to $0$. Put
$V(n)=W(n_j)$. Then $V$ is periodic for $\Gamma'=\{n:n_j\in g_j\Z\}\supset\Gamma$.
For $\Gamma'$ the operator separates: in Lemma~\ref{R3:lem:planewave} the dual
group is $\mu_{g_j}$ in the $j$-th coordinate, and $D(x)+B$ is the $1$-dimensional
Floquet matrix plus scalars. So $V$ is $\Gamma'$-isospectral to $0$, and hence
$\Gamma$-isospectral by the sublattice remark.
\end{proof}

\begin{theorem}[first-order structure, all lattices]\label{R3:thm:structure}
Let $d\ge2$, let $\Gamma\subset\Z^d$ be any full-rank sublattice, and let $V$ be
$\Gamma$-periodic and Floquet isospectral to $0$. Then:
\begin{enumerate}
\item[(i)] $b_1=0$;
\item[(ii)] \emph{(averaging)} for each $j$, the potential on $\Z^{d-1}/p_j(\Gamma)$
($p_j$ deletes the $j$-th coordinate) whose Fourier coefficients are
$b|_{\{h\in H:h_j=1\}}$ (equivalently: the average of $V$ over the fibres of
$\Z^d/\Gamma\to\Z^{d-1}/p_j(\Gamma)$) is $p_j(\Gamma)$-Floquet isospectral to $0$;
\item[(iii)] \emph{(diagonal fibres)} $b_h=0$ for every $h\in\Hdiag{\varepsilon}$ and
every $\varepsilon\in\{\pm1\}^d$;
\item[(iv)] \emph{(pair products)} $b_hb_{h^{-1}}=0$ for every $h\in\Hcirc$.
\end{enumerate}
\end{theorem}

\begin{proof}
(i) is the trace of Lemma~\ref{R3:lem:planewave} at order $\lambda^{N-1}$.

(ii) Take $j=d$. Fix $x'=(x_1,\dots,x_{d-1})$ and let $x_d=s\to\infty$. Group
$H$ by the value $r$ of $\xi_d$. Each group $P_r$ is a coset $\xi^0H^{(d)}$ of
$H^{(d)}=\{h:h_d=1\}$, and $H^{(d)}$ is the dual group of $p_d(\Gamma)$. For
$\xi\in P_{r'}$, $D_\xi=sr'+O(1)$. Inside $P_r$, the $O(1)$ part of
$D_{\xi^0h}$ is $\sum_{i<d}\bigl((x_i\xi^0_i)h_i+((x_i\xi^0_i)h_i)^{-1}\bigr)$, and
the coupling is $b_{hh'^{-1}}$ with $hh'^{-1}\in H^{(d)}$. This block is the
plane-wave Floquet matrix $M'(x'\xi^0)$ of the averaged potential for the lattice
$p_d(\Gamma)$. Substituting $\lambda=sr+\mu$ and taking the Schur complement with
respect to the other groups, whose diagonal entries are $s(r-r')+O(1)$ and whose
couplings are $O(1)$, gives
\[
 \det(\lambda-D-B)=\prod_{r'\neq r}(s(r-r'))^{|P_{r'}|}\bigl(\det(\mu-M'(x'\xi^0))+O(s^{-1})\bigr),
\]
and the same with $B=0$. Equality of the left sides for all $s$ forces
$\det(\mu-M'(y))=\det(\mu-M'_0(y))$ for all $y\in(\C^*)^{d-1}$, $\mu\in\C$.

(iii) Fix $\varepsilon$ and $\xi^0\in H$. Since $d\ge2$, choose
$a\in(\C^*)^d$ with $\sum_ja_j(\xi^0_j)^{\varepsilon_j}=0$ and otherwise generic.
Then, for $\xi\in H$, $\sum_ja_j\xi_j^{\varepsilon_j}=0$ holds exactly when
$(\xi_j^{\varepsilon_j})_j\in\C\cdot((\xi^0_j)^{\varepsilon_j})_j$, i.e.\ when
$\xi\in P:=\xi^0\Hdiag{\varepsilon}$. Indeed, for $\xi$ outside $P$ the condition
cuts out a proper subspace of the hyperplane, and there are finitely many such
$\xi$. Put $x_j=sa_j$ if $\varepsilon_j=1$ and $x_j=(sa_j)^{-1}$ if
$\varepsilon_j=-1$. Then
$D_\xi=s\alpha_\xi+s^{-1}\beta_\xi$ with $\alpha_\xi=\sum_ja_j\xi_j^{\varepsilon_j}$
and $\beta_\xi=\sum_ja_j^{-1}\xi_j^{-\varepsilon_j}$. Thus $\alpha=0$ on $P$ and
$\alpha\ne0$ off $P$. For bounded $\mu$ and large $s$, the Schur complement of
$\mu-D-B$ onto $P$ is $\mu-B_{PP}+O(s^{-1})$. Hence $D+B$ has exactly $|P|$
eigenvalues in a fixed disc containing the spectrum of $B_{PP}$, and they tend
to the eigenvalues of $B_{PP}$. By isospectrality these eigenvalues are
$s^{-1}\beta_\xi\to0$ ($\xi\in P$). So $B_{PP}$ is nilpotent. But $B_{PP}$ is the
convolution matrix of $b|_{\Hdiag{\varepsilon}}$ on the finite abelian group
$\Hdiag{\varepsilon}$, whose eigenvalues are the Fourier transforms
$\sum_{h\in\Hdiag{\varepsilon}}b_h\chi(h)$. All of them vanish, so
$b|_{\Hdiag{\varepsilon}}=0$.

(iv) Take $x=sa$ with generic $a$, so $D=s\diag(\alpha)+s^{-1}\diag(\beta)$ with
pairwise distinct $\alpha_\xi=\sum_ja_j\xi_j$. With $\epsilon=1/s$, the matrix
$(D+B)/s=\diag(\alpha)+\epsilon B+\epsilon^2\diag(\beta)$ has simple eigenvalue
branches analytic in $\epsilon$~\cite{R3-Kato}. The branch through $\alpha_1$ must
equal $\alpha_1+\epsilon^2\beta_1$, by isospectrality and continuity at
$\epsilon=0$. Comparing the $\epsilon^2$ Rayleigh--Schr\"odinger coefficients and
using (i) gives
\begin{equation}\label{R3:eq:second}
 \sum_{h\ne1}\frac{b_hb_{h^{-1}}}{L_h(a)}=0,\qquad L_h(a)=\sum_ja_j(1-h_j),
\end{equation}
identically in $a$. For $h,h'\in\Hcirc$ the vectors $(1-h_j)_j$ and $(1-h'_j)_j$
are proportional only if $h=h'$ or both $h,h'$ are constant vectors. In two
coordinates, write $h_k=e^{i\theta_k}$ with $\theta_k\in(0,2\pi)$; then
$(1-h_2)/(1-h_1)=e^{i(\theta_2-\theta_1)/2}\sin(\theta_2/2)/\sin(\theta_1/2)$
determines $\theta_2-\theta_1\in(-2\pi,2\pi)$. When $\theta_2\ne\theta_1$ it
determines $\theta_1$ as well, by strict monotonicity of $\cot(\theta_1/2)$.
Apply this to every pair of coordinates. Hence for non-constant $h\in\Hcirc$ the
hyperplane $L_h=0$ is a simple pole of~\eqref{R3:eq:second} carried by the single
term $h$. It is not the pole of any $h'$ with a coordinate equal to $1$, since
$L_{h'}$ then omits that variable. The residue gives $b_hb_{h^{-1}}=0$. A
constant $h\in\Hcirc$ lies in $\Hdiag{(1,\dots,1)}$, so $b_h=0$ by (iii).
\end{proof}

For $h\in H$ let $J(h)=\{j:h_j\neq1\}$. Call $h$ \emph{sign-constant} if there are
signs $\varepsilon_j$ ($j\in J(h)$) with $h_j^{\varepsilon_j}$ independent of
$j\in J(h)$.

\begin{corollary}[a rigidity criterion]\label{R3:cor:criterion}
Suppose $g_j(\Gamma)\le3$ for all $j$ and every $h\in H$ is sign-constant. Then $0$
is the only $\Gamma$-periodic potential Floquet isospectral to $0$. This applies
in particular to:
\begin{enumerate}
\item[(a)] every $\Gamma\supset3\Z^d$ or $\Gamma\supset2\Z^d$. This includes the theorem
for $(3\Z)^d$ in all dimensions (\cite[Part~XII]{K-034}), and gives a proof of it
that does not use the cubic perturbation identity employed there;
\item[(b)] for every $d\ge2$, $a\ge1$ and $\varepsilon\in\{\pm1\}^d$, the lattice
$\Gamma_{a,\varepsilon}=\{n\in\Z^d:\ \sum_j\varepsilon_jn_j\equiv0\pmod a\}$. This
family contains the lattices of~\cite[Theorem~6.1]{R3-FFORS} and the planar family
$(a,0)\Z\oplus(1,\pm1)\Z$.
\end{enumerate}
\end{corollary}

\begin{proof}
Let $h\in H$ and $J=J(h)$. Iterating Theorem~\ref{R3:thm:structure}(ii) over the
coordinates outside $J$ shows the following: the $|J|$-dimensional potential for
the projected lattice $p_J(\Gamma)\subset\Z^J$, with Fourier coefficients
$b|_{H_J}$ where $H_J=\{h':h'_i=1\ (i\notin J)\}$, is Floquet isospectral to $0$.
If $|J|=0$, then $b_h=0$ by (i). If $J=\{j\}$, this is a one-dimensional problem
of period $g=g_j\le3$, which is rigid (\cite[Theorem~C]{R3-CFK}; directly, for
$g\le3$ the characteristic polynomial at $z=1$ yields $e_1(V)=\dots=e_g(V)=0$,
hence $V=0$). So $b|_{K_j}=0$. If $|J|\ge2$, then $h$
lies in the $\varepsilon$-diagonal subgroup of the $|J|$-dimensional problem, and
part (iii) applied there gives $b_h=0$.

For (a), every coordinate of $h\in H\subset\mu_3^d$ (resp.\ $\mu_2^d$) different
from $1$ is $\omega^{\pm1}$ (resp.\ $-1$), and $g_j\mid3$ (resp.\ $2$).

For (b), $H=\{(\zeta^{\varepsilon_1},\dots,\zeta^{\varepsilon_d}):\zeta^a=1\}$, and
$e_1-\varepsilon_1\varepsilon_ke_k\in\Gamma_{a,\varepsilon}$ gives $g_j=1$.
\end{proof}

\begin{proposition}[reduction to a core statement]\label{R3:prop:reduction}
For $d\ge1$ let $\mathrm{CS}(d)$ be the statement: \emph{for every full-rank
$\Gamma\subset\Z^d$ with all $g_j\le3$ and every $\Gamma$-periodic $V$ that is Floquet
isospectral to $0$ and satisfies $\supp b\subset\Hcirc$, one has $V=0$.} If
$\mathrm{CS}(d')$ holds for all $d'\le d$, then
Conjecture~\ref{R3:conj:classification} holds in all dimensions $\le d$.
\end{proposition}

\begin{proof}
The ``if'' direction is Proposition~\ref{R3:prop:if}. For ``only if'', argue by
induction on $d$; the case $d=1$ is periods $\le3$. Let all $g_j\le3$ and
$h\notin\Hcirc$, so $J(h)\subsetneq\{1,\dots,d\}$. As in the proof of
Corollary~\ref{R3:cor:criterion}, the projected $|J(h)|$-dimensional potential is
isospectral to $0$. Its lattice has $g$-values $g_j\le3$ ($j\in J(h)$), so by the
induction hypothesis it vanishes, and $b_h=0$. Thus $\supp b\subset\Hcirc$ and
$\mathrm{CS}(d)$ gives $V=0$.
\end{proof}

\begin{remark}[a tool for the core statement]\label{R3:rem:minors}
Let $k$ be the least size of a nonzero principal minor of $B$; by (i) and (iv),
$k\ge3$ when $\supp b\subset\Hcirc$. Expanding
$\det(\lambda-D-B)=\sum_{M}\det(-B_{MM})\prod_{\xi\notin M}(\lambda-D_\xi)$,
dividing by $\prod_\xi(\lambda-D_\xi)$ and inserting $x$ as in (iii) with
$\lambda=s\mu$ shows the following. For every $\varepsilon$, the functions
$\sum_{|M|=k'}\det(B_{MM})\prod_{\xi\in M}(\mu-\alpha^\varepsilon_\xi(a))^{-1}$,
$k'\in\{k,k+1\}$, vanish identically in $(\mu,a)$; the next corrections occur at
relative order $s^{-2}$. When $\supp b=\{h\}$ is a single element this gives
$b_h^{\operatorname{ord}h}=0$. A general argument would need to control the
cancellations among these minors.
\end{remark}

\begin{proposition}[computer-assisted, $d=2$; partial]\label{R3:prop:computations}
\textup{(a) (exact.)} Let $\Gamma=(a,0)\Z\oplus(b,c)\Z$ with $0\le b<a$,
$\gcd(a,b)\le3$, $c\le3$ and $ac\le12$. Then $0$ is the only $\Gamma$-periodic
potential Floquet isospectral to $0$. This extends the symbolic check $ac\le8$
reported in~\cite{R3-FFORS}.

\textup{(b) (modular evidence, not a proof.)} For the same lattices with
$13\le ac\le14$, every branch ideal described below has the origin as its only zero
over $\overline{\mathbb F}_p$ for one prime $p\equiv1\pmod{ac}$ with
$p\approx10^6$.
\end{proposition}

\begin{proof}[Method]
The symmetries of $\Z^2$ preserve $\Delta$, isospectrality and the conditions
$g_1,g_2\le3$, so it suffices to treat one Hermite normal form in each symmetry
class. Fix $\Gamma$. By Theorem~\ref{R3:thm:structure}(i)--(iii) and the rigidity
of the one-dimensional periods $g_1=\gcd(a,b)\le3$ and $g_2=c\le3$, $b_h=0$
for every $h$ outside $A=\Hcirc\setminus(\Hdiag{(1,1)}\cup\Hdiag{(1,-1)})$.
If $A=\emptyset$, $\Gamma$ is rigid. Otherwise $A$ is a disjoint union of
inverse pairs $\{h,h^{-1}\}$. By (iv), $V$ lies in one of the
$2^{\#\text{pairs}}$ \emph{branch} subspaces, on which one chosen member of every
pair has $b=0$; let $S$ be the set of surviving members. For each branch, the
identity of Lemma~\ref{R3:lem:planewave} at $ac+8$ random rational points $x$ gives an
ideal $I_S\subset\Q(\zeta_{ac})[b_h:h\in S]$, and it suffices to show that its
only zero is the origin. Two symmetries reduce the work. If $\sigma_k$ is the
automorphism $\zeta_{ac}\mapsto\zeta_{ac}^k$, then $\sigma_k(I_S)=I_{S^k}$ after
renaming $b_h\mapsto b_{h^k}$, because the points $x$ are rational; so $I_S$ and
$I_{S^k}$ have the origin as only zero simultaneously. A lattice symmetry $T$ (a
signed permutation with $T\Gamma=\Gamma$) maps isospectral potentials of branch
$S$ to isospectral potentials of branch $T(S)$. Hence one branch per orbit
suffices; the orbits include $V\mapsto V(-\,\cdot)$, i.e.\ $k=-1$. For each
representative, \textsc{Singular}~\cite{R3-Singular} computes an exact standard
basis of $I_S$ over $\Q(\zeta_{ac})$. It checks that $\Q(\zeta_{ac})[b]/I_S$ is
finite-dimensional and that every variable is nilpotent modulo $I_S$. For index
$\le10$ it also checks that the total length equals the local length at the
origin, computed with a local ordering. Hence the origin is the only zero over $\C$,
and $V=0$. As
controls, the same code finds the known nonzero solutions for $4\Z\oplus\Z$ and
$5\Z\oplus\Z$ when the constraints on $K_1$ are dropped (lengths $24$ vs.\ $16$
and $120$ vs.\ $60$), and none for $3\Z\oplus\Z$. For (b) the same branch ideals,
built from $ac+8$ random points, are treated over $\mathbb F_p$, with
$\zeta_{ac}\in\mathbb F_p$. A solution over $\C$ could reduce to the origin
modulo an unlucky prime, so (b) is evidence only. The scripts are
\texttt{R3\_\allowbreak{}problem1\_\allowbreak{}exact\_\allowbreak{}branches.py} and \texttt{explore/branch2d\_\allowbreak{}fast.py}.
\end{proof}

\begin{table}[ht]
\centering\footnotesize
\begin{tabular}{lll}
\hline
$N=ac$ & rigid by & exact branch computation:\\
 & Cor.~\ref{R3:cor:criterion} & $[a,b,c]$ (orbits/branches; max.\ length)\\
\hline
$1$--$3$ & all & --\\
$4$ & $[2,0,2]$, $[4,1,1]$ & $[2,1,2]$ $(1/2;\,4)$\\
$5$ & $[5,1,1]$ & $[5,2,1]$ $(1/4;\,12)$\\
$6$ & $[6,1,1]$ & $[2,0,3]$ $(1/2;\,6)$, $[2,1,3]$ $(1/2;\,6)$,\\
 & & $[3,1,2]$ $(1/2;\,6)$\\
$7$ & $[7,1,1]$ & $[7,2,1]$ $(2/8;\,44)$\\
$8$ & $[4,2,2]$, $[8,1,1]$ & $[4,1,2]$ $(3/8;\,43)$, $[8,3,1]$ $(2/4;\,32)$\\
$9$ & $[3,0,3]$, $[9,1,1]$ & $[3,1,3]$ $(2/8;\,61)$, $[9,2,1]$ $(2/8;\,61)$\\
$10$ & $[10,1,1]$ & $[5,1,2]$ $(4/16;\,118)$, $[5,2,2]$ $(1/4;\,29)$,\\
 & & $[10,3,1]$ $(4/16;\,119)$\\
$11$ & $[11,1,1]$ & $[11,2,1]$ $(4/32;\,332)$, $[11,3,1]$ $(4/32;\,330)$\\
$12$ & $[12,1,1]$ & $[4,1,3]$ $(3/8;\,98)$, $[4,2,3]$ $(4/16;\,164)$,\\
 & & $[6,1,2]$ $(4/16;\,167)$, $[6,2,2]$ $(2/4;\,28)$, $[12,5,1]$ $(2/4;\,40)$\\
\hline
\end{tabular}
\caption{Exact verification of Proposition~\ref{R3:prop:computations}(a). Here
$[a,b,c]$ is the Hermite normal form $(a,0)\Z\oplus(b,c)\Z$, and one
representative is listed per class under the symmetries of $\Z^2$. There are
$37$ classes: $17$ are rigid by Corollary~\ref{R3:cor:criterion} (the allowed set
$A$ is empty) and $20$ are settled by the computation.
Lengths are $\dim_{\Q(\zeta_{ac})}$ of the branch quotient ring at $ac+8$ sample
points; for $N\le10$ they also equal the local length at the origin. Logs:
\texttt{R3\_\allowbreak{}problem1\_\allowbreak{}exact\_\allowbreak{}branches\_\allowbreak{}1\_\allowbreak{}4.log} and
\texttt{R3\_\allowbreak{}problem1\_\allowbreak{}exact\_\allowbreak{}branches.log}.}
\label{R3:tab:exact}
\end{table}

\subsection{Status of the two problems}\label{R3:ss:status}

\begin{itemize}
\item \cite[Conjecture~2]{R3-FFORS}: \textbf{partial}. Proved for every even $q>3$
(known) and every $q$ with a divisor in $[5,69]$
(Corollary~\ref{R3:cor:conj2}). It is exact for $q=5,9$ and computer-assisted
(Krawczyk) for the other odd $q\le69$. Open for $q=n$ and $q=3n$ with all prime
factors of $n$ larger than $69$, in particular for every prime $q>69$.
\item \cite[Problem~1]{R3-FFORS}: \textbf{partial}. The structure theorem holds
for all lattices and all $d$; rigidity is proved for the class of
Corollary~\ref{R3:cor:criterion}; the classification reduces to the core statement
$\mathrm{CS}(d)$ (Proposition~\ref{R3:prop:reduction}). The planar case is checked
exactly for index $\le12$ and modulo a prime for index $13$ and $14$. The core
statement itself remains open.
\end{itemize}

\begin{remark}[candidates not pursued]
Some neighbouring questions were considered and set aside. Liu asked whether
Fermi isospectrality is rigid for two-dimensional discrete periodic operators (see
also~\cite{R3-Liu23}); this was answered negatively by a numerically certified
counterexample~\cite{R3-BFL}, which also disproves a conjecture of Gieseker,
Kn\"orrer and Trubowitz. Conjectures~1 and~3 of~\cite{R3-FFORS}, and the
$(3\Z)^d$ case of the rectangular problem, are already treated in~\cite{K-034}.
\end{remark}


%% ==================================================================
%% Section: Spectral volumes
%% ==================================================================
% R4: spectral regions of weighted graph Laplacians -- edge addition and suspension
% (fragment, no preamble; uses amsmath/amssymb/amsthm environments)
\providecommand{\R}{\mathbb{R}}
\providecommand{\one}{\mathbf{1}}
\providecommand{\RG}{\mathcal{R}}
\providecommand{\tr}{\operatorname{tr}}
\providecommand{\spec}{\operatorname{spec}}
\providecommand{\Exp}{\operatorname{Exp}}

\section{Spectral volumes of weighted Laplacians under edge addition and suspension}
\label{R4:sec}

\noindent\textbf{Summary.}
\begin{itemize}
\item Problem~\ref{R4:prob:Q3} (EJT Question~3): \emph{partial}.  It is
proved for all connected graphs on at most six vertices, and in every
order for two infinite families.
\item Problem~\ref{R4:prob:Q4} (EJT Question~4): \emph{partial}.  It is
proved for all graphs on at most four vertices, together with a uniform
multiplicative bound.
\item EJT Question~5 and Fallat--Gupta--Lin's boundary-linearity question
were already answered in \cite{K-015}, and are not treated here.
\end{itemize}

%======================================================================
\subsection{Setting, the two target questions, and summary}
\label{R4:subsec:setting}

Let $G$ be a simple graph on $[n]=\{1,\dots,n\}$, $n\ge 2$.  For an edge
$f=ab$ put $b_f=e_a-e_b$ and $B_f=b_fb_f^{\mathsf T}$.  For
$w\in\R_{\ge0}^{E(G)}$ let $L_G(w)=\sum_{f\in E(G)}w_fB_f$.  Every such
matrix satisfies $L_G(w)\one=0$; we list its eigenvalues as
$0=\lambda_1\le\lambda_2\le\dots\le\lambda_n$ and call
$(\lambda_2,\dots,\lambda_n)$ its \emph{reduced spectrum} (zeros other than
$\lambda_1$ are allowed).  Write
\[
 C_{m}=\Bigl\{x\in\R^{m}:0\le x_1\le\dots\le x_{m},\ \textstyle\sum_i x_i=1\Bigr\},
 \qquad m=n-1,
\]
and define the (relaxed, trace-normalized, ordered) spectral region
\[
 \RG(G)=\bigl\{(\lambda_2,\dots,\lambda_n)(L_G(w)) : w\in\R_{\ge0}^{E(G)},\
 \tr L_G(w)=1\bigr\}\subseteq C_{n-1}.
\]
It is compact (continuous image of a compact simplex).  Following
Erchenko--Jakobson--Tsypin \cite[Section~8]{R4-EJT}, $P(G)$ is the
proportion of the standard simplex
$\{x\in\R^{n-1}_{\ge0}:\sum x_i=1\}$ occupied by realizable
(unordered) reduced spectra; since that set is the union of the $(n-1)!$
coordinate permutations of $\RG(G)$, which overlap only in a null set,
\begin{equation}\label{R4:eq:P}
 P(G)=\frac{\operatorname{vol}\RG(G)}{\operatorname{vol}C_{n-1}}
 =\Pr\Bigl[\tfrac{\operatorname{sort}(X)}{X_1+\dots+X_{n-1}}\in\RG(G)\Bigr],
 \qquad X_1,\dots,X_{n-1}\ \text{i.i.d. }\Exp(1),
\end{equation}
the last equality being the classical exponential representation of the
uniform distribution on a simplex.  If $G$ is disconnected then
$\lambda_2\equiv0$ on $\RG(G)$ and $P(G)=0$.

\begin{problem}[Erchenko--Jakobson--Tsypin {\cite[Section~8, Question~3]{R4-EJT}}]
\label{R4:prob:Q3}
``Consider a non-complete graph $G$. Let $G'$ be a graph obtained from $G$
by adding an edge. Is $P(G')>P(G)$?''
\end{problem}

\begin{problem}[Erchenko--Jakobson--Tsypin {\cite[Section~8, Question~4]{R4-EJT}}]
\label{R4:prob:Q4}
``Let $G$ be a non-complete graph and $G'$ be the suspension of $G$. Is
$P(G')\ge P(G)$?''  (The suspension adds one vertex joined to every vertex
of $G$, \cite[Definition~7.2]{R4-EJT}.)
\end{problem}

Taken literally, Problem~\ref{R4:prob:Q3} has a trivial negative answer
(add an edge to the empty graph on $n\ge3$ vertices: both volumes are
$0$), as already observed in \cite{K-015}; the meaningful question
concerns connected $G$.  For $n\le4$ the connected version follows from
the complete census of four-vertex volumes
\cite{R4-EJT,R4-FGL,K-015}; for every fixed $n\ge5$ it was open
(the concluding section of \cite{K-015} lists the positivity of the
fixed-order gain $c_n^{\rm conn}$, $n\ge5$, as open).  Problem~\ref{R4:prob:Q4} is trivial
for disconnected $G$, known along the family $K_n-e$
\cite{K-015}, and was open in general (EJT Question~5, on iterated
suspensions, is answered affirmatively in \cite{K-015}).

\medskip\noindent\textbf{Results of this section.}
\begin{itemize}
\item[(a)] A reduction: for connected $G$ and a non-edge $e$,
$P(G+e)>P(G)$ \emph{if and only if} $\RG(G+e)\neq\RG(G)$
(Proposition~\ref{R4:prop:reduction}); together with an exact,
finitely checkable linear (von Neumann) certificate of non-realizability
(Lemma~\ref{R4:lem:cert}).
\item[(b)] Problem~\ref{R4:prob:Q3} has a positive answer, in every order
$n$, for every connected $G$ whose complement minus the added edge is a
matching (possibly empty) (Theorem~\ref{R4:thm:matching}), and for every
connected $G$ for which adding $e$ increases the Hadwiger number
(Theorem~\ref{R4:thm:hadwiger}; e.g.\ all trees).
\item[(c)] Problem~\ref{R4:prob:Q3} has a positive answer for
\emph{every} connected non-complete graph with at most $6$ vertices
(Theorem~\ref{R4:thm:small}); this is computer-assisted, with exact
rational certificates and rigorous ball arithmetic.  In particular the
fixed-order constants $c^{\rm conn}_5$ and $c^{\rm conn}_6$ are
positive.  \textbf{Partial}: the general case $n\ge7$ remains open.
\item[(d)] For Problem~\ref{R4:prob:Q4}, we obtain the following.
\begin{itemize}
\item An exact ``lift identity'' for the uniform join
(Proposition~\ref{R4:prop:lift}).
\item The unconditional bound $P(G')\ge\gamma_n\,P(G)$, with an explicit
constant $\gamma_n\ge\frac{n+1}{2n}$ (Theorem~\ref{R4:thm:Q4partial}).
\item A positive answer for every graph with at most four vertices
(Theorem~\ref{R4:thm:Q4small}; computer-assisted: exact rational box
certificates for the lower bounds, based on the exact four-vertex regions
of \cite{R4-EJT}; the volumes of the graphs themselves are the closed forms
of \cite{K-015}, re-checked by quadrature, and the case $K_4-e$ uses the
$K_5-e$ region of \cite{K-015}).
\item A reduction of the full question to a purely probabilistic
``cone-lifting inequality'' (Conjecture~\ref{R4:conj:cone}), supported by
computation.
\end{itemize}
\textbf{Partial}: Problem~\ref{R4:prob:Q4} remains open for $n\ge5$
(except along $K_n-e$, \cite{K-015}).
\end{itemize}

%======================================================================
\subsection{Regular points and the reduction of Question 3}
\label{R4:subsec:reduction}

For $w\in\R^{E(G)}$ write
$\det(tI+L_G(w))=\sum_{k=0}^{n}q_k(w)t^{n-k}$, so
$q_k=e_k(\lambda_2,\dots,\lambda_n)$ and $q_n=0$.  By the all-minors
matrix-tree theorem \cite{R4-Chaiken},
\begin{equation}\label{R4:eq:forest}
 q_k(w)=\sum_{F}\Bigl(\prod_{C\in\pi(F)}|C|\Bigr)\prod_{f\in F}w_f ,
\end{equation}
the sum running over $k$-edge forests $F\subseteq E(G)$ and $\pi(F)$ being
the set of components of the spanning forest $([n],F)$.

\begin{lemma}[A regular weighting]\label{R4:lem:regular}
Let $H$ be a connected graph on $n\ge2$ vertices.  There is
$w\in\R_{>0}^{E(H)}$ such that $L_H(w)$ has $n-1$ distinct non-zero
eigenvalues and the Jacobian matrix
$\bigl(\partial\lambda_i/\partial w_f\bigr)_{2\le i\le n,\,f\in E(H)}$ has
rank $n-1$.
\end{lemma}

\begin{proof}
This is a weak form of the universal lacunary cone of \cite{K-016} (its
tree lacunary chart and regular spanning-subgraph lift).  That result gives,
for every connected graph, an open cone of strict realizations at regular
points.  We include the short argument only to keep the section
self-contained.  Choose a spanning tree $T$ and order its edges $f_1,\dots,f_{n-1}$ so that
every prefix $\{f_1,\dots,f_k\}$ spans a subtree on $k+1$ vertices (list
the vertices in breadth-first order from a root and let $f_k$ join the
$(k+1)$st vertex to its parent).  For $0<s<1$ put $w_{f_j}=s^{j}$ and
$w_f=0$ for $f\notin E(T)$.  Every $k$-subset of $E(T)$ is a forest, and
the prefix forest is the unique $k$-subset minimizing $\sum_{f_j\in F}j$;
its coefficient in \eqref{R4:eq:forest} is $k+1$.  Hence
$q_k=(k+1)s^{k(k+1)/2}(1+O(s))$, and in logarithmic coordinates
\[
 \frac{\partial\log q_k}{\partial\log w_{f_j}}
 =\begin{cases}1+O(s),&j\le k,\\ O(s),&j>k,\end{cases}
\]
because every term of $q_k$ containing $f_j$ with $j>k$ is $O(s)$ times
the leading term.  The matrix $(\mathbf 1[j\le k])_{k,j}$ is unitriangular,
so for small $s$ the map $(w_{f_j})_j\mapsto(q_k)_k$ is a local
diffeomorphism (in the original coordinates too, since
$\partial q_k/\partial w_j=(q_k/w_j)\,\partial\log q_k/\partial\log w_j$).
Moreover $q_k/q_{k-1}=\frac{k+1}{k}s^{k}(1+O(s))$, so by the Newton
polygon the reduced eigenvalues are $\frac{k+1}{k}s^{k}(1+O(s))$,
$1\le k\le n-1$, pairwise distinct for small $s$; near a polynomial with
simple roots, roots and coefficients are local diffeomorphic coordinates.
Thus the $(n-1)\times(n-1)$ minor of the eigenvalue Jacobian indexed by
$E(T)$ is non-singular.  Now give every edge of $E(H)\setminus E(T)$ the
weight $\varepsilon>0$; for $\varepsilon$ small, simplicity of the
spectrum and non-singularity of that minor persist by continuity
(eigenvalues and eigenvectors of a simple spectrum depend smoothly on the
matrix, and $\partial\lambda_i/\partial w_f=(b_f^{\mathsf T}x_i)^2$ for
unit eigenvectors $x_i$).  All weights are then positive.
\end{proof}

\begin{lemma}[Regions are closures of their interiors]\label{R4:lem:interior}
If $H$ is connected, then $\RG(H)$ is the closure of its interior
relative to the affine hull of $C_{n-1}$.
\end{lemma}

\begin{proof}
Let $\Lambda(w)=(\lambda_2,\dots,\lambda_n)(L_H(w))$ and let $Z\subset
\R^{E(H)}$ be the set where the reduced spectrum is not simple or the
coefficient map $w\mapsto(q_1,\dots,q_{n-1})$ has rank $<n-1$.  It is the
zero set of the product of the discriminant with the sum of squares of
the $(n-1)$-minors of the coefficient Jacobian, a polynomial, and it is
proper by Lemma~\ref{R4:lem:regular}.  It is also a cone ($q_k$ is
homogeneous of degree $k$), so it cannot contain the hyperplane
$\{\sum_f w_f=1/2\}$ (otherwise it would contain the open set
$\{\sum_fw_f>0\}$); hence the weights $w\in\Delta(H)=\{w\ge0,\ \sum
w_f=1/2\}$ (i.e.\ $\tr L=1$) with all $w_f>0$ and $w\notin Z$ are dense in
$\Delta(H)$.  At such $w$ the spectrum is simple and positive ($H$ is
connected) and $d\Lambda$ is onto $\R^{n-1}$.  Since $\sum_i\lambda_i=\tr
L_H(w)=2\sum_fw_f$, the differential of $\Lambda$ maps the tangent space
$\{\sum_f\dot w_f=0\}$ of $\{\sum_fw_f=1/2\}$ into $\{\sum_i\dot x_i=0\}$, and
it maps the remaining (radial) direction $w$ to $\Lambda(w)$, which has
coordinate sum $1$; hence the restriction of $\Lambda$ to
$\{w>0,\ \sum_fw_f=1/2\}$ is a submersion onto the affine hyperplane
$\{\sum_ix_i=1\}$ near $\Lambda(w)$.  As the image point has strictly increasing positive
coordinates, it is an interior point of $\RG(H)$ relative to the affine
hull of $C_{n-1}$; every point of $\RG(H)$ is a limit of such points.
\end{proof}

\begin{proposition}[Reduction]\label{R4:prop:reduction}
Let $G$ be a graph on $[n]$ and $e\notin E(G)$ such that $G+e$ is
connected.  Then $\RG(G)\subseteq\RG(G+e)$, and
\[
 P(G+e)>P(G)\iff \RG(G+e)\neq\RG(G).
\]
\end{proposition}

\begin{proof}
Inclusion: give $e$ weight zero.  If the two compact sets are equal the
volumes agree.  Conversely let $p\in\RG(G+e)\setminus\RG(G)$.  As
$\RG(G)$ is closed, some relatively open ball $U\ni p$ misses $\RG(G)$; by
Lemma~\ref{R4:lem:interior}, $U$ contains an interior point $p'$ of
$\RG(G+e)$, hence a small ball around $p'$ lies in
$\RG(G+e)\setminus\RG(G)$, which therefore has positive volume.
\end{proof}

Note that $p$ may be any point of $\RG(G+e)$, including degenerate
spectra with extra zeros.

\begin{lemma}[Linear certificate of non-realizability]\label{R4:lem:cert}
Let $Y$ be a real symmetric $n\times n$ matrix with
$b_f^{\mathsf T}Yb_f\le 0$ for every $f\in E(G)$, and let
$\theta_1\le\dots\le\theta_{n-1}$ be the eigenvalues of $HYH$ on $\one^\perp$,
$H=I-\frac1nJ$.  If $\sigma\in\R^{n-1}_{\ge0}$ satisfies
\begin{equation}\label{R4:eq:cert}
 \sum_{i=1}^{n-1}\theta_i\,\sigma_i^{\downarrow}>0
\end{equation}
($\sigma^\downarrow$ = decreasing rearrangement), then no
$L\in\{L_G(w):w\ge0\}$ has reduced spectrum $\sigma$.
\end{lemma}

\begin{proof}
If $L=L_G(w)$, $w\ge0$, has reduced spectrum $\sigma$, then, since
$L=HLH$ and $L\one=0$, von Neumann's trace inequality on $\one^\perp$
\cite[Ch.~III]{R4-Bhatia} gives
$\tr(YL)=\tr\bigl((HYH)L\bigr)\ge\sum_i\theta_i\sigma_i^\downarrow>0$, while
$\tr(YL)=\sum_f w_f\,b_f^{\mathsf T}Yb_f\le0$.
\end{proof}

The certificates of Lemma~\ref{R4:lem:cert} detect exactly the spectra
that do not majorize any spectrum of $\RG(G)$: by convexity of
$\{L_G(w):w\ge0\}$ and of the orbitope $\operatorname{conv}\{QD_\sigma Q^{\mathsf
T}\}$ and the separating hyperplane theorem, \eqref{R4:eq:cert} is
solvable iff $\sigma$ does not majorize the reduced spectrum of any
$L_G(w)$ of the same trace.  Equivalently they are dual to maximizing a
concave spectral functional $\sum_ic_i\lambda_{i+1}$
($c_1\ge\dots\ge c_{n-1}$) over trace-one weightings.

%======================================================================
\subsection{An all-order family: complements that become matchings}
\label{R4:subsec:matching}

\begin{lemma}[Missing-edge inequality and its equality case]\label{R4:lem:missing}
Let $L=L_G(w)$, $w\ge0$, with $\lambda_2>0$, and let $uv$ be a pair with
$L_{uv}=0$, $n\ge3$.  Then $\lambda_n\ge\frac{n}{n-2}\lambda_2$.  If
equality holds, then $e_u-e_v$ is a $\lambda_2$-eigenvector and
$L_{uk}=L_{vk}=-\lambda_n/n<0$ for every $k\notin\{u,v\}$.
\end{lemma}

\begin{proof}
(The inequality and its equality form are proved in \cite{K-015}
(the missing-edge inequality and the block form at equality); we include a
short argument for completeness.)  Put $m=\frac{\lambda_2+\lambda_n}2$,
$\delta=\frac{\lambda_n-\lambda_2}2$, $T=L-mH$, $p=He_u$, $q=He_v$.  On
$\one^\perp$, $\|T\|\le\delta$; $|p|=|q|=\sqrt{(n-1)/n}$ and
$p^{\mathsf T}Tq=L_{uv}-mH_{uv}=m/n$.  Cauchy--Schwarz gives
$m/n\le\delta\frac{n-1}{n}$, i.e.\ $n\lambda_2\le(n-2)\lambda_n$.  At
equality $Tq=\delta p$; writing $q=q_++q_-$ along the $\pm\delta$
eigenspaces of $T$ (forced by $|Tq|=\delta|q|$) we get $p=q_+-q_-$ and
$Tp=\delta q$.  Hence $L(p-q)=\lambda_2(p-q)$ and
$L(p+q)=\lambda_n(p+q)$ with $p-q=e_u-e_v$,
$p+q=e_u+e_v-\frac2n\one$.  Reading coordinate $k\notin\{u,v\}$ of both
identities gives $L_{ku}=L_{kv}$ and $L_{ku}+L_{kv}=-2\lambda_n/n$.
\end{proof}

\begin{theorem}[Positive answer when the new complement is a matching]
\label{R4:thm:matching}
Let $G$ be connected and non-complete on $n\ge3$ vertices, let
$F=E(\overline G)$ and $e\in F$.  If $F\setminus\{e\}$ is a matching
(possibly empty), then $P(G+e)>P(G)$.
\end{theorem}

\begin{proof}
By Proposition~\ref{R4:prop:reduction} it suffices to exhibit a point of
$\RG(G+e)\setminus\RG(G)$.

\emph{Case $F=\{e\}$.} The uniformly weighted $K_n=G+e$ has reduced
spectrum $(\frac1{n-1},\dots,\frac1{n-1})$.  A relaxed $G$-realization
would have $\lambda_2=\lambda_n>0$ and a zero entry at $e$, contradicting
Lemma~\ref{R4:lem:missing}.

\emph{Case $F\setminus\{e\}=M$, a matching with $r\ge1$ edges.} Give every
edge of $G+e=K_n-M$ weight $\alpha>0$.  Then
$L=\alpha\bigl(nH-\sum_{xy\in M}B_{xy}\bigr)$, whose reduced spectrum is
$\sigma_r=\{((n-2)\alpha)^{(r)},(n\alpha)^{(n-1-r)}\}$, attaining equality
in Lemma~\ref{R4:lem:missing}.  Suppose $L'=L_G(w')$, $w'\ge0$, had reduced
spectrum $\sigma_r$, and let $S$ be its positive support (a spanning
subgraph of $G$, connected as $\lambda_2>0$).  Every non-edge of $S$ is a
pair with $L'_{xy}=0$ and equality holds, so by Lemma~\ref{R4:lem:missing}
every vertex outside a non-edge $xy$ of $S$ is joined to both $x$ and $y$;
thus the non-edges of $S$ form a matching $M'$, and the vectors
$e_x-e_y$ ($xy\in M'$) are pairwise orthogonal $\lambda_2$-eigenvectors, so
$\operatorname{mult}(\lambda_2)\ge|M'|$.  Now $M'\supseteq F$.  If $F$ is not a
matching this is impossible.  If $F$ is a matching then
$|M'|\ge|F|=r+1>r=\operatorname{mult}_{\sigma_r}((n-2)\alpha)$, a
contradiction.
\end{proof}

\begin{remark}[Relation to {\cite{K-015}}]
The case $F=\{e\}$ is the last-edge gap computed exactly in
\cite{K-015} ($P(K_n)-P(K_n-e)=\binom{(n-2)(n+1)/2}{n-2}^{-1}$).  The
general case is a short consequence of the equality analysis of
\cite{K-015} (its ``support rigidity on the universal facet'': equality in
Lemma~\ref{R4:lem:missing} forces a matching complement and the facet
spectra $\sigma_r$ of $K_n-M_r$), combined with the multiplicity count in
the proof above; \cite{K-015} does not draw this edge-addition
consequence, so we record it.  The theorem covers all graphs $K_n-F$ in
which $F$ is a matching plus one edge, for every $n$.
\end{remark}

%======================================================================
\subsection{An all-order family on the sparse side: growth of the Hadwiger number}
\label{R4:subsec:hadwiger}

For a graph $H$ let $\operatorname{had}(H)$ be the largest $h$ for which
$V(H)$ can be partitioned into $h$ sets, each inducing a connected
subgraph, that are pairwise joined by at least one edge (for connected $H$
this is the Hadwiger number).

\begin{lemma}[Singular limits]\label{R4:lem:singular}
Let $H$ be a graph on $[n]$ and $w^{(k)}\in\R^{E(H)}_{\ge0}$ with
$w^{(k)}_f\to w^0_f\in[0,\infty]$ for every $f$.  Let
$\Gamma=\{f:w^0_f=\infty\}$, let $V_1,\dots,V_h$ be the vertex sets of the
components of $([n],\Gamma)$ (singletons allowed), $M=\operatorname{diag}(|V_i|)$,
and let $L_Q$ be the Laplacian on $[h]$ with the finite weights
$W_{ij}=\sum_{f\in E_H(V_i,V_j)}w^0_f$.  Let
$\lambda^{(k)}_1\le\dots\le\lambda^{(k)}_n$ be the eigenvalues of
$L_H(w^{(k)})$ and $\xi_1\le\dots\le\xi_h$ those of the pencil
$L_Qy=\xi My$.  Then $\lambda^{(k)}_j\to\infty$ for $j>h$ and
$\lambda^{(k)}_j\to\xi_j$ for $j\le h$.
\end{lemma}

\begin{proof}
Let $P$ be the space of vectors constant on every $V_i$ and $\iota:\R^h\to
P$ the obvious isomorphism; $|\iota y|^2=y^{\mathsf T}My$, and for $x=\iota
y$ the edges inside a $V_i$ contribute nothing, so
$x^{\mathsf T}L_H(w^{(k)})x\to y^{\mathsf T}L_Qy$ (edges between different
$V_i$ are not in $\Gamma$).  Courant--Fischer with the test spaces
$\iota(\text{first $j$ pencil eigenvectors})$ gives
$\limsup\lambda_j^{(k)}\le\xi_j$ for $j\le h$.  Let $L_\Gamma$ be the
unweighted Laplacian of $\Gamma$; its kernel is $P$, so $x^{\mathsf
T}L_\Gamma x\ge\kappa|x|^2$ on $P^\perp$ for some $\kappa>0$.  Every
$(h+1)$-dimensional subspace contains a unit $x\perp P$, and
$x^{\mathsf T}L_H(w^{(k)})x\ge(\min_{f\in\Gamma}w^{(k)}_f)\kappa\to\infty$,
whence $\lambda^{(k)}_{h+1}\to\infty$.  For the lower bound take
orthonormal eigenvectors $x^{(k)}_1,\dots,x^{(k)}_h$ and a subsequence
along which they converge to $x_1,\dots,x_h$.  Boundedness of
$\lambda^{(k)}_j\ge(\min_\Gamma w^{(k)})\,x_j^{(k)\mathsf T}L_\Gamma x^{(k)}_j$
forces $x_j\in P$.  For a unit $x=\sum_{i\le j}c_ix_i$ put
$x^{(k)}=\sum_{i\le j}c_ix^{(k)}_i$; then $\lambda_j^{(k)}\ge x^{(k)\mathsf
T}L_H(w^{(k)})x^{(k)}\ge\sum_{f\notin\Gamma}w^{(k)}_f(x^{(k)}_a-x^{(k)}_b)^2
\to y^{\mathsf T}L_Qy$ with $y=\iota^{-1}x$, so $\liminf\lambda^{(k)}_j$ is at
least the maximum of the pencil Rayleigh quotient on the $j$-dimensional
space $\iota^{-1}\operatorname{span}(x_1,\dots,x_j)$, which is $\ge\xi_j$.
As every subsequence has a further subsequence with these limits, the
whole sequence converges.
\end{proof}

\begin{theorem}[Hadwiger growth forces a strict volume gain]\label{R4:thm:hadwiger}
Let $G$ be connected and $e\notin E(G)$.  If
$\operatorname{had}(G+e)>\operatorname{had}(G)$, then $P(G+e)>P(G)$.
\end{theorem}

\begin{proof}
Let $h=\operatorname{had}(G+e)$ and fix a partition $V_1,\dots,V_h$ of
$[n]$ witnessing it in $G+e$, spanning trees $T_i$ of $(G+e)[V_i]$, and
for $i<j$ an edge $f_{ij}\in E(G+e)$ between $V_i$ and $V_j$.  Put
$m_i=|V_i|$.  For $t>0$ weight every edge of $\bigcup_iT_i$ by $t$, every
$f_{ij}$ by $m_im_j/n$, and all other edges by $0$.  By
Lemma~\ref{R4:lem:singular} (as $t\to\infty$) the quotient pencil is
$L_Q=M-mm^{\mathsf T}/n$, $M^{-1/2}L_QM^{-1/2}$ is the orthogonal
projection onto $(M^{1/2}\one)^\perp$, and so the reduced spectrum
$\sigma(t)$ has $h-1$ entries tending to $1$ and $n-h$ entries tending to
$\infty$.

Suppose $\sigma(t_k)$ were the reduced spectrum of $L_G(w^{(k)})$,
$w^{(k)}\ge0$, for a sequence $t_k\to\infty$.  Passing to a subsequence
with $w^{(k)}\to w^0\in[0,\infty]^{E(G)}$ and applying
Lemma~\ref{R4:lem:singular} to $G$: exactly $h$ eigenvalues stay bounded,
so $([n],\Gamma)$ has exactly $h$ components $V'_1,\dots,V'_h$, each
spanning a connected subgraph of $G$, and the quotient pencil
$(L_{Q'},M')$ has eigenvalues $0,1^{(h-1)}$.  Then
$M'^{-1/2}L_{Q'}M'^{-1/2}$ is the projection onto $(M'^{1/2}\one)^\perp$, so
$L_{Q'}=M'-m'm'^{\mathsf T}/n$ has all off-diagonal entries non-zero: every
two of the $V'_i$ are joined by an edge of $G$.  Thus
$\operatorname{had}(G)\ge h$, a contradiction.  Hence for some $t$ the
normalized $\sigma(t)$ lies in $\RG(G+e)\setminus\RG(G)$, and
Proposition~\ref{R4:prop:reduction} applies.
\end{proof}

\begin{remark}
Theorem~\ref{R4:thm:hadwiger} contains, for instance, every tree $G$ (adding
any edge creates a cycle, i.e.\ a $K_3$ minor), and every graph with no
$K_h$ minor to which $e$ adds one.  The same argument shows more
generally that $\RG(G+e)=\RG(G)$ would force, for every $h$, the union of
the quotient (``mass-minor'') regions of all $h$-part contractions of
$G+e$ to coincide with that of $G$; we use only the equal-spectrum point
of complete quotients.  The two families of
Theorems~\ref{R4:thm:matching} and~\ref{R4:thm:hadwiger} together cover
$21$ of the $42$ non-isomorphic pairs $(G,e)$ ($G$ connected, $e$ up to
$\operatorname{Aut}(G)$) with $n=5$, $131$ of the $402$ with $n=6$, and
$1348$ of the $5381$ with $n=7$; for $n\le6$ the remaining pairs are
handled by the certificates of Theorem~\ref{R4:thm:small}, while the
$4033$ remaining pairs with $n=7$ have not been certified.
\end{remark}

%======================================================================
\subsection{All connected graphs on at most six vertices}
\label{R4:subsec:small}

\begin{theorem}[Question 3 for $n\le6$; computer-assisted]\label{R4:thm:small}
Let $G$ be a connected non-complete graph on $n\le6$ vertices and
$e\notin E(G)$.  Then $P(G+e)>P(G)$.  Consequently, for $n\le 6$ the
fixed-order gain $c^{\rm conn}_n=\min\{P(G+e)-P(G)\}$, the minimum over
connected non-complete $G$ with $|V(G)|=n$ and $e\notin E(G)$, is
positive.
\end{theorem}

\begin{proof}
For $n=3$ the only pair is $P_3\subset K_3$, and for $n=4$ the six
connected graphs ($5$ non-complete, $6$ pairs $(G,e)$ up to isomorphism)
have the known volumes
$P(K_{1,3})\approx0.180$, $P(P_4)\approx0.384$, $P(\mathrm{paw})=9/13$,
$P(C_4)=3/4$, $P(K_4-e)=9/10$, $P(K_4)=1$
(the four-vertex census of \cite{K-015}, based on \cite{R4-EJT,R4-FGL}), so every edge
addition is strict.  (Our certificates below re-prove $n=3,4$ as well.)

For $n\in\{5,6\}$, for each of the connected non-complete graphs on $n$
vertices (20 for $n=5$, 111 for $n=6$, taken from the complete atlas of
graphs on at most seven vertices \cite{R4-Atlas}) and each orbit of
non-edges under $\operatorname{Aut}(G)$ ($42$ resp.\ $402$ pairs), the
file of certificates contains a pair $(w,Y)$
with: $w$ a vector of non-negative rationals on $E(G)\cup\{e\}$, so that
the reduced spectrum $\sigma$ of $L_{G+e}(w)$ (normalized) lies in
$\RG(G+e)$; and $Y$ a rational symmetric matrix with
$b_f^{\mathsf T}Yb_f\le0$ for all $f\in E(G)$ (checked exactly) and
\eqref{R4:eq:cert} (checked rigorously: the characteristic polynomials of
$L_{G+e}(w)$ and $HYH$ are computed exactly over $\mathbb Q$, divided by
$t$, their roots enclosed in certified balls (FLINT/Arb
\cite{R4-Arb}), and a lower bound
for $\min_\pi\sum_i\theta_i\sigma_{\pi(i)}=\sum_i\theta_i\sigma_i^\downarrow$ is
obtained from the ball lower bounds of all products, minimized over all
permutations).  By Lemma~\ref{R4:lem:cert}, $\sigma\notin\RG(G)$, and
Proposition~\ref{R4:prop:reduction} gives $P(G+e)>P(G)$.  Isomorphic
pairs have equal volumes, so all pairs are covered.  Finitely many pairs
give a positive minimum.
\end{proof}

\begin{remark}[The two classical extremal problems do not suffice]\label{R4:rem:dichotomy}
A natural sufficient condition for $P(G+e)>P(G)$ is that the
trace-normalized absolute algebraic connectivity $\max\lambda_2/\tr$
increases, or the trace-normalized minimal spectral radius
$\min\lambda_n/\tr$ decreases, when $e$ is added (each is a certificate of
the type in Lemma~\ref{R4:lem:cert}).  Numerically (to $10^{-5}$) this
holds for all $42$ pairs with $n=5$ and for $401$ of the $402$ pairs with
$n=6$.  It fails, provably, for the remaining pair: let $D=P_2\times P_3$
be the domino, with rows $1\,0\,5$ and $2\,3\,4$ (edges
$01,05,23,34$ and rungs $12,03,54$), and $e=14$ the long diagonal.  The
weights $\tfrac1{22}$ on the three rungs and $\tfrac1{11}$ on the four row
edges give $\tr L_D=1$ and spectrum $\{0,\tfrac1{11},\tfrac1{11},\tfrac2{11},
\tfrac3{11},\tfrac4{11}\}$, while the vector $x=(1,-1,-3,-1,1,3)$ (indexed
by $0,\dots,5$) satisfies $\sum x_i=0$, $|x|^2=22$ and $|x_a-x_b|=2$ on all
eight edges of $D+e$, so $\lambda_2\le\frac{x^{\mathsf T}Lx}{|x|^2}=\frac1{11}$
for every trace-one weighting of $D+e$.  Similarly the weights
$\tfrac19$ on the outer rungs and $\tfrac1{18}$ on the other five edges give
$\lambda_n(L_D)=\tfrac13$, while the bipartition vector
($+1$ on $\{0,2,4\}$, $-1$ on $\{1,3,5\}$; $e$ joins the two classes) gives
$\lambda_n\ge\frac13$ for every trace-one weighting of $D+e$.  So both
extremal values coincide for $D$ and $D+e$, yet $P(D+e)>P(D)$ by
Theorem~\ref{R4:thm:small}; its certificate excludes a $(D+e)$-spectrum
with $\lambda_2/\tr\approx\frac1{11}$ (within $10^{-8}$) but
$\lambda_3/\tr\approx\frac2{11}$.  (All of
this is checked in exact arithmetic by \texttt{R4\_\allowbreak{}domino.py}.)  Thus any
general proof of Problem~\ref{R4:prob:Q3} must use spectral information
beyond $\lambda_2$ and $\lambda_n$.
\end{remark}

%======================================================================
\subsection{Suspension: a lift identity and a partial answer to Question 4}
\label{R4:subsec:Q4}

\begin{lemma}[Uniform join {\cite[Lemma~7.3]{R4-EJT}}]\label{R4:lem:join}
If $L$ is a weighted Laplacian on $G$ (vertices $1,\dots,n$) with reduced
spectrum $\mu=(\mu_2,\dots,\mu_n)$ and the new vertex $0$ is joined to every
vertex with weight $t\ge0$, the resulting Laplacian on $G'=G\vee K_1$ has
reduced spectrum $\{\mu_i+t\}_{i=2}^n\cup\{(n+1)t\}$.
\end{lemma}

\begin{proof}
$L'(0,x)=(0,(L+tI)x)$ for $x\perp\one$, and
$L'(n,-\one)=(n+1)t\,(n,-\one)$.
\end{proof}

For $\nu\in\R^n_{\ge0}$ call an index $j$ \emph{admissible} if
$\nu_j\le(n+1)\min_i\nu_i$, and put
\[
 \mu^{(j)}(\nu)=\bigl(\nu_i-\tfrac{\nu_j}{n+1}\bigr)_{i\ne j}\in\R^{n-1}_{\ge0}.
\]
Admissibility is exactly $\mu^{(j)}\ge0$, and
$\nu=\{\mu^{(j)}_i+t\}\cup\{(n+1)t\}$ with $t=\nu_j/(n+1)$.  For a symmetric
cone $A\subseteq\R^{n-1}_{\ge0}$ let
$J(A)=\{\nu:\ \mu^{(j)}(\nu)\in A$ for some admissible $j\}$ and let
$U_{m}(A)$ denote the probability that $m$ i.i.d.\ $\Exp(1)$ variables
lie in $A$.  By Lemma~\ref{R4:lem:join} and \eqref{R4:eq:P},
\begin{equation}\label{R4:eq:Jbound}
 P(G\vee K_1)\ \ge\ U_n\bigl(J(\widehat{\RG}(G))\bigr),\qquad
 P(G)=U_{n-1}\bigl(\widehat{\RG}(G)\bigr),
\end{equation}
where $\widehat{\RG}(G)$ is the cone over all permutations of $\RG(G)$.

\begin{proposition}[Lift identity]\label{R4:prop:lift}
Let $\nu_1,\dots,\nu_n$ be i.i.d.\ $\Exp(1)$, $N(\nu)$ the number of
admissible indices and $N_A(\nu)$ the number of admissible $j$ with
$\mu^{(j)}\in A$.  Then $N$ is uniform on $\{1,\dots,n\}$ and, for every
measurable symmetric $A$ (more generally for every measurable symmetric
$f\ge0$),
\[
 \mathbb E\Bigl[\sum_{j\ \rm adm}f(\mu^{(j)})\Bigr]=\frac{n+1}{2}\,
 \mathbb E f(X),\qquad X\sim\Exp(1)^{\otimes(n-1)};
\]
in particular $\mathbb E N_A=\tfrac{n+1}2U_{n-1}(A)$.
More precisely, for each admissible $j$ the map
$(\nu)\mapsto(\mu^{(j)},t)$, $t=\nu_j/(n+1)$, has Jacobian $n+1$ and
transports the density $e^{-\sum\nu}$ to
$(n+1)e^{-\sum\mu}e^{-2nt}$.
\end{proposition}

\begin{proof}
By symmetry the left side is $n\,\mathbb E[f(\mu^{(n)})\mathbf
1\{n\ \text{adm}\}]$.  Substitute $\nu_i=\mu_i+t$ ($i<n$),
$\nu_n=(n+1)t$; the region $\{n$ admissible$\}$ becomes
$\{\mu\ge0,t\ge0\}$, the Jacobian is $n+1$ and
$\sum\nu=\sum\mu+2nt$.  Integrating $t$ gives the factor
$\frac{n+1}{2n}$.  For the law of $N$: with $a=\min\nu$ (so $e^{-na}$ is
uniform on $(0,1)$) the other $n-1$ variables exceed $a$ by i.i.d.\
$\Exp(1)$ amounts (R\'enyi's representation \cite{R4-Renyi}), so $N-1$ is
binomial$(n-1,1-e^{-na})$ mixed over a uniform parameter, i.e.\ uniform on
$\{0,\dots,n-1\}$.
\end{proof}

Because the transported density factorizes, a \emph{fixed} index is
exact: conditionally on $j$ being admissible (probability
$\frac{n+1}{2n}$), $\mu^{(j)}$ is exactly $\Exp(1)^{\otimes(n-1)}$ and
independent of $t$.  Geometrically, on the simplex of directions
$\Delta_{n-1}=\{\omega\ge0,\ \sum\omega_i=1\}$ the set where $j$ is admissible
is the pyramid with base the facet $\{\omega_j=0\}$ and apex
$p_j^*=\frac1{2n}(1,\dots,1,n+1,1,\dots,1)$ ($n+1$ in position $j$), and the
direction of $\mu^{(j)}$ is the central projection from $p_j^*$ onto that
facet, which pushes the uniform measure of the pyramid to the uniform
measure of the facet; the $n$ pyramids cover $\Delta_{n-1}$ with average
multiplicity $\frac{n+1}2$.

\begin{theorem}[Partial answer to Question 4]\label{R4:thm:Q4partial}
For every graph $G$ on $n\ge2$ vertices,
\[
 P(G\vee K_1)\ \ge\ \gamma_n\,P(G),\qquad
 \gamma_n:=\inf_{\theta}\ \frac{n+1}{2}\,
 \mathbb E_{r\sim\rho_n}\Bigl[\frac1{N(\theta,r)}\Bigr]\ \ge\ \frac{n+1}{2n},
\]
where $\theta$ ranges over the open simplex (directions of
$\mu\in\R^{n-1}_{>0}$), $\rho_n(r)=2(n-1)(1+2r)^{-n}$ on $[0,\infty)$, and
$N(\theta,r)$ is the number of admissible indices of the vector
$\nu(\theta,r)=\{\theta_i+r/n\}\cup\{(n+1)r/n\}$, namely
$N=1+\#\{i:\theta_i\le(n+2)r\}$ if $r<\min\theta$ and
$N=1+\#\{i:\theta_i\le(n+1)\min\theta+r\}$ otherwise.
\end{theorem}

\begin{proof}
Select, for each $\nu$, one admissible index uniformly at random.  Then
$\Pr[\nu\in J(A)]\ge\mathbb E\bigl[\sum_{j\ \rm adm}N^{-1}\mathbf
1_A(\mu^{(j)})\bigr]$.  By Proposition~\ref{R4:prop:lift} (with
$f=\mathbf 1_A\cdot N^{-1}$ read along the fibre) this equals
$\frac{n+1}{2}\mathbb E[\mathbf 1_A(X)\,\mathbb E_\tau N^{-1}]$ with
$\tau=nt\sim\Exp(2)$ independent of $X$; writing $X=S\theta$
($S\sim\mathrm{Gamma}(n-1)$, $\theta$ uniform and independent) and
$r=\tau/S$, whose density is $\rho_n$, the inner expectation depends only
on $(\theta,r)$, and the formula for $N$ follows from
$\min\nu=\min(\min\mu+\tau/n,(n+1)\tau/n)$.  Use \eqref{R4:eq:Jbound} with
$A=\widehat{\RG}(G)$.  The lower bound $\frac{n+1}{2n}$ is $N\le n$.
\end{proof}

\begin{remark}[Size of $\gamma_n$; numerical]\label{R4:rem:gamma}
The integrand is piecewise explicit, since
$\rho_n([a,\infty))=(1+2a)^{-(n-1)}$.  Numerically
(\texttt{R4\_\allowbreak{}q4\_\allowbreak{}numerics.py}) $\gamma_3\approx0.934$ (numerically the
minimum is at $\theta=(\frac16,\frac56)$, where the integrand equals
$\frac{239}{256}$),
$\gamma_4\approx0.826$, $\gamma_5\approx0.748$, $\gamma_6\approx0.700$,
$\gamma_7\approx0.667$, $\gamma_8\approx0.643$; for $4\le n\le8$ the
infimum is approached at a vertex of the simplex, where the integrand
tends to $\frac{n+1}2\bigl(\frac{1-3^{1-n}}{n-1}+\frac{3^{1-n}}n\bigr)$, so
presumably $\gamma_n\to\frac12$.  Only $\gamma_n\ge\frac{n+1}{2n}$ is
proved.  Since the $\theta$-average of the integrand is $1$
(Proposition~\ref{R4:prop:lift}), $\gamma_n\le1$ in any case: a uniformly
random choice of the removed index cannot answer
Problem~\ref{R4:prob:Q4}.  Nor can any rule that chooses the removed
entry by its rank among the admissible ones: a linear program over such
rules gives a best guaranteed density of only about $0.997$, $0.984$,
$0.979$, $0.977$ for $n=3,4,5,6$.
\end{remark}

For small $n$ the bound \eqref{R4:eq:Jbound} can be evaluated directly,
because the four-vertex regions are known exactly.

\begin{theorem}[Question~4 for at most four vertices; computer-assisted]
\label{R4:thm:Q4small}
Let $G$ be a non-complete graph on $n\le4$ vertices.  Then
$P(G\vee K_1)\ge P(G)$, with strict inequality when $G$ is connected.
The case $G=K_4-e$ uses the exact region of $K_5-e$ from \cite{K-015}.
\end{theorem}

\begin{proof}
If $G$ is disconnected, then $P(G)=0$.  The connected non-complete graphs
with $n\le4$ are $P_3$, $K_{1,3}$, $P_4$, the paw, $C_4$ and $K_4-e$.
Their suspensions satisfy $P_3\vee K_1=K_4-e$ and $(K_4-e)\vee K_1=K_5-e$.
For $P_3$: the paw is $K_4-e$ minus an edge, and setting that weight to $0$
shows that the region of the paw lies in that of $K_4-e$, so
$P(K_4-e)\ge P(\mathrm{paw})=\frac9{13}>\frac12=P(P_3)$ by
\cite[Theorem~6.1]{R4-EJT}; here $P(P_3)=\frac12$ is a direct computation, since
$\lambda_2\lambda_3=3ab$ and $\lambda_2+\lambda_3=2(a+b)$ give exactly
$\lambda_2\le\lambda_3/3$. For $K_4-e$: the exact regions of $K_4-e$ and
$K_5-e$ in \cite{K-015} give $P(K_5-e)=1-\binom93^{-1}=\frac{83}{84}>\frac9{10}=P(K_4-e)$.

For $G\in\{K_{1,3},P_4,\mathrm{paw},C_4\}$ we use $P(G\vee K_1)\ge
U_4\bigl(J(\widehat\RG(G))\bigr)$ from \eqref{R4:eq:Jbound}.  The regions
are known exactly for nonnegative weights \cite[Theorems~3.1, 4.1, 5.1,
6.1]{R4-EJT}.  For an ordered triple $0<a\le b\le c$ we use the following
sufficient conditions for membership.
\begin{itemize}
\item $K_{1,3}$: $F_a>0>F_b$, where $F_v=6v^2-2vs_v+p_v$, with $s_v$ and $p_v$ the sum and the product
of the other two weights.  EJT's condition is $F_aF_bF_c\le0$,
and $F_c\ge6c^2-4c^2+ab>0$.
\item $P_4$: $b>2a$ and $g=c(b-2a)-b(b-a)>0$.  EJT's conditions are
$c\ge\frac{a+b+c}2$ together with a sextic inequality $\mathrm{lhs}\le
\mathrm{rhs}$.  Sympy factors the difference as
$\mathrm{rhs}-\mathrm{lhs}=-256f_af_bf_c$, where
$f_v=v^2-vs_v+2p_v$ (notation as above).
Here $f_b=-g<0$, and $f_a\ge(a-b)^2+b^2>0$.  Also $g>0$ forces
$c>\frac{b(b-a)}{b-2a}>a+b$, which gives $c\ge\frac{a+b+c}2$ and
$f_c=c(c-a-b)+2ab>0$.
\item paw: $c>\kappa a$ with a rational $\kappa>2+\sqrt3$.  This gives
$\sqrt3(c-a)\ge a+c$.
\item $C_4$: $c>a+b$.
\end{itemize}

Write the order statistics of $\nu\sim\Exp(1)^{\otimes4}$ through
R\'enyi's representation \cite{R4-Renyi}: $s_1=\frac{E_1}4$,
$s_2-s_1=\frac{E_2}3$, $s_3-s_2=\frac{E_3}2$, $s_4-s_3=E_4$, with $E_i$
i.i.d.\ $\Exp(1)$.  The event $\nu\in J(\widehat\RG(G))$ is invariant under
scaling, so it is an event for the point $E/\sum_iE_i$, which is uniform on
the $3$-simplex.  In the coordinates $(e_2,e_3,e_4)$ of that point, the
four candidates $\mu^{(k)}$ (remove $s_k$ and subtract $s_k/5$) and their
admissibility conditions $s_k<5s_1$ are affine.

We cover the simplex by dyadic boxes.  A box is counted only if, for some
$k$, exact rational interval arithmetic proves on the whole box that $k$
is admissible and that $\mu^{(k)}$ satisfies the strict conditions above.
Its mass $6\operatorname{vol}(\text{box}\cap\text{simplex})$ is computed
exactly by inclusion--exclusion.  The total counted mass is therefore an
exact rational lower bound $\ell_G$ for $U_4(J(\widehat\RG(G)))$.

With boxes of side $2^{-8}$ (\texttt{R4\_\allowbreak{}q4\_\allowbreak{}small.py}), the bound $\ell_G$
exceeds the volume of $G$ in each case:
\begin{itemize}
\item $\ell_{K_{1,3}}=0.24663\ldots>0.17997\ldots=P(K_{1,3})$;
\item $\ell_{P_4}=0.52378\ldots>0.38350\ldots=P(P_4)$;
\item $\ell_{\mathrm{paw}}=0.77563\ldots>\frac9{13}$;
\item $\ell_{C_4}=0.83521\ldots>\frac34$.
\end{itemize}
The volumes of $G$ are the closed forms of the four-vertex census
\cite{K-015} (for the paw and $C_4$ also \cite{R4-EJT}), which we
re-computed independently as $P(G)=12\iint_\Omega(1+p+q)^{-3}\,dp\,dq$ by
quadrature; the margins exceed $0.06$ in every case.
\end{proof}

\begin{remark}
Monte Carlo estimates of $U_4(J(\widehat\RG(G)))$ are about $0.290$, $0.567$,
$0.785$ and $0.847$ for $K_{1,3}$, $P_4$, paw and $C_4$.  So the uniform
join alone exceeds $P(G)$ by $0.09$--$0.18$.  The same method gives
nothing for $n\ge5$, because the five-vertex regions are not known.
\end{remark}

\begin{conjecture}[Cone-lifting inequality]\label{R4:conj:cone}
For every $n\ge3$ and every closed symmetric cone
$A\subseteq\R^{n-1}_{\ge0}$, $\ U_n(J(A))\ \ge\ U_{n-1}(A)$.
\end{conjecture}

By \eqref{R4:eq:Jbound}, Conjecture~\ref{R4:conj:cone} implies a positive
answer to Problem~\ref{R4:prob:Q4} for every $G$.  By Strassen's theorem
\cite{R4-Strassen} it is equivalent, up to closedness issues at the
apices where some $\mu^{(j)}$ vanishes, to the existence of a coupling of a uniform direction
$\theta'$ in the $(n-1)$-simplex and a uniform direction $\theta$ in the
$(n-2)$-simplex (up to permutation) with $\theta'\in J(\R_{>0}\theta)$
almost surely, i.e.\ to a randomized choice, for each point of
$\Delta_{n-1}$, of one pyramid containing it such that the central
projections described after Proposition~\ref{R4:prop:lift} produce exactly
the uniform law on the facet; by
Proposition~\ref{R4:prop:lift} the ``multiplicity budget'' is exactly
$\frac{n+1}2$ per target point, so the question is whether the overlaps
can be arranged without loss.  Two facts are provable: (i) choosing the
minimal coordinate as the removed index gives, by
Proposition~\ref{R4:prop:lift}, $\Pr[\mu^{(\arg\min)}\in A]=\frac{n+1}{2}
\mathbb E[\mathbf 1_A(X)(1-e^{-2\min X})]$, and since
$\mu^{(\arg\min)}=X'+\frac{n\min\nu}{n+1}\one$ with $X'$ an exact
$\Exp(1)^{\otimes(n-1)}$ sample, the conjecture holds for every cone that
is closed under $x\mapsto x+c\one$ ($c\ge0$); (ii) the conjecture holds for
$A$ equal to the whole orthant (with equality).  Regions of non-complete
graphs are, however, never closed under adding constants: by
Lemma~\ref{R4:lem:missing} they omit a neighbourhood of the
equal-spectrum ray.  So neither fact settles Problem~\ref{R4:prob:Q4}.
The cone-lifting inequality has been tested numerically
(\texttt{R4\_\allowbreak{}q4\_\allowbreak{}numerics.py}).  Monte Carlo estimates of both sides for
sixteen families of cones and $n\in\{3,4,5,6,8\}$ give ratios $\ge1$ up
to sampling error; the tightest cases are the graph-like cones
$\{\max/\min\ge\kappa\}$.  A discretized Hall/max-flow test shows no
deficit for $n=3,4,5$.  What remains open is exactly
Conjecture~\ref{R4:conj:cone}, or any graph-specific substitute for it.

%======================================================================
\subsection{Computations accompanying this section}
\label{R4:subsec:computations}

All scripts are run with \texttt{python3 -I}, and each prints what it
checks.

\emph{Used in proofs.}
\begin{itemize}
\item \texttt{q3\_\allowbreak{}verify.py}: the independent verifier for
Theorem~\ref{R4:thm:small}.  It works in exact $\mathbb Q$-arithmetic with
Arb balls.  For $N=3,\dots,6$ it checks the file
\texttt{q3\_\allowbreak{}certs\_\allowbreak{}n$N$.json}, reporting $1$, $6$, $42$ and $402$
valid certificates, with no pair uncovered.  The smallest certified margins
are $0.316$, $1.08\cdot10^{-3}$, $1.84\cdot10^{-2}$ and $9.9\cdot10^{-6}$.
\item \texttt{R4\_\allowbreak{}q4\_\allowbreak{}small.py}: Theorem~\ref{R4:thm:Q4small}, using
exact rational interval arithmetic and exact box masses.
\item \texttt{R4\_\allowbreak{}domino.py}: Remark~\ref{R4:rem:dichotomy}, in exact
arithmetic.
\end{itemize}

\emph{Not used in proofs.}
\begin{itemize}
\item The certificate generators \texttt{q3\_\allowbreak{}certify*.py}.
\item The sanity checks \texttt{R4\_\allowbreak{}checks.py} and
\texttt{R4\_\allowbreak{}regions\_\allowbreak{}check.py}.
\item The numerical evidence \texttt{R4\_\allowbreak{}q4\_\allowbreak{}numerics.py} for
Remark~\ref{R4:rem:gamma} and Conjecture~\ref{R4:conj:cone}.
\end{itemize}


\begin{thebibliography}{999}

\bibitem{K-006} Keio, \emph{Minimal recurrence orders and integral solution lattices for Franel sequences}, manuscript, September 2026.

\bibitem{K-007} Keio, \emph{Dyadic Lucas prime cascades and composite Fibonacci shifts}, manuscript, September 2026.

\bibitem{K-015} Keio, \emph{One missing edge: the exact $K_n-e$ cone, strict support, and sharp defect laws}, manuscript, September 2026.

\bibitem{K-016} Keio, \emph{Universal weighted-Laplacian cones and the maximal-cusp spectral link}, manuscript, September 2026.

\bibitem{K-024} Keio, \emph{Transcendence of algebraic arc lengths and inverse periods: division dichotomies, beta quantiles, and arithmetic transfers}, manuscript, October 2026.

\bibitem{K-029} Keio, \emph{Observation and recovery: geometry, mixtures, and spectral inversion}, unified manuscript, October 2026. Part~VI, section ``The complete binomial-mixture likelihood transition''; its results are cited here by their labels in that manuscript.

\bibitem{K-031} Keio, \emph{Prime-representing constants and related structures}, manuscript, October 2026.

\bibitem{K-032} Keio, \emph{Minimal differential and recurrence orders: Bessel periods, bounded-height tableaux, and banana integrals}, manuscript, October 2026.

\bibitem{K-033} Keio, \emph{Euler's constant: arithmetic, monodromy, and Gamma-related structures}, collected manuscript, October 2026.

\bibitem{K-034} Keio, \emph{Global proof transfers and structural classifications}, collected manuscript, October 2026.

\bibitem{K-037} Keio, \emph{Isotropic Fuchs squeezes and exact differential complexity: hypercubic Green functions, Regev sums, and random walk poles}, manuscript, October 2026.

\bibitem{K-041} Keio, \emph{Ewald points of monotone polytopes: the minimum in dimension eight and branching simplex bundles}, manuscript, October 2026.

\bibitem{K-042} Keio, \emph{Cubic Pisot floors: composite windows, sparse exceptions, and finite-sieve barriers}, manuscript, October 2026.

\bibitem{K-047} Keio, \emph{A monic resultant for Bessel derivative coefficients}, manuscript, October 2026.

\bibitem{R1-KZ}
M.~Kauers and D.~Zeilberger,
\emph{The computational challenge of enumerating high-dimensional rook walks},
Adv. in Appl. Math. \textbf{47} (2011), no.~4, 813--819; arXiv:1011.4671.

\bibitem{R1-CK}
S.~Chen and M.~Kauers,
\emph{Some open problems related to creative telescoping},
J. Syst. Sci. Complex. \textbf{30} (2017), no.~1, 154--172; arXiv:1609.03768.

\bibitem{R1-BCvHP}
A.~Bostan, F.~Chyzak, M.~van~Hoeij and L.~Pech,
\emph{Explicit formula for the generating series of diagonal 3D rook paths},
S\'em. Lothar. Combin. \textbf{66} (2011), Art.~B66a; arXiv:1105.4456.

\bibitem{R1-Lipshitz}
L.~Lipshitz,
\emph{The diagonal of a D-finite power series is D-finite},
J. Algebra \textbf{113} (1988), 373--378.

\bibitem{R1-Malgrange}
B.~Malgrange,
\emph{Sur les points singuliers des \'equations diff\'erentielles},
Enseign. Math. (2) \textbf{20} (1974), 147--176.

\bibitem{R1-DLMF}
\emph{NIST Digital Library of Mathematical Functions}, \texttt{https://dlmf.nist.gov}, \S15.8
(release 1.2.x).

\bibitem{R1-OEIS}
OEIS Foundation, \emph{The On-Line Encyclopedia of Integer Sequences}, sequences
A051708, A144045, A181731, \texttt{https://oeis.org}.

\bibitem{R1-EG}
S.~Even and J.~Gillis,
\emph{Derangements and Laguerre polynomials},
Math. Proc. Cambridge Philos. Soc. \textbf{79} (1976), 135--143.

\bibitem{R1-KZfiles}
M.~Kauers and D.~Zeilberger, recurrences for $d$-dimensional rook walks, data files at
\url{http://www.risc.jku.at/people/mkauers/walks/} and
\url{http://sites.math.rutgers.edu/~zeilberg/tokhniot/ManuelRookRecurrences} (accessed October 2026).

\bibitem{R7-MRS16}
N.~Mani and S.~Rubinstein-Salzedo,
\emph{$N$-division points of hypocycloids},
Arnold Math.\ J.\ \textbf{2} (2016), no.~2, 149--158.
DOI 10.1007/s40598-015-0034-2; arXiv:1503.00566 (2015)
[Remark~5.7 of the published version is Remark~4.3 of the arXiv version].

\bibitem{R7-CoxShurman05}
D.~A.~Cox and J.~Shurman,
\emph{Geometry and number theory on clovers},
Amer.\ Math.\ Monthly \textbf{112} (2005), no.~8, 682--704.

\bibitem{R7-CoxHyde14}
D.~A.~Cox and T.~Hyde,
\emph{The Galois theory of the lemniscate},
J.\ Number Theory \textbf{135} (2014) (online 2013), Corollary~4.8; arXiv:1208.2653.

\bibitem{R7-Alperin00}
R.~C.~Alperin,
\emph{A mathematical theory of origami constructions and numbers},
New York J.\ Math.\ \textbf{6} (2000), 119--133; arXiv:math/9912039.

\bibitem{R7-CoxGalois}
D.~A.~Cox,
\emph{Galois Theory}, 2nd ed.,
Wiley, Hoboken, NJ, 2012.

\bibitem{R7-JS13}
W.~Jedidi and T.~Simon,
\emph{Further examples of GGC and HCM densities},
Bernoulli \textbf{19} (2013), no.~5A, 1818--1838; arXiv:1312.2809.

\bibitem{R7-Bosch15}
P.~Bosch,
\emph{HCM property and the half-Cauchy distribution},
Probab.\ Math.\ Statist.\ \textbf{35} (2015), no.~2, 191--200; arXiv:1402.1059.

\bibitem{R7-BoschSimon15}
P.~Bosch and T.~Simon,
\emph{On the infinite divisibility of inverse Beta distributions},
Bernoulli \textbf{21} (2015), no.~4, 2552--2568; arXiv:1405.4176.

\bibitem{R7-BoschSimon16}
P.~Bosch and T.~Simon,
\emph{A proof of Bondesson's conjecture on stable densities},
Ark.\ Mat.\ \textbf{54} (2016), no.~1, 31--38; arXiv:1411.3369.

\bibitem{R7-Simon11}
T.~Simon,
\emph{Multiplicative strong unimodality for positive stable laws},
Proc.\ Amer.\ Math.\ Soc.\ \textbf{139} (2011), no.~7, 2587--2595.

\bibitem{R7-Bondesson92}
L.~Bondesson,
\emph{Generalized Gamma Convolutions and Related Classes of Distributions and
Densities}, Lecture Notes in Statistics \textbf{76}, Springer, New York, 1992.

\bibitem{R7-BondessonSimon18}
L.~Bondesson and T.~Simon,
\emph{Stieltjes functions of finite order and hyperbolic monotonicity},
Trans.\ Amer.\ Math.\ Soc.\ \textbf{370} (2018), no.~6, 4201--4222; arXiv:1604.05267.

\bibitem{R7-Widder41}
D.~V.~Widder,
\emph{The Laplace Transform},
Princeton Univ.\ Press, Princeton, NJ, 1941.

\bibitem{R7-DLMF}
F.~W.~J.~Olver et al.\ (eds.),
\emph{NIST Digital Library of Mathematical Functions},
\texttt{https://dlmf.nist.gov}, Chapters~10, 13 and~16.

\bibitem{R7-LM95}
L.~Lorch and M.~E.~Muldoon,
\emph{Transcendentality of zeros of higher derivatives of functions involving
Bessel functions},
Internat.\ J.\ Math.\ Math.\ Sci.\ \textbf{18} (1995), no.~3, 551--560.

\bibitem{R7-Videla}
C.~R.~Videla, On points constructible from conics, \emph{Math. Intelligencer} \textbf{19} (1997), no.~2, 53--57.

\bibitem{R7-BCJLS}
S.~Baek, I.~Choe, Y.~Jung, D.~Lee and J.~Seo, Constructions by ruler and compass, together with a fixed conic, \emph{Bull. Aust. Math. Soc.}; arXiv:1210.8046 (2012).

\bibitem{R5-VKVK19}
R.~K. Vinayak, W.~Kong, G.~Valiant and S.~Kakade,
\emph{Maximum likelihood estimation for learning populations of parameters},
Proceedings of the 36th International Conference on Machine Learning (ICML 2019),
PMLR \textbf{97} (2019); arXiv:1902.04553.
Problem~\ref{R5:prob:wp} is quoted from Section~5 (``Discussion and future directions'').

\bibitem{R5-PW20}
Y.~Polyanskiy and Y.~Wu,
\emph{Self-regularizing property of nonparametric maximum likelihood estimator in mixture models},
arXiv:2008.08244v2 (2020).
Conjecture~\ref{R5:conj:gompertz} is from Section~5.2 (Example~3 and the paragraph after eq.~(37));
Problem~\ref{R5:prob:modes} is from Section~5.5; Lemma~9 and Corollaries~10--11 are in Section~5.4--5.5.

\bibitem{R5-MMRS25}
C.~Musco, C.~Musco, L.~Rosenblatt and A.~V. Singh,
\emph{Sharper bounds for Chebyshev moment matching, with applications},
Proceedings of the 38th Conference on Learning Theory (COLT 2025), PMLR \textbf{291} (2025);
arXiv:2408.12385v3 (May 2026).
The quoted sentences are footnote~3 (after Theorem~1) and the paragraph after Algorithm~1 in Section~3.1.

\bibitem{R5-BKM22}
V.~Braverman, A.~Krishnan and C.~Musco,
\emph{Sublinear time spectral density estimation},
Proceedings of the 54th ACM Symposium on Theory of Computing (STOC 2022);
arXiv:2104.03461 (full version with appendices: ``Linear and sublinear time spectral density estimation'').

\bibitem{R5-Favard37}
J.~Favard,
\emph{Sur les meilleurs proc\'ed\'es d'approximation de certaines classes de fonctions par des polyn\^omes trigonom\'etriques},
Bulletin des Sciences Math\'ematiques \textbf{61} (1937), 209--224.

\bibitem{R5-AK37}
N.~I. Akhiezer and M.~G. Krein,
\emph{On the best approximation of differentiable periodic functions by trigonometric sums},
Doklady Akademii Nauk SSSR \textbf{15} (1937) (in Russian).

\bibitem{R5-Korneichuk}
N.~Korneichuk,
\emph{Exact Constants in Approximation Theory},
Encyclopedia of Mathematics and its Applications \textbf{38}, Cambridge University Press, 1991.

\bibitem{R5-KK20}
N.~Kashyap and M.~Krishnapur,
\emph{How many modes can a mixture of Gaussians with uniformly bounded means have?},
arXiv:2005.01580 (2020).

\bibitem{R5-MWY24}
Y.~Ma, Y.~Wu and P.~Yang,
\emph{On the best approximation by finite Gaussian mixtures},
arXiv:2404.08913v2 (2025).

\bibitem{R5-SGS25}
J.~A. Soloff, A.~Guntuboyina and B.~Sen,
\emph{Multivariate, heteroscedastic empirical Bayes via nonparametric maximum likelihood},
Journal of the Royal Statistical Society Series B \textbf{87} (2025), 1--32.

\bibitem{R5-Wang26}
H.~Wang,
\emph{Almost-sure uniqueness of the Gaussian location NPMLE},
arXiv:2609.23903 (2026).

\bibitem{R5-LBK25}
J.~Lee, V.~Ba\'cak and E.~H. Kennedy,
\emph{Learning smooth populations of parameters with trial heterogeneity},
arXiv:2507.23140v2 (2026).

\bibitem{R5-LK26}
J.~Lee and E.~H. Kennedy,
\emph{Minimax rates for learning smooth populations of parameters},
arXiv:2610.04135v2 (2026).

\bibitem{R2-QiCR2022}
F.~Qi, Decreasing properties of two ratios defined by three and four polygamma functions,
\emph{C. R. Math. Acad. Sci. Paris} \textbf{360} (2022), 89--101.
doi:10.5802/crmath.296.

\bibitem{R2-YHQ2024}
H.-P.~Yin, L.-X.~Han, F.~Qi, Decreasing and complete monotonicity of functions defined by
derivatives of completely monotonic function involving trigamma function,
\emph{Demonstr. Math.} \textbf{57} (2024), Art.~20240041.
doi:10.1515/dema-2024-0041; arXiv:2405.19361.

\bibitem{R2-QiAgarwal2019}
F.~Qi, R.~P.~Agarwal, On complete monotonicity for several classes of functions related to
ratios of gamma functions, \emph{J. Inequal. Appl.} \textbf{2019}, Paper No.~36.
doi:10.1186/s13660-019-1976-z.

\bibitem{R2-QiAgarwal2024}
F.~Qi, R.~P.~Agarwal, Several functions originating from Fisher--Rao geometry of Dirichlet
distributions and involving polygamma functions, \emph{Mathematics} \textbf{12} (2024), no.~1,
Art.~44. doi:10.3390/math12010044.

\bibitem{R2-QiRM2021}
F.~Qi, Lower bound of sectional curvature of Fisher--Rao manifold of beta distributions and
complete monotonicity of functions involving polygamma functions, \emph{Results Math.}
\textbf{76} (2021), no.~4, Paper No.~217. doi:10.1007/s00025-021-01530-2
(preprint: doi:10.20944/preprints202011.0315.v1).

\bibitem{R2-QiMIA2021}
F.~Qi, Necessary and sufficient conditions for a difference constituted by four derivatives of
a function involving trigamma function to be completely monotonic, \emph{Math. Inequal. Appl.}
\textbf{24} (2021), no.~3, 845--855. doi:10.7153/mia-2021-24-58.

\bibitem{R2-QiAIMS2020}
F.~Qi, Completely monotonic degree of a function involving trigamma and tetragamma functions,
\emph{AIMS Math.} \textbf{5} (2020), no.~4, 3391--3407. doi:10.3934/math.2020219.

\bibitem{R2-Yang2025}
Z.-H.~Yang, Extension of a complete monotonicity theorem with applications,
arXiv:2507.10954 (2025).

\bibitem{R2-Qi2026}
F.~Qi, The decrease in sectional curvature of the Fisher--Rao manifold of beta distributions,
arXiv:2609.10628 (2026).

\bibitem{R2-Pudim}
Pudim Project, public research archive \texttt{pudim-ai-demo-zetalaw} (GitHub), ledger entries
APP-0020 and APP-0041 (Qi's degree-four conjecture for $h_\lambda$), accessed October 2026.

\bibitem{R2-Widder}
D.~V.~Widder, \emph{The Laplace Transform}, Princeton University Press, Princeton, 1941.

\bibitem{R2-DLMF}
NIST Digital Library of Mathematical Functions, \texttt{https://dlmf.nist.gov}, \S5.15 and \S25.11.

\bibitem{R6-HIKK}
S.~Hahn, D.~Ismailescu, G.~Kim and M.~Kim,
\emph{Explicit algebraic numbers all whose integer parts of powers are composite},
arXiv:2608.05309v1 (August 2026).

\bibitem{R6-Siurys}
J.~\v{S}iurys,
\emph{A linear recurrence sequence of composite numbers},
LMS Journal of Computation and Mathematics \textbf{15} (2012), 360--373.

\bibitem{R6-SiurysThesis}
J.~\v{S}iurys,
\emph{Linear recurrence sequences of composite numbers} (Tiesin\.es rekurenčiosios sekos sudarytos i\v{s} sud\.etini\c{u} skai\v{c}i\c{u}),
doctoral dissertation, Vilnius University, 2013.

\bibitem{R6-Graham}
R.~L.~Graham,
\emph{A Fibonacci-like sequence of composite numbers},
Mathematics Magazine \textbf{37} (1964), 322--324.

\bibitem{R6-Lunev}
I.~Lunev,
\emph{A Tribonacci-like sequence of composite numbers},
Journal of Integer Sequences \textbf{20} (2017).

\bibitem{R6-IKLP}
D.~Ismailescu, A.~Ko, C.~Lee and J.~Y.~Park,
\emph{On second order linear sequences of composite numbers},
arXiv:1812.08041 (2018).

\bibitem{R6-Zsigmondy}
K.~Zsigmondy,
\emph{Zur Theorie der Potenzreste},
Monatshefte f\"ur Mathematik und Physik \textbf{3} (1892), 265--284;
see also A.~S.~Bang, \emph{Taltheoretiske Unders\o gelser}, Tidsskrift for Mathematik (5) \textbf{4} (1886), 70--80, 130--137.

\bibitem{R6-Lucas}
E.~Lucas,
\emph{Th\'eorie des fonctions num\'eriques simplement p\'eriodiques},
American Journal of Mathematics \textbf{1} (1878), 184--240, 289--321.

\bibitem{R6-Varona}
J.~L.~Varona,
\emph{A couple of transcendental prime-representing constants},
American Mathematical Monthly (2022), doi:10.1080/00029890.2021.1977885; arXiv:2012.11750.

\bibitem{R6-Ridout}
D.~Ridout,
\emph{Rational approximations to algebraic numbers},
Mathematika \textbf{4} (1957), 125--131.

\bibitem{R6-Saito}
K.~Saito,
\emph{Intervals without primes near an iterated linear recurrence sequence},
arXiv:2504.14968 (2025); International Journal of Number Theory (2026).

\bibitem{R6-HR}
G.~H.~Hardy and S.~Ramanujan,
\emph{The normal number of prime factors of a number $n$},
Quarterly Journal of Mathematics \textbf{48} (1917), 76--92.

\bibitem{R3-FFORS}
M.~Faust, L.~Friedman, G.~O'Malley, R.~Ramos and A.~Sharma,
\emph{Algebraic properties of the ideal of spectral invariants for the discrete Laplacian},
arXiv:2602.08304v1 (2026).

\bibitem{R3-CFK}
J.~Cobb, M.~Faust and A.~Kretschmer,
\emph{Inverse eigenvalue problems, Floquet isospectrality and the Hilbert--Chow morphism},
arXiv:2601.09604v1 (2026).

\bibitem{R3-FLMPRTTZ}
M.~Faust, W.~Liu, R.~Matos, J.~Plute, J.~Robinson, Y.~Tao, E.~Tran and C.~Zhuang,
\emph{Floquet isospectrality of the zero potential for discrete periodic Schr\"odinger operators},
J.~Math.~Phys.~\textbf{65} (2024), 073501; arXiv:2401.09731.

\bibitem{R3-Liu23}
W.~Liu,
\emph{Floquet isospectrality for periodic graph operators},
arXiv:2302.13103 (2023).

\bibitem{R3-BFL}
T.~Brysiewicz, M.~Faust and W.~Liu,
\emph{A counterexample to Fermi isospectral rigidity for two dimensional discrete periodic Schr\"odinger operators},
arXiv:2603.11183 (2026).

\bibitem{R3-Krawczyk}
R.~Krawczyk,
\emph{Newton-Algorithmen zur Bestimmung von Nullstellen mit Fehlerschranken},
Computing \textbf{4} (1969), 187--201.

\bibitem{R3-Moore}
R.~E.~Moore,
\emph{A test for existence of solutions to nonlinear systems},
SIAM J.~Numer.~Anal.~\textbf{14} (1977), 611--615.

\bibitem{R3-Arb}
F.~Johansson,
\emph{Arb: efficient arbitrary-precision midpoint-radius interval arithmetic},
IEEE Trans.~Comput.~\textbf{66} (2017), 1281--1292.

\bibitem{R3-Singular}
W.~Decker, G.-M.~Greuel, G.~Pfister and H.~Sch\"onemann,
\emph{Singular 4.3.2 --- A computer algebra system for polynomial computations} (2022).

\bibitem{R3-msolve}
J.~Berthomieu, C.~Eder and M.~Safey El Din,
\emph{msolve: a library for solving polynomial systems},
in: Proc.\ ISSAC 2021, ACM, 2021, 51--58.

\bibitem{R3-Kato}
T.~Kato,
\emph{Perturbation Theory for Linear Operators}, 2nd ed., Springer, Berlin, 1976.

\bibitem{R4-EJT}
A.~Erchenko, D.~Jakobson, and A.~Tsypin,
\emph{Flexibility of eigenvalues for graph Laplacians arising from genus 3 surfaces},
arXiv:2604.26308 [math.SP], v1 29 April 2026, v2 14 October 2026.

\bibitem{R4-FGL}
S.~M.~Fallat, H.~Gupta, and J.~C.-H.~Lin,
\emph{Inverse eigenvalue problem for Laplacian matrices of a graph},
SIAM J. Matrix Anal. Appl. \textbf{46} (2025), no.~3, 1866--1886;
arXiv:2411.00292.

\bibitem{R4-Bhatia}
R.~Bhatia,
\emph{Matrix Analysis},
Graduate Texts in Mathematics, vol.~169, Springer, New York, 1997.

\bibitem{R4-Chaiken}
S.~Chaiken,
\emph{A combinatorial proof of the all minors matrix tree theorem},
SIAM J. Algebraic Discrete Methods \textbf{3} (1982), 319--329.

\bibitem{R4-Strassen}
V.~Strassen,
\emph{The existence of probability measures with given marginals},
Ann. Math. Statist. \textbf{36} (1965), 423--439.

\bibitem{R4-Renyi}
A.~R\'enyi,
\emph{On the theory of order statistics},
Acta Math. Acad. Sci. Hungar. \textbf{4} (1953), 191--231.

\bibitem{R4-Atlas}
R.~C.~Read and R.~J.~Wilson,
\emph{An Atlas of Graphs},
Oxford University Press, 1998 (as distributed with NetworkX, \texttt{graph\_atlas\_g}).

\bibitem{R4-Arb}
F.~Johansson,
\emph{Arb: efficient arbitrary-precision midpoint-radius interval arithmetic},
IEEE Trans. Comput. \textbf{66} (2017), 1281--1292; together with the FLINT library
(python-flint~0.9).

\end{thebibliography}
\end{document}
